\documentclass[11pt]{amsart}
\usepackage[T1]{fontenc}
\usepackage{lmodern}
\usepackage{amsmath,amssymb,amsthm,mathtools}
\usepackage[margin=1in]{geometry}
\usepackage{microtype}
\usepackage{enumitem}
\usepackage[colorlinks=true,linkcolor=blue,citecolor=blue,urlcolor=blue]{hyperref}
\newcommand{\R}{\mathbb R}
\newcommand{\C}{\mathbb C}
\newcommand{\N}{\mathbb N}
\newcommand{\Z}{\mathbb Z}
\newcommand{\Q}{\mathbb Q}
\newcommand{\dd}{\,\mathrm d}
\newcommand{\one}{\mathbf 1}
\newcommand{\e}{\mathrm e}
\newcommand{\norm}[1]{\left\|#1\right\|}
\newcommand{\abs}[1]{\left|#1\right|}
\newcommand{\ip}[2]{\left\langle #1,#2\right\rangle}
\DeclareMathOperator{\supp}{supp}
\DeclareMathOperator{\dist}{dist}
\DeclareMathOperator{\Lip}{Lip}
\DeclareMathOperator{\TV}{TV}
\newtheorem{theorem}{Theorem}[section]
\newtheorem{lemma}[theorem]{Lemma}
\newtheorem{proposition}[theorem]{Proposition}
\newtheorem{corollary}[theorem]{Corollary}
\theoremstyle{definition}
\newtheorem{definition}[theorem]{Definition}
\theoremstyle{remark}
\newtheorem{remark}[theorem]{Remark}
\newtheorem*{remark*}{Remark}
\newtheorem{maintheorem}{Theorem}
\renewcommand{\themaintheorem}{\Alph{maintheorem}}
\numberwithin{equation}{section}
\setlist[enumerate]{leftmargin=*,itemsep=4pt}
\setlist[itemize]{leftmargin=*,itemsep=3pt}
\allowdisplaybreaks[2]
\providecommand{\lean}[1]{}
\providecommand{\leanok}{}
\emergencystretch=2em
\title[Blueprint: triangles in the plane]{Triangles in the plane:\\a self-contained proof blueprint}
\author{}
\date{November 6, 2026}
\hypersetup{pdftitle={Triangles in the plane: a self-contained proof blueprint},
 pdfsubject={Expanded triangle proof including three-letter density Hales-Jewett and the analytic inputs}}
\begin{document}
\begin{abstract}
This blueprint expands the argument of the supplied manuscript on planar
triangles.  It includes proofs of the three-letter density Hales--Jewett
theorem and the analytic representation and rigidity facts needed by the
argument.  Weighted circle compactness, measurable stationarity, annular
cancellation, and frequency summation then give uniform smoothing of the
triangle counting form and the geometric conclusion.  All results are
arranged in dependency order, and each theorem, proposition, lemma, and
corollary after the introduction is followed immediately by its proof.
The foundational measure theory and calculus assumed are specified below.
\end{abstract}
\maketitle
\tableofcontents

\section{Introduction and foundations}\label{sec:introduction}

The purpose of this document is to make the proof structure and the
measure-theoretic, functional-analytic, and summation steps of the supplied
paper explicit.  The notation is kept close to that paper.  Definitions
are made before use, and the geometric density argument is placed after
the analytic estimates on which it depends.

For a measurable set $E\subset\R^2$, write
\[
 d^*(E)=\limsup_{L\to\infty}\sup_{z\in\R^2}
       \frac{|E\cap(z+[0,L]^2)|}{L^2}.
\]
Here and throughout, $|\cdot|$ denotes Lebesgue measure when applied
to a measurable subset of Euclidean space.

\begin{maintheorem}[Geometric conclusion]\label{intro:density}
\lean{Auto.intro_density}
\leanok
Let $\Delta\subset\R^2$ consist of three noncollinear points.  If
$E\subset\R^2$ is measurable and $d^*(E)>0$, there is $\lambda_0>0$
such that, for every $\lambda\ge\lambda_0$, the set $E$ contains a
translate of a rotation of $\lambda\Delta$.
\end{maintheorem}

The analytic assertion used to prove this is as follows.  Let
$J(x_1,x_2)=(-x_2,x_1)$ and $M=\alpha I+\beta J$ with $\beta\ne0$.
Let $\sigma$ be uniform probability measure on the unit circle.  For
a finite measure $\nu$ define
\[
 \Lambda_\nu(f_0,f_1,f_2)
  =\iint f_0(x)f_1(x+u)f_2(x+Mu)\dd\nu(u)\dd x.
\]
Write $\nu_s$ for the pushforward of $\nu$ under $u\mapsto su$,
and write $\|\Lambda_\nu\|$ for the
$L^2\times L^2\times L^\infty$ operator norm.

\begin{maintheorem}[Uniform smoothing]\label{intro:smoothing}
\lean{Auto.intro_smoothing}
\leanok
For every nonnegative $\rho\in L^1(\R^2)$ with $\int\rho=1$, put
$\rho_\varepsilon(u)=\varepsilon^{-2}\rho(u/\varepsilon)$.  Then
\[
 \lim_{\varepsilon\downarrow0}\sup_{s>0}
  \big\|\Lambda_{(\sigma*\rho_\varepsilon-\sigma)_s}\big\|=0.
\]
The same statement holds with the $L^\infty$ slot at any of the
three vertices.
\end{maintheorem}

These two introductory statements are proved in the main text, after
their prerequisites.  The analytic proof has three main parts: a rigidity
theorem for measurable stationary frequency fields; a compactness argument
for cancellation on a single frequency annulus; and a positive
telescoping estimate that sums the annuli without a decay rate.

\subsection*{Foundations and the scope of self-containment}

We work over the usual construction of Lebesgue measure, integration,
$L^p$ spaces, and Hilbert spaces.  The foundational facts used are
Tonelli and Fubini, monotone and dominated convergence, regularity of
finite Borel measures on Euclidean spaces, completion of Lebesgue
measure, measure extension, Radon--Nikodym, and monotone classes.
We also use elementary
Euclidean calculus, including change of variables, the inverse function
theorem, and convergent power series.  The Hilbert-space foundations
include orthogonal projection onto a closed subspace and the existence
of adjoints of bounded operators.  Basic inner-product algebra and
finite combinatorics are understood.  Thus ``self-contained'' here means
a proof over these standard analysis foundations, rather than a
reconstruction of measure theory itself.

All additional inputs used by the triangle argument are proved in the
text.  The analytic preliminaries include the needed Fourier facts,
sequential weak-* extraction, elementary compact-operator facts, the
positive-definite representation theorem, the exponential form of local
matrix groups, and the real-analytic zero-set argument.  These statements
are restricted to the settings in which they are applied.

The combinatorial preliminaries prove the density Hales--Jewett theorem
for three letters.  This is the theorem of Furstenberg and Katznelson
\cite{FK}; the proof here follows the finite density-increment approach
of Polymath \cite{Polymath}, with a tiling argument in the form developed
by Dodos, Kanellopoulos, and Tyros \cite{DKT}.  The two-letter input,
the comparisons of finite probability measures, the insensitive-set
tilings, and the density-increment iteration are all proved.  The
three-letter theorem is therefore an established earlier result when it
is used in the final geometric argument.



\section{Analytic prerequisites with proofs}\label{ext:section}

This section supplies the Fourier, compactness, spectral-measure,
matrix-group, and analytic-zero-set facts used in the argument.  The
foundations are Lebesgue measure and integration, the ordinary extension
theorem for a finite premeasure, elementary Euclidean calculus, and the
elementary Hilbert-space operations of taking adjoints and orthogonal
projections.  No spectral theorem for compact operators, representation
theorem for positive-definite functions, or local Lie-group theorem is
assumed.

\subsection{Fourier analysis in the required form}

\begin{definition}[Fourier convention and Gaussian kernels]
\label{ext:fourier-convention}
\lean{Auto.FCMeasure,Auto.FCMeasure.fourier,Auto.gaussKernel}
For $f\in L^1(\R^2)$ and a finite complex Borel measure $\eta$, put
\[
 \widehat f(\xi)=\int_{\R^2}f(x)\exp(-2\pi i x\cdot\xi)\dd x,
 \qquad
 \widehat\eta(\xi)=\int_{\R^2}\exp(-2\pi i x\cdot\xi)\dd\eta(x).
\]
The inverse transform uses the positive sign in the exponential and
is denoted by a check.  The Schwartz space $\mathcal S(\R^2)$ consists
of the smooth functions for which every derivative, multiplied by any
polynomial, is bounded.  For $t>0$ write
\[
 \gamma_t(x)=t^{-1}\exp(-\pi|x|^2/t).
\]
\end{definition}

\begin{lemma}[The convolution inequalities used below]\label{ext:young}
\lean{Auto.ext_young_one,Auto.ext_young_two,Auto.ext_young_top}
\leanok
If $\kappa$ is a finite complex Borel measure and
$p\in\{1,2,\infty\}$, then
\[
 \|\kappa*f\|_p\le\|\kappa\|_{\mathrm{TV}}\|f\|_p.
\]
In particular this holds when $\kappa$ has an $L^1$ density.
\end{lemma}
\begin{proof}
The assertion for $p=\infty$ follows by taking absolute
values inside the convolution integral; for $p=1$ it follows
by Tonelli and translation invariance.  For $p=2$,
Cauchy--Schwarz against the finite measure $|\kappa|$ gives
\[
 |(\kappa*f)(x)|^2
 \le\|\kappa\|_{\mathrm{TV}}
             \int|f(x-y)|^2\dd|\kappa|(y).
\]
Tonelli and translation invariance give the claimed $L^2$
bound after integration in $x$.  The same estimate shows
that the defining integral is absolutely convergent for
almost every $x$, by another application of Cauchy--Schwarz.
\end{proof}

\begin{lemma}[Density and continuity under translations and dilations]
\label{ext:approximation}
\lean{Auto.ext_approximation_dense,Auto.ext_approximation_translation,Auto.ext_approximation_translation_config,Auto.ext_approximation_dilation}
\leanok
Smooth compactly supported functions are dense in $L^p(\R^2)$ for
$p=1,2$.  For these two exponents, translations are continuous:
\[
 \|f(\cdot-h)-f\|_p\longrightarrow0\qquad(h\to0).
\]
For $F\in L^p(\R^2\times[0,2\pi))$, the same conclusion
holds for translation in the first two variables, with the
product measure $\dd x\dd\theta/(2\pi)$.
The normalized dilations $D_rf(x)=r^{-2}f(x/r)$ satisfy
$\|D_rf-f\|_1\to0$ as $r\to1$.
\end{lemma}
\begin{proof}
Truncating both the support and the magnitude of a function reduces
density to bounded functions on a bounded set.  Simple functions
approximate these in $L^p$.  By regularity of Lebesgue measure, their
level sets can be approximated in measure by finite unions of bounded
rectangles.  The indicator of a bounded rectangle can in turn be
approximated in $L^p$ by a smooth function between zero and one that
equals one away from a thin boundary layer and vanishes outside a
slightly enlarged rectangle.  These steps prove the density assertion.

For a smooth compactly supported function, translation continuity
follows from uniform continuity and a common bounded support for small
translations.  Translation is an isometry on $L^p$, so approximation
proves the assertion for every $f\in L^p$.  The same argument proves
dilation continuity in $L^1$, using $\|D_rf\|_1=\|f\|_1$ and the
common bounded support of $D_rg$ for $r$ near one and $g\in C_c^\infty$.
For the product-space assertion, apply planar translation
continuity to $F(\cdot,\theta)$ for almost every $\theta$.
Fubini writes the $p$th power of the product norm as the
integral of the $p$th powers of these sectional norms.
They are dominated by
$2^p\|F(\cdot,\theta)\|_p^p$, an integrable function of
$\theta$.  Dominated convergence gives the assertion.
\end{proof}

\begin{lemma}[The Gaussian Fourier calculation]\label{ext:gaussian}
\lean{Auto.ext_gaussian_integral,Auto.ext_gaussian_fourier,Auto.ext_gaussian_pointwise,Auto.ext_gaussian_uniform}
\leanok
The Gaussian kernels have integral one and satisfy
\[
 \widehat{\gamma_t}(\xi)=\exp(-\pi t|\xi|^2).
\]
If $f$ is bounded and continuous, then $\gamma_t*f(x)\to f(x)$
at every $x$ as $t\downarrow0$.  The convergence is uniform when
$f$ is bounded and uniformly continuous.
\end{lemma}
\begin{proof}
The elementary Gaussian integral is
$\int_\R\exp(-\pi x^2)\dd x=1$: square the integral and use polar
coordinates in the plane.  If
\[
 G(\xi)=\int_\R\exp(-\pi x^2)\exp(-2\pi i x\xi)\dd x,
\]
differentiation under the integral and integration by parts give
$G'(\xi)=-2\pi\xi G(\xi)$.  Since $G(0)=1$, differentiating
$\exp(\pi\xi^2)G(\xi)$ shows that $G(\xi)=\exp(-\pi\xi^2)$.
Taking products of the one-dimensional identity and then scaling
proves the displayed formula and the integral normalization.

After the substitution $y=\sqrt t\,z$,
\[
 (\gamma_t*f)(x)-f(x)
 =\int\exp(-\pi|z|^2)
               \big(f(x-\sqrt t\,z)-f(x)\big)\dd z.
\]
Pointwise convergence follows by dominated convergence.  For uniform
convergence, use uniform continuity on $|z|\le A$ and bound the
remaining integral by
$2\|f\|_\infty\int_{|z|>A}\exp(-\pi|z|^2)\dd z$; first increase
$A$, and then decrease $t$.
\end{proof}

\begin{lemma}[Fourier inversion and Plancherel]\label{ext:fourier}
\lean{Auto.ext_fourier_maps,Auto.ext_fourier_onto,Auto.ext_fourier_inversion,Auto.ext_fourier_plancherel,Auto.ext_fourier_unitary}
\leanok
The Fourier transform maps $\mathcal S(\R^2)$ onto itself, and
\[
 f(x)=\int\widehat f(\xi)\exp(2\pi i x\cdot\xi)\dd\xi
 \qquad(f\in\mathcal S).
\]
For $f,g\in\mathcal S$,
\[
 \int\widehat f(\xi)\overline{\widehat g(\xi)}\dd\xi
       =\int f(x)\overline{g(x)}\dd x.
\]
Consequently the Fourier transform extends uniquely to a unitary
operator on $L^2(\R^2)$.
\end{lemma}
\begin{proof}
Differentiation under the integral expresses derivatives of
$\widehat f$ as transforms of polynomial multiples of $f$.
Integration by parts expresses polynomial multiples of these
derivatives as transforms of derivatives of polynomial multiples
of $f$.  All the latter functions are integrable.  This proves
that $\widehat f$ is Schwartz.  The identical argument applies
to the inverse transform.

Fubini and Lemma~\ref{ext:gaussian} give, for $t>0$,
\[
 \int\widehat f(\xi)\exp(2\pi i x\cdot\xi)
                         \exp(-\pi t|\xi|^2)\dd\xi
       =(f*\gamma_t)(x).
\]
The left side converges to the inverse-transform integral by
dominated convergence, because $\widehat f\in L^1$.  The right
side converges to $f(x)$ by Lemma~\ref{ext:gaussian}.  This proves
inversion, including that the transform is onto the Schwartz space.

Apply inversion to $g$, take its complex conjugate, multiply by
$f$, and integrate.  Fubini is legitimate because
$\|f\|_1\|\widehat g\|_1<\infty$.  The result is the inner-product
identity.  In particular the transform is an $L^2$ isometry on
the Schwartz space.  This space is dense in $L^2$ by
Lemma~\ref{ext:approximation}, so the isometry extends uniquely
to $L^2$.  Its range is closed and contains the dense Schwartz
space, and is therefore all of $L^2$.
\end{proof}

\begin{lemma}[Uniqueness for finite measures]\label{ext:fourier-uniqueness}
\lean{Auto.ext_fourier_uniqueness,Auto.ext_fourier_uniqueness_dense}
\leanok
A finite complex Borel measure on $\R^2$ is determined by its
Fourier transform.  In particular, if $\nu$ is a finite positive
measure, the characters
$\xi\mapsto\exp(2\pi i s\cdot\xi)$ with $s\in\Q^2$ span a
dense subspace of $L^2(\nu)$.
\end{lemma}
\begin{proof}
It is enough to consider a finite complex measure $\eta$ with
$\widehat\eta=0$.  The Gaussian formula and Fubini imply
\[
 (\eta*\gamma_t)(x)
 =\int\exp(2\pi i x\cdot\xi)\exp(-\pi t|\xi|^2)
                          \widehat\eta(\xi)\dd\xi=0.
\]
For $\varphi\in C_c(\R^2)$, another use of Fubini gives
\[
 0=\int\varphi(x)(\eta*\gamma_t)(x)\dd x
   =\int(\varphi*\gamma_t)(y)\dd\eta(y).
\]
By uniform convergence in Lemma~\ref{ext:gaussian}, the last
integral tends to $\int\varphi\dd\eta$.  Thus this integral is
zero for every continuous compactly supported $\varphi$.
Regularity of a finite Borel measure on Euclidean space shows
that these test functions determine the measure: approximate
indicators of compact sets from above by continuous cutoffs in
arbitrarily small open neighborhoods, and then use inner
regularity on Borel sets.  Hence $\eta=0$.

If $g\in L^2(\nu)$ is orthogonal to all the stated characters,
the finite complex measure $\overline g\,\nu$ has Fourier
transform zero on $\Q^2$, with the sign of the argument
immaterial.  Its transform is continuous by dominated
convergence, so it is zero everywhere.  The first assertion
therefore gives $g=0$ almost everywhere.  The orthogonal
complement of the character span is zero, proving density.
\end{proof}

\begin{definition}[Fourier support for bounded functions]
\label{ext:bounded-support}
\lean{Auto.HasFourierSupportIn,Auto.HasL2FourierSupportIn}
For a compact set $K\subset\R^2$, a bounded measurable function $g$
has Fourier support in $K$ if
\[
 \int_{\R^2}g(x)\widehat\varphi(x)\dd x=0
 \quad\text{for every }\varphi\in C_c^\infty(\R^2)
 \text{ with }\supp\varphi\cap K=\varnothing.
\]
This is the test-function definition of distributional Fourier support;
the displayed integrals are absolutely convergent because
$\widehat\varphi$ is Schwartz.  For an $L^2$ function, Fourier support
in $K$ means that its $L^2$ Fourier transform vanishes almost everywhere
outside $K$.
\end{definition}

\begin{lemma}[Reproduction of an input with compact Fourier support]
\label{ext:reproduction}
\lean{Auto.ext_reproduction_L2,Auto.ext_reproduction_bounded,Auto.ext_reproduction_agree}
\leanok
Suppose $p\in\mathcal S(\R^2)$ and $\widehat p=1$ on a
neighborhood of a compact set $K$.  If $g\in L^2$ or $g\in L^\infty$
has Fourier support in $K$, then $p*g=g$ almost everywhere.
For a function in $L^2\cap L^\infty$, the two support conditions in
Definition~\ref{ext:bounded-support} agree.
\end{lemma}
\begin{proof}
For $L^2$ inputs, the identity
$\widehat{p*g}=\widehat p\,\widehat g$ holds first for Schwartz
functions by Fubini and then by $L^2$ approximation, using
Lemmas~\ref{ext:young} and~\ref{ext:fourier}.  The support assumption
then proves the assertion by injectivity of the $L^2$ Fourier transform.

For bounded inputs we first extend the defining support condition to
Schwartz tests $\varphi$ that vanish on a neighborhood of $K$.
Choose $\eta\in C_c^\infty$ equal to one near zero and put
$\varphi_R(\xi)=\eta(\xi/R)\varphi(\xi)$.  Each $\varphi_R$ is an
allowed compactly supported test.  Write
$\Delta=\partial_1^2+\partial_2^2$ for the Laplacian.
If $v\in\mathcal S$, integration by parts gives
\[
 (1+4\pi^2|x|^2)^2\widehat v(x)
       =\widehat{(1-\Delta)^2v}(x),
 \qquad
 \|\widehat v\|_1
       \le C\sum_{|\alpha|\le4}\|\partial^\alpha v\|_1.
\]
The second inequality follows by bounding the transform on the right
by the $L^1$ norm of its input and integrating
$(1+4\pi^2|x|^2)^{-2}$ over the plane.
The product rule and rapid decay of $\varphi$ give
$\|\partial^\alpha(\varphi_R-\varphi)\|_1\to0$ for every
$|\alpha|\le4$: the undifferentiated cutoff contributes a tail
integral, and each differentiated cutoff contributes a factor
$R^{-1}$ or smaller times a rapidly decaying derivative.
Consequently $\|\widehat\varphi_R-\widehat\varphi\|_1\to0$.
Pairing with bounded $g$ proves $\int g\widehat\varphi=0$.

For any $\psi\in\mathcal S$, Fubini now gives
\[
 \int(p*g-g)(x)\widehat\psi(x)\dd x
   =\int g(y)\widehat{(\widehat p-1)\psi}(y)\dd y=0.
\]
For the convolution term, absolute integrability follows from
$\|g\|_\infty\|p\|_1\|\widehat\psi\|_1<\infty$.
The final equality holds because $(\widehat p-1)\psi$ is Schwartz
and vanishes near $K$.  Lemma~\ref{ext:fourier} says that Fourier
transforms of Schwartz tests exhaust the Schwartz space.  Thus the
bounded function $p*g-g$ annihilates every $C_c^\infty$ test and,
by their $L^1$ density from Lemma~\ref{ext:approximation}, every
$L^1$ test.  It is zero almost everywhere.

Finally, for $g\in L^2\cap L^\infty$, Plancherel gives
$\int g\widehat\varphi=\int\widehat g\,\varphi$.
Vanishing of the latter pairings for all smooth tests supported
outside $K$ is equivalent to $\widehat g=0$ there almost everywhere.
For the nontrivial direction, cover $\R^2\setminus K$ by countably
many balls whose closures are disjoint from $K$.  Smooth tests supported
in a ball are dense in its $L^2$ space: first multiply a function by
smooth cutoffs increasing to one inside that ball, and then approximate
each resulting function in $L^2(\R^2)$ using
Lemma~\ref{ext:approximation}, multiplying the approximants by another
cutoff equal to one on its compact support and supported in the ball.
Thus $\widehat g$ vanishes on each ball almost everywhere.
This proves agreement of the definitions.
\end{proof}

\subsection{The compactness facts used for circle operators}

\begin{definition}[Compact operator]\label{ext:compact-definition}
\lean{IsCompactOperator}
A bounded linear operator between Hilbert spaces is compact if
the image of every bounded sequence has a norm-convergent
subsequence.  Equivalently, the image of its unit ball has a
finite $\varepsilon$-net for every $\varepsilon>0$.
\end{definition}

\begin{lemma}[Elementary compactness principles]\label{ext:compact-facts}
\lean{Auto.ext_compact_facts_norm_adjoint,Auto.ext_compact_facts_norm_comp_adjoint,Auto.ext_compact_facts_finite_rank,Auto.ext_compact_facts_limit,Auto.ext_compact_facts_integral,Auto.ext_compact_facts_weak}
\leanok
For a bounded operator between Hilbert spaces,
$\|T\|=\|T^*\|$ and $\|TT^*\|=\|T\|^2$.
The following compactness facts also hold.
\begin{enumerate}
\item Finite-rank operators are compact, and an operator-norm
limit of compact operators is compact.
\item An integral operator with kernel in $L^2(X\times Y)$,
where $X,Y$ are bounded measurable subsets of Euclidean space,
is compact from $L^2(Y)$ to $L^2(X)$.
\item If $T$ is compact and a bounded sequence $f_n$ converges
weakly to $f$, then $Tf_n\to Tf$ in norm.
\end{enumerate}
\end{lemma}
\begin{proof}
Taking the supremum of $|\langle Tx,y\rangle|$ over unit
vectors in the two spaces gives $\|T\|=\|T^*\|$.
The inequality $\|TT^*\|\le\|T\|^2$ follows by composition,
and $\|T^*y\|^2=\langle TT^*y,y\rangle$ gives the reverse
inequality after taking the supremum over unit $y$.

The first assertion for finite rank follows from compactness of
bounded closed sets in finite-dimensional Euclidean space.
If $T_m\to T$ in operator norm, approximate the image of the
unit ball under $T$ by its image under one $T_m$ and then by a
finite net for that image.  This gives a finite net for $T$ at
every positive accuracy.  A successive-net or diagonal argument
shows that every sequence in this image has a Cauchy subsequence;
completeness makes it convergent.

For the second assertion, finite sums of kernels
$a(x)b(y)$ are dense in $L^2(X\times Y)$.  Indeed, rectangular
simple functions are dense in the product measure space, by
approximation of measurable sets using the algebra of finite
unions of measurable rectangles.  Such finite sums define
finite-rank operators.  Cauchy--Schwarz shows that the operator
norm of the difference of two integral operators is at most
the $L^2$ norm of the difference of their kernels.  Apply the
first assertion.

For the third, suppose norm convergence fails.  A subsequence
stays a positive distance from $Tf$, and a further subsequence
converges in norm to some $v$ by compactness.  For every test
vector $h$, the identity
$\langle Tf_n,h\rangle=\langle f_n,T^*h\rangle$ shows that the
weak limit is $Tf$.  The norm limit $v$ must therefore equal
$Tf$, a contradiction.
\end{proof}

\begin{lemma}[Compactness from the product with the adjoint]
\label{ext:ttstar}
\lean{Auto.ext_ttstar}
\leanok
If $T:H\to H'$ is bounded and $TT^*:H'\to H'$ is compact,
then $T$ is compact.
\end{lemma}
\begin{proof}
For a bounded sequence $u_n$ in $H'$, choose a subsequence on
which $TT^*u_n$ converges.  On that subsequence,
\[
 \|T^*(u_n-u_m)\|^2
 =\langle TT^*(u_n-u_m),u_n-u_m\rangle\longrightarrow0.
\]
Thus $T^*$ is compact.  Given $\varepsilon>0$, choose a finite
$\varepsilon$-net for $T^*$ of the unit ball.  Let $F\subset H$
be the span of its finitely many centers and $P_F$ the orthogonal
projection onto $F$.  Then
$\|(I-P_F)T^*\|\le\varepsilon$.  Taking adjoints gives
$\|T(I-P_F)\|\le\varepsilon$.  The operator $TP_F$ has finite
rank and approximates $T$ in norm.  Lemma~\ref{ext:compact-facts}
proves compactness.
\end{proof}

\begin{lemma}[A gap for a nonnegative compact perturbation]
\label{ext:compact-gap}
\lean{Auto.ext_compact_gap_ker,Auto.ext_compact_gap}
\leanok
Let $K$ be a compact self-adjoint operator on a Hilbert space,
and suppose $B=I+K$ is nonnegative.  Then $\ker B$ is finite
dimensional, and there is $\gamma>0$ such that
\[
 \langle Bf,f\rangle\ge
       \gamma\operatorname{dist}(f,\ker B)^2
       \qquad(f\in H).
\]
\end{lemma}
\begin{proof}
If the kernel were infinite dimensional, it would contain an
orthonormal sequence $u_n$.  Since $Ku_n=-u_n$, its image under
$K$ would have no convergent subsequence, contrary to compactness.

For any bounded nonnegative operator $B$, positivity of
$\langle B(x+ty),x+ty\rangle$ for real or complex $t$ gives
the Cauchy--Schwarz inequality for its quadratic form:
\[
 |\langle Bx,y\rangle|^2
       \le\langle Bx,x\rangle\langle By,y\rangle.
\]
When the second diagonal term is positive, minimize the
quadratic polynomial in $t$; when it is zero, letting $t$
vary first shows that the cross term is zero.  Taking the
supremum over unit vectors $y$ gives
\[
 \|Bx\|^2\le\|B\|\langle Bx,x\rangle.
\]
If there were no positive lower bound for the quadratic form
on unit vectors orthogonal to $\ker B$, there would be such
vectors $x_n$ with $\langle Bx_n,x_n\rangle\to0$.  The last
inequality gives $Bx_n\to0$.  On a subsequence $Kx_n$ converges,
so $x_n=Bx_n-Kx_n$ converges in norm to a unit vector $x$.
This vector is still orthogonal to the kernel, but $Bx=0$,
a contradiction.  Finally decompose an arbitrary $f$ into
its kernel component and its orthogonal component.  The
quadratic form vanishes on the first component and has no
cross term, proving the stated estimate.
\end{proof}

\begin{lemma}[Sequential weak-* extraction]\label{ext:weakstar}
\lean{Auto.ext_weakstar_plane,Auto.ext_weakstar_plane_countable,Auto.ext_weakstar_config,Auto.ext_weakstar_config_countable}
\leanok
On $\R^2$, and on $\R^2\times[0,2\pi)$ with measure
$\dd x\dd\theta/(2\pi)$, every sequence bounded in
$L^\infty$ has a weak-* convergent subsequence.  A common
subsequence can be chosen for any countable family of such
bounded sequences.
\end{lemma}
\begin{proof}
Write the measure as $m$.  Step functions on rectangles with
rational endpoints and rational real and imaginary values
form a countable dense subset of $L^1$, after allowing finite
linear combinations.  Density follows by truncation and
regularity of Lebesgue measure, as in
Lemma~\ref{ext:approximation}.  If $\|f_n\|_\infty\le A$,
the pairings of $f_n$ with each of these tests form a bounded
scalar sequence.  Successive subsequence extraction and a
diagonal choice make all these pairings converge.  The bound
$A\|\varphi\|_1$ then extends the limit to a bounded linear
functional $L$ on all of $L^1$, and gives convergence against
every $L^1$ test.

For clarity, this functional is represented by a bounded
function without invoking a compactness theorem for dual
spaces.  Exhaust the underlying space by sets $X_k$ of finite
measure.  On each $X_k$, the set function
$A\mapsto L(\one_A)$ is a complex measure: for disjoint sets,
countable additivity follows from convergence of the partial
union indicators in $L^1(X_k)$.  Its total variation is at
most $A m$, by the defining norm bound.  The Radon--Nikodym
theorem gives a density $g_k$ with $|g_k|\le A$.
Uniqueness of densities makes these functions agree on
overlaps.  They define a global $g\in L^\infty$ satisfying
$L(\varphi)=\int g\varphi\dd m$, first for simple functions
of finite support and then for all $L^1$ functions by density.
This is the required weak-* limit.

For countably many sequences, perform the scalar diagonal
selection over the countable set of all pairs consisting of
a sequence index and a test index.  The preceding construction
then supplies the limit for each sequence on the same
subsequence.
\end{proof}

\subsection{Spectral measures of the correlation functions}

\begin{definition}[Positive-definite function]\label{ext:pd-definition}
\lean{Auto.IsPositiveDefinite}
A function $c:\R^2\to\C$ is positive definite if
\[
 \sum_{a,b=1}^N z_a\overline{z_b}\,c(s_a-s_b)\ge0
\]
for every finite family $s_a\in\R^2$, $z_a\in\C$; the
displayed quantity is required to be real.
\end{definition}

\begin{lemma}[Integrable positive-definite functions]
\label{ext:pd-integrable}
\lean{Auto.ext_pd_integrable_elementary,Auto.ext_pd_integrable_fourier}
\leanok
If $c$ is positive definite, then $c(0)\ge0$,
$c(-s)=\overline{c(s)}$, and $|c(s)|\le c(0)$.
If, in addition, $c$ is continuous and integrable, then
$\widehat c$ is a nonnegative integrable function,
\[
 \int\widehat c(\xi)\dd\xi=c(0),
 \qquad
 c(s)=\int\exp(2\pi i s\cdot\xi)\widehat c(\xi)\dd\xi.
\]
\end{lemma}
\begin{proof}
Use one point and then the two points $0,s$ in the defining
inequality.  Varying their complex coefficients shows that
the corresponding matrix is Hermitian and nonnegative.
Minimizing its quadratic form gives
$|c(s)|^2\le c(0)^2$.  This also shows that $c$ is identically
zero if $c(0)=0$.

For a continuous compactly supported function $\varphi$, the
defining finite inequalities and Riemann-sum approximation give
\[
 \iint c(x-y)\varphi(x)\overline{\varphi(y)}\dd x\dd y\ge0.
\]
Fix $\xi\in\R^2$ and a real $\varphi\in C_c^\infty$ with
$\|\varphi\|_2=1$.  Apply this inequality to
$L^{-1}\varphi(x/L)\exp(-2\pi i\xi\cdot x)$.
After writing $h=x-y$, the resulting integral is
\[
 \int c(h)\exp(-2\pi i\xi\cdot h)
       \left(\int\varphi(t+h/L)\varphi(t)\dd t\right)\dd h.
\]
The inner correlation has absolute value at most one and
tends to one for each fixed $h$, by translation continuity
in $L^2$.  Since $c\in L^1$, dominated convergence shows
that the last integral tends to $\widehat c(\xi)$.  Thus
$\widehat c$ is nonnegative everywhere.

For $t>0$, Fubini and Lemma~\ref{ext:gaussian} give
\[
 \int\widehat c(\xi)\exp(-\pi t|\xi|^2)\dd\xi
       =\int c(s)\gamma_t(s)\dd s\longrightarrow c(0).
\]
Monotone convergence on the left as $t\downarrow0$ proves
integrability and the mass formula.  Finally,
\[
 \int\widehat c(\xi)\exp(2\pi i s\cdot\xi)
                         \exp(-\pi t|\xi|^2)\dd\xi
       =(c*\gamma_t)(s).
\]
The left side converges by the newly proved integrability,
and the right side converges by continuity and boundedness
of $c$.  This proves the representation.
\end{proof}

\begin{lemma}[A subsequence of uniformly tight measures]
\label{ext:tight-measures}
\lean{Auto.ext_tight_measures}
\leanok
Let $\nu_n$ be positive Borel measures on $\R^2$ with the
same finite mass $m$, and suppose
\[
 \lim_{A\to\infty}\sup_n\nu_n(\{x:|x|>A\})=0.
\]
There are a subsequence and a positive Borel measure $\nu$
of mass $m$ such that
\[
 \int g\dd\nu_n\longrightarrow\int g\dd\nu
\]
for every bounded continuous function $g$ on $\R^2$.
\end{lemma}
\begin{proof}
The case $m=0$ is immediate, so divide by $m$ and assume
$m=1$.  The coordinatewise map
\[
 \Phi(x_1,x_2)=
 \left(\frac12+\frac{\arctan x_1}{\pi},
       \frac12+\frac{\arctan x_2}{\pi}\right)
\]
is a homeomorphism onto $(0,1)^2$.  Push the measures forward
to the compact square $K=[0,1]^2$.

Here is an explicit subsequence construction on $K$.  Encode
each coordinate by its binary expansion, taking the expansion
ending in zeros at dyadic points and taking all ones at the
endpoint one.  This is a measurable map from $K$ into
$\Omega=(\{0,1\}^2)^{\N}$, and the continuous binary-sum
map $\pi:\Omega\to K$ recovers the original point.
For each finite word $w$ on the four-symbol alphabet, let
$p_n(w)$ be the mass of the cylinder specified by $w$ in
the encoded measure.  There are countably many such words.
A diagonal subsequence makes all $p_n(w)$ converge to numbers
$p(w)$.  They satisfy
\[
 p(\varnothing)=1,
 \qquad p(w)=\sum_{a\in\{0,1\}^2}p(wa),
 \qquad p(w)\ge0.
\]
These numbers define a probability on the algebra of finite
unions of cylinders, by refining to a common word length.

This algebra probability is a premeasure.  Indeed, cylinders
are both open and closed for the metric in which agreement
in the first $k$ coordinates gives distance at most $2^{-k}$.
The space $\Omega$ is compact: successive choice of constant
coordinates gives a convergent subsequence of any sequence,
or equivalently the finitely many length-$k$ cylinders give
finite nets and the space is complete.  If an algebra set is
a disjoint countable union of algebra sets, compactness gives
a finite subcover, so only finitely many of the disjoint
pieces are nonempty.  Finite additivity is therefore already
the countable additivity needed on the algebra.  The ordinary
measure-extension theorem gives a probability $P$ on the
Borel sets of $\Omega$.  Put $\widetilde\nu=\pi_*P$.

For a continuous function $g$ on $K$, approximate $g\circ\pi$
by a function constant on cylinders of length $k$, using any
point in the corresponding closed dyadic square.  The
uniform error is at most the modulus of continuity of $g$
at $\sqrt2\,2^{-k}$.  The integrals of this cylinder function
converge because they are finite sums involving $p_n(w)$.
First take the subsequence limit and then let $k\to\infty$.
This proves convergence of the pushed-forward measures
against every continuous function on $K$.

The measure $\widetilde\nu$ gives no mass to $\partial K$.
Given $\delta>0$, tightness supplies $A$ for which all original
measures have mass at most $\delta$ outside $[-A,A]^2$.
There is a continuous function on $K$, between zero and one,
equal to zero on $\Phi([-A,A]^2)$ and equal to one on
$\partial K$; for example use the quotient of the distance
to the first set by the sum of distances to the two disjoint
closed sets.  Its integrals against the pushed-forward
measures are at most $\delta$, and the same is true for
$\widetilde\nu$.  Hence $\widetilde\nu(\partial K)\le\delta$.
Let $\delta\downarrow0$.

Pull $\widetilde\nu$ back to $(0,1)^2$ by $\Phi^{-1}$ to
obtain a probability $\nu$ on $\R^2$.  Compactly supported
continuous tests on $\R^2$ extend by zero to continuous
tests on $K$ after this change of coordinates, so their
integrals converge.  For a bounded continuous test, insert
a continuous compact cutoff equal to one on $[-A,A]^2$.
The discarded integrals are uniformly small for the original
measures by tightness, and small for $\nu$ because its finite
mass is exhausted by these squares.  This proves convergence
for every bounded continuous test and finishes the proof.
\end{proof}

\begin{theorem}[Spectral representation in the required form]
\label{ext:bochner}
\lean{Auto.ext_bochner,Auto.ext_bochner_mass_le_one}
\leanok
Let $c:\R^2\to\C$ be a Lipschitz positive-definite function.
There is a unique finite positive Borel measure $\nu$ such that
\[
 c(s)=\int\exp(2\pi i s\cdot\xi)\dd\nu(\xi)
 \qquad(s\in\R^2),\qquad \nu(\R^2)=c(0).
\]
In particular, $c(0)\le1$ implies that the representing measure
has mass at most one.
\end{theorem}
\begin{proof}
Put $m=c(0)$.  The elementary part of
Lemma~\ref{ext:pd-integrable} gives $m\ge0$ and $|c|\le m$;
if $m=0$ there is nothing to prove.  For $0<\varepsilon\le1$
put
\[
 c_\varepsilon(s)=c(s)\exp(-\pi\varepsilon|s|^2).
\]
This function is integrable and positive definite.  To see the
latter directly, the Gaussian formula represents
$\exp(-\pi\varepsilon|s_a-s_b|^2)$ as the integral of
$\exp(2\pi i(s_a-s_b)\cdot\xi)$ against the positive
probability density $\gamma_\varepsilon(\xi)$.  Inside that
integral the positive-definite inequality for
$c_\varepsilon$ is the one for $c$ with coefficients
$z_a\exp(2\pi i s_a\cdot\xi)$.

By Lemma~\ref{ext:pd-integrable},
\[
 \dd\nu_\varepsilon(\xi)
       =\widehat{c_\varepsilon}(\xi)\dd\xi
\]
is a positive measure of mass $m$, with characteristic function
$c_\varepsilon$.  These measures are uniformly tight.  Indeed,
for $t>0$ the Gaussian identity gives
\begin{align*}
 \int\big(1-\exp(-\pi t|\xi|^2)\big)\dd\nu_\varepsilon(\xi)
 &=\int\gamma_t(s)
       \big(m-\operatorname{Re}c_\varepsilon(s)\big)\dd s\\
 &\le\int\gamma_t(s)
       \big(L|s|+m\pi|s|^2\big)\dd s\\
 &\le C\big(L\sqrt t+mt\big),
\end{align*}
where $L$ is a Lipschitz constant for $c$.  In the second line
we used $1-\exp(-\pi\varepsilon|s|^2)\le\pi|s|^2$ and
$|c(s)-c(0)|\le L|s|$.  The last line follows by scaling
in the first and second Gaussian moments.  Set $t=A^{-2}$.
On $|\xi|\ge A$, the integrand on the left is at least
$1-\exp(-\pi)$.  Consequently
\[
 \sup_{0<\varepsilon\le1}\nu_\varepsilon(\{|\xi|\ge A\})
       \le C\big(L/A+m/A^2\big)\longrightarrow0.
\]

Apply Lemma~\ref{ext:tight-measures} to
$\nu_{1/n}$.  The limiting measure $\nu$ has mass $m$.
For each $s$, the character $\xi\mapsto\exp(2\pi i s\cdot\xi)$
is bounded and continuous, so
\[
 \int\exp(2\pi i s\cdot\xi)\dd\nu(\xi)
   =\lim_n c(s)\exp(-\pi|s|^2/n)=c(s).
\]
Uniqueness follows from Lemma~\ref{ext:fourier-uniqueness}.
\end{proof}

\begin{lemma}[Measurable dependence of spectral measures]
\label{ext:measurable-spectral}
\lean{Auto.ext_measurable_spectral}
\leanok
Suppose $z\mapsto c_z(s)$ is measurable for each $s\in\R^2$,
and every $c_z$ is Lipschitz and positive definite with
$c_z(0)\le1$.  Let $\nu_z$ be the measures supplied by
Theorem~\ref{ext:bochner}.  Then $z\mapsto\nu_z(B)$ is
measurable for every Borel set $B\subset\R^2$.
\end{lemma}
\begin{proof}
The functions $(z,s)\mapsto c_z(s)$ are jointly measurable:
replace $s$ by a sequence of rational-grid approximations
and use continuity in $s$.  For $\varphi\in\mathcal S$,
Fourier inversion and Fubini give
\[
 \int\varphi(\xi)\dd\nu_z(\xi)
      =\int\widehat\varphi(s)c_z(s)\dd s.
\]
The right side is measurable in $z$, since
$|c_z(s)|\le c_z(0)\le1$ and $\widehat\varphi\in L^1$.

For an open set $O$, choose smooth compactly supported
functions $0\le\varphi_m\le1$ in $O$ increasing pointwise
to $\one_O$.  One explicit construction takes cutoffs
$\eta_m$ equal to one on a compact exhaustion of $O$, with
support inside $O$, and sets
$\varphi_m=1-\prod_{j=1}^m(1-\eta_j)$.
Thus $\nu_z(O)=\lim_m\int\varphi_m\dd\nu_z$ is measurable.
The class of sets with measurable $\nu_z$-mass is closed
under complements and countable disjoint unions, since the
total mass is the measurable function $c_z(0)$.  Open sets
are closed under finite intersections and generate the
Borel sets.  The elementary monotone-class argument therefore
extends measurability to all Borel sets.
\end{proof}

\begin{definition}[Measurable measure kernels]\label{ext:kernel-definition}
\lean{Auto.IsMeasurableKernel}
Let $(X,\mathcal A)$ be a measurable space.  A family
$z\mapsto\nu_z$ of finite positive Borel measures on $\R^q$ is a
\emph{measurable kernel} if $z\mapsto\nu_z(B)$ is
$\mathcal A$-measurable for every Borel set $B\subset\R^q$.
\end{definition}

\begin{lemma}[Integration against measurable kernels]
\label{ext:kernel-integration}
\lean{Auto.ext_kernel_integration_measurable,Auto.ext_kernel_integration_measurable_complex,Auto.ext_kernel_integration_measure,Auto.ext_kernel_integration_prod}
\leanok
Let $(X,\mathcal A,\lambda)$ be a sigma-finite measure space and let
$z\mapsto\nu_z$ be a measurable kernel on $\R^q$.
For every nonnegative $\mathcal A\otimes\mathcal B(\R^q)$-measurable
function $F$, the function
$z\mapsto\int F(z,\xi)\dd\nu_z(\xi)$ is measurable.
There is a sigma-finite measure $\mathfrak m$ on the product space
satisfying
\[
 \int F\dd\mathfrak m
   =\int_X\int_{\R^q}F(z,\xi)\dd\nu_z(\xi)\dd\lambda(z).
\]
If $z\mapsto\eta_z$ is another measurable finite positive kernel,
then $z\mapsto\nu_z\otimes\eta_z$ is a measurable kernel.
These identities also hold for integrable complex functions.
\end{lemma}
\begin{proof}
For the indicator of a rectangle $A\times B$, the inner integral is
$\one_A(z)\nu_z(B)$ and is measurable.  The class of product-measurable
sets for which the section integral is measurable is closed under
complements, because $\nu_z(\R^q)$ is finite and measurable, and under
countable disjoint unions by monotone convergence.  Rectangles form
a system closed under finite intersections and generate the product
sigma-algebra.  The monotone-class theorem therefore proves the
claim for indicators of all product-measurable sets.  Increasing simple
approximations prove it for all nonnegative $F$.

For such a set $A$, define
\[
 \mathfrak m(A)=\int_X\nu_z(A_z)\dd\lambda(z),
 \qquad A_z=\{\xi:(z,\xi)\in A\}.
\]
Sections are Borel measurable, and the preceding argument makes this
definition meaningful.  Countable additivity follows from countable
additivity of each $\nu_z$ and monotone convergence under the outer
integral.  If $X_n$ increase to $X$ and $\lambda(X_n)<\infty$, the sets
\[
 \bigl(X_n\cap\{z:\nu_z(\R^q)\le n\}\bigr)\times\R^q
\]
exhaust the product and have $\mathfrak m$-measure at most
$n\lambda(X_n)$.  Thus $\mathfrak m$ is sigma-finite.
The integral identity follows first for indicators from its definition,
then for simple functions, and then by monotone convergence.

For product kernels, the mass of a rectangle $B\times C$ is
$\nu_z(B)\eta_z(C)$ and is measurable.  The same monotone-class
argument, using the finite measurable total mass, extends this to every
Borel subset of the product Euclidean space.  Finally split an
integrable complex function into positive and negative real and
imaginary parts.  Ordinary Tonelli applied to the constructed measure
and an additional independent measure permits the angular integrations
used below to be interchanged as well.
\end{proof}

\subsection{Local matrix groups and analytic zero sets}

\begin{lemma}[A continuous local matrix group is exponential]
\label{ext:matrix-group}
\lean{Auto.ext_matrix_group,Auto.ext_matrix_group_analytic}
\leanok
Let $U$ be a neighborhood of zero in $\R^2$.  Suppose
$C:U\to\mathrm{GL}_d(\R)$ is continuous, $C(0)=I$, and
$C(h+k)=C(h)C(k)$ whenever $h,k,h+k\in U$.
There are commuting real matrices $T_1,T_2$ such that,
on a possibly smaller neighborhood of zero,
\[
 C(h)=\exp(h_1T_1+h_2T_2).
\]
The expression on the right is real analytic for all
$h\in\R^2$.
\end{lemma}
\begin{proof}
For $i=1,2$ write $C_i(t)=C(te_i)$ near zero.
Choose $a>0$ sufficiently small that the matrix
\[
 A_i=\int_0^a C_i(t)\dd t
\]
is invertible: $a^{-1}A_i\to I$ as $a\downarrow0$, and
invertibility is open because the determinant is continuous.
For $h$ in a small interval around zero, the local group law
gives
\[
 C_i(h)A_i=\int_h^{h+a}C_i(t)\dd t.
\]
The right side is differentiable even though differentiability
of $C_i$ was not assumed.  Multiplying by $A_i^{-1}$ first
shows that $C_i$ is differentiable, and differentiating gives
\[
 C_i'(h)=C_i(h)T_i,
 \qquad T_i=(C_i(a)-I)A_i^{-1}.
\]
The derivative of $C_i(h)\exp(-hT_i)$ is zero.  Its value at
zero is $I$, so $C_i(h)=\exp(hT_i)$ near zero.  The matrix
exponential and its derivative here are obtained directly
from its absolutely convergent power series.

Commutativity of addition and the local group law yield
\[
 \exp(sT_1)\exp(tT_2)=\exp(tT_2)\exp(sT_1)
\]
for small $s,t$.  Differentiate at $s=t=0$ to get
$T_1T_2=T_2T_1$.  Hence
\[
 C(h)=C(h_1e_1)C(h_2e_2)
     =\exp(h_1T_1)\exp(h_2T_2)
     =\exp(h_1T_1+h_2T_2).
\]
The last equality follows by multiplying the exponential
series and using the binomial formula for commuting matrices.
Those series converge locally uniformly with all derivatives,
which also proves real analyticity on the whole plane.
\end{proof}

\begin{lemma}[Zero sets of real analytic functions]
\label{ext:analytic-zero}
\lean{Auto.ext_analytic_zero,Auto.ext_analytic_zero_config}
\leanok
Let $F:U\to\R$ be real analytic on a connected open set
$U\subset\R^d$, where $d=1,2,$ or $3$.  If $F$ vanishes
on a set of positive Lebesgue measure, it vanishes identically.
In particular, a real analytic function on
$\R^2\times(\R/(2\pi\Z))$ which vanishes on a set of positive
product measure is identically zero.
\end{lemma}
\begin{proof}
Let $Z_\infty$ be the set of points at which all partial
derivatives of $F$, including $F$ itself, vanish.  This set
is closed in $U$ by continuity of the derivatives.  It is
also open: the convergent Taylor series at one of its points
is zero, so $F$ is identically zero in a neighborhood, where
all of its derivatives vanish as well.  Connectedness shows
that $Z_\infty$ is either empty or all of $U$.  In the latter
case the assertion is proved.

Suppose therefore that $F$ is not identically zero.  At
every point $x$ some Taylor coefficient is nonzero.  If the
first nonzero homogeneous term has degree $k$, it takes a
nonzero value in some direction $v$, so
$\partial_v^kF(x)\ne0$.  For $k=0$ this simply means
$F(x)\ne0$.  By continuity there is a ball about $x$ on
which this derivative never vanishes.  On every line
parallel to $v$ intersecting the ball, the restriction of
$F$ has at most $k$ zeros.  Otherwise repeated Rolle's
theorem would give a zero of its $k$th derivative on that
interval.  Fubini, after an orthogonal change of coordinates,
therefore shows that the zero set of $F$ in the ball has
$d$-dimensional measure zero.  A countable subcover of these
balls covers $U$, so the entire zero set has measure zero,
contradicting the hypothesis.

For the final assertion, lift the periodic angular variable
to $\theta\in\R$.  This gives a real analytic function on
the connected space $\R^3$.  A positive-measure zero set in
the quotient has a positive-measure lift in one angular
period.  Apply the result just proved to the lifted function.
\end{proof}


\section{Finite words, binary density, and tilings}\label{sec:comb-prelim}

The elementary combinatorial material in this section supplies the binary
base case and the tiling step of the density argument.  The organization
of the tiling proof follows the uniform-measure argument of
Dodos, Kanellopoulos, and Tyros~\cite{DKT}; every ingredient needed here is proved
below.

\begin{definition}[Words, subspaces, and uniform density]\label{cj:def-words}
\lean{Auto.unifDensity,Auto.subspacePts,Auto.relDensity,Auto.binEmb,Auto.binaryCube,Auto.binaryFace,Auto.concatW,Auto.suffixFiber,Auto.ContainsSubspace,Auto.ContainsLine}
For a positive integer $k$, let $[k]=\{1,\ldots,k\}$.
For $k\ge2$ and $n\ge0$, let
$[k]^n$ be the set of words of length $n$ over $[k]$; $[k]^0$ consists
of the empty word.  Write $\mu_{k,n}(A)=|A|/k^n$ for uniform density.
The concatenation of words $x$ and $y$ is denoted by $x\frown y$.

An $m$-dimensional combinatorial subspace of $[k]^n$ is specified by
pairwise disjoint, nonempty coordinate sets $I_1,\ldots,I_m$ and by
fixed letters on their complement.  Its coordinate map
$\phi:[k]^m\longrightarrow[k]^n$ puts the letter $a_j$ on every
coordinate in $I_j$ and the specified fixed letters elsewhere.  The
subspace is $V=\phi([k]^m)$.  We order the variable sets by their least
coordinates when a definite choice of coordinate map is required.
A one-dimensional subspace is called a \emph{combinatorial line}.
For $A\subseteq[k]^n$, put
$\mu_V(A)=|A\cap V|/|V|$.
If $V\subseteq[3]^n$, its \emph{binary face} is
$V^{(2)}=\phi([2]^m)$; fixed letters of $V$ may still equal $3$.
For $x\in[k]^\ell$, the suffix fiber of $A\subseteq[k]^n$ at $x$ is
\[
 A_x=\{y\in[k]^{n-\ell}:x\frown y\in A\}.
\]
\end{definition}

\begin{lemma}[Subspace bookkeeping]\label{cj:bookkeeping}
\lean{Auto.cj_bookkeeping_card,Auto.cj_bookkeeping_comp,Auto.cj_bookkeeping_concat,Auto.cj_bookkeeping_count_subspaces,Auto.cj_bookkeeping_count_lines}
\leanok
An $m$-dimensional subspace has $k^m$ points.  The image of an
$r$-dimensional subspace under the coordinate map of an
$m$-dimensional subspace is an $r$-dimensional subspace.  Concatenating
subspaces of dimensions $r$ and $s$ gives a subspace of dimension $r+s$.
There are at most $(k+m)^n$ $m$-dimensional subspaces of $[k]^n$ and
at most $(k+1)^n$ lines.
\end{lemma}
\begin{proof}
Each variable set is nonempty, so its common letter can be read from
any one of its coordinates.  Thus the coordinate map is injective.
Under composition, each new variable set is a union of distinct old
variable sets; these unions remain nonempty and disjoint.  Fixed
parameters simply turn the corresponding old variable sets into fixed
coordinates.  Concatenation places the two collections of variable
sets in disjoint coordinate intervals.  Finally, encode a subspace
by a word over $[k]\cup\{v_1,\ldots,v_m\}$, with $v_j$ marking its
$j$th variable set.  Each subspace has such an encoding among the
$(k+m)^n$ possible words.  Taking $m=1$ gives the line bound.
\end{proof}

\begin{lemma}[Binary density theorem]\label{cj:binary}
\lean{Auto.C_cj_binary,Auto.cj_binary}
\leanok
For $0<\alpha\le1$, set
\[
 H_2(\alpha)=2\lceil\alpha^{-2}\rceil.
\]
If $n\ge H_2(\alpha)$ and $A\subseteq[2]^n$ has
$\mu_{2,n}(A)\ge\alpha$, then $A$ contains a combinatorial line.
\end{lemma}
\begin{proof}
Identify a word $x$ with the set $S(x)$ of coordinates carrying the
letter $2$.  Two distinct words form a binary combinatorial line
exactly when one of their associated sets is contained in the other.
Indeed, the variable set of the line is the difference of the two
sets.  A line-free family therefore corresponds to an antichain
$\mathcal A$ of subsets of $[n]$.

Choose a uniformly random permutation of $[n]$ and take its chain of
initial segments, including the empty set and $[n]$.  The probability
that this chain contains a specified $j$-element set is
$\binom nj^{-1}$.  A chain meets an antichain at most once, so taking
expectations gives
\[
 \sum_{S\in\mathcal A}\binom n{|S|}^{-1}\le1,
 \qquad
 |\mathcal A|\le\binom n{\lfloor n/2\rfloor}.
\]
The binomial coefficients are maximal at the middle because
$\binom n{j+1}/\binom nj=(n-j)/(j+1)$.
For $n=2r$, their normalized maximum satisfies
\[
 \frac{\binom{2r}r}{2^{2r}}
 =\prod_{j=1}^r\left(1-\frac1{2j}\right)
 \le \exp\left(-\frac12\sum_{j=1}^r\frac1j\right)
 \le(r+1)^{-1/2}.
\]
Here $1-t\le e^{-t}$ and
$\sum_{j=1}^r j^{-1}\ge\int_1^{r+1}t^{-1}\dd t$; the displayed
bound is also valid when $r=0$.  For $n=2r+1$, the normalized maximum
is $(2r+1)/(2r+2)$ times that for $2r$, and is no larger.
Consequently a line-free subset of $[2]^n$ has density at most
$(\lfloor n/2\rfloor+1)^{-1/2}$.  When
$n\ge2\lceil\alpha^{-2}\rceil$, this last quantity is strictly less
than $\alpha$, a contradiction.
\end{proof}

\begin{lemma}[Multidimensional binary density theorem]\label{cj:binary-multi}
\lean{Auto.C_cj_binary_multi,Auto.cj_binary_multi}
\leanok
For every integer $m\ge1$ and $0<\alpha\le1$, there is an integer
$B(m,\alpha)$ such that every $A\subseteq[2]^n$ with
$n\ge B(m,\alpha)$ and $\mu_{2,n}(A)\ge\alpha$ contains an
$m$-dimensional subspace.  One possible recursive choice is
\[
 B(1,\alpha)=H_2(\alpha),
\]
\[
 \begin{split}
 M&=H_2(\alpha/2),\\
 B(m+1,\alpha)&=M+B\left(m,\frac{\alpha}{2\,3^M}\right).
 \end{split}
\]
\end{lemma}
\begin{proof}
The case $m=1$ is Lemma~\ref{cj:binary}.  Suppose the assertion is
known for $m$, and split a word of length $n$ into a prefix $x$ of
length $n-M$ and a suffix of length $M$.  The average suffix-fiber
density is at least $\alpha$.  The prefixes whose fibers have density
at least $\alpha/2$ have density at least $\alpha/2$: if their density
is $p$, the average is at most $p+\alpha/2$.
Every such fiber contains a line by the choice of $M$.
Choose one line from each of these fibers.  There are at most $3^M$
possible lines by Lemma~\ref{cj:bookkeeping}, so some fixed line
$L\subseteq[2]^M$ lies in every fiber above a set $C$ of prefixes of
density at least $\alpha/(2\,3^M)$.  The induction hypothesis gives
an $m$-dimensional subspace $U\subseteq C$.  Then
$U\frown L\subseteq A$ is an $(m+1)$-dimensional subspace.
\end{proof}

\begin{lemma}[Uniform lower bounds for prefix fibers]\label{cj:fiber}
\lean{Auto.C_cj_fiber,Auto.cj_fiber}
\leanok
For $k\ge2$, $m\ge1$, and $0<\varepsilon<1$, define
\[
 P(k,m,\varepsilon)
 =m\left(\left\lceil\frac{k^m-1}{\varepsilon}\right\rceil+2\right).
\]
If $n\ge P(k,m,\varepsilon)$ and $A\subseteq[k]^n$, there are
$\ell<n$ and an $m$-dimensional subspace $U\subseteq[k]^\ell$ such
that
\[
 \mu_{k,n-\ell}(A_u)\ge\mu_{k,n}(A)-\varepsilon
 \qquad(u\in U).
\]
\end{lemma}
\begin{proof}
Put $\alpha=\mu_{k,n}(A)$, $K=k^m$, and
$\rho=\varepsilon/(K-1)$; let $J=\lceil1/\rho\rceil$.
We construct a fixed prefix a block of $m$ coordinates at a time.
Initially the prefix is empty and its suffix fiber has density
$\alpha$.  At stage $j$, suppose the chosen prefix has length $jm$
and its remaining fiber has density $a_j\ge\alpha+j\rho$.
Split that remaining fiber according to its next $m$ coordinates.
There are $K$ child fibers whose densities average to $a_j$.
If every child has density at least $\alpha-\varepsilon$, the desired
$U$ is obtained by appending all $[k]^m$ to the current prefix, and
we stop.

Otherwise one child has density less than
$\alpha-\varepsilon\le a_j-\varepsilon$.  The sum of the other
$K-1$ densities is therefore greater than
$(K-1)a_j+\varepsilon$.  At least one of those children has density
greater than $a_j+\rho$.  Select its $m$-letter word as the next
fixed block.  This establishes the induction hypothesis for stage
$j+1$.

The process cannot fail to stop at all stages $0,\ldots,J$, since
$J+1$ failures would produce a fiber of density greater than
$\alpha+(J+1)\rho>1$.  The assumed dimension permits all these
stages: at stage $j\le J$ the length
$\ell=(j+1)m$ is at most $(J+1)m<n$.
Thus the terminal prefix subspace and all of its suffix fibers are
well defined and have the stated property.
\end{proof}

\begin{lemma}[A binary face in a dense ternary cube]\label{cj:binary-face}
\lean{Auto.C_cj_binary_face,Auto.cj_binary_face}
\leanok
For $d\ge1$ and $0<\alpha\le1$, put
\[
 M=B(d,\alpha/2),\qquad
 F_2(d,\alpha)=P(3,M,\alpha/2).
\]
If $n\ge F_2(d,\alpha)$ and $A\subseteq[3]^n$ has
$\mu_{3,n}(A)\ge\alpha$, there is a $d$-dimensional ternary
subspace $V$ such that $V^{(2)}\subseteq A$.
\end{lemma}
\begin{proof}
By Lemma~\ref{cj:fiber}, there are $\ell<n$ and an
$M$-dimensional $U\subseteq[3]^\ell$ such that every suffix fiber
above $U$ has density at least $\alpha/2$.  The same is true above
$U^{(2)}$.  Averaging first over $u\in U^{(2)}$ and then over suffixes
shows that some suffix $y$ satisfies
\[
 \frac{|\{u\in U^{(2)}:u\frown y\in A\}|}{2^M}\ge\alpha/2.
\]
Via the coordinate map of $U$, this is a subset of $[2]^M$ of density
at least $\alpha/2$.  Lemma~\ref{cj:binary-multi} gives a
$d$-dimensional binary subspace in this subset.  Its encoding uses
fixed letters $1,2$ and $d$ variables.  Allow those variables also
to take the letter $3$, compose with the coordinate map of $U$, and
append $y$.  Lemma~\ref{cj:bookkeeping} shows that the result is a
$d$-dimensional ternary subspace, and its binary face is contained
in $A$ by construction.
\end{proof}

\begin{definition}[Insensitive sets]\label{cj:insensitive}
\lean{Auto.InsEquiv,Auto.IsInsensitive,Auto.IsInsensitiveIn}
For $i\in\{1,2\}$, two words in $[3]^n$ are
\emph{$(i,3)$-equivalent} if, at each coordinate, either their letters
are equal or both letters belong to $\{i,3\}$.  A set
$D\subseteq[3]^n$ is \emph{$(i,3)$-insensitive} if it is a union of
these equivalence classes.  A subset of a subspace is called
$(i,3)$-insensitive in that subspace when its inverse image under the
coordinate map is $(i,3)$-insensitive.
\end{definition}

\begin{lemma}[Elementary properties of insensitivity]\label{cj:insensitive-facts}
\lean{Auto.cj_insensitive_facts_union,Auto.cj_insensitive_facts_inter,Auto.cj_insensitive_facts_compl,Auto.cj_insensitive_facts_fix,Auto.cj_insensitive_facts_preimage,Auto.cj_insensitive_facts_face}
\leanok
Unions, intersections, and complements of $(i,3)$-insensitive sets
are $(i,3)$-insensitive.  Fixing any collection of coordinates
preserves insensitivity on the remaining coordinates.  The inverse
image of an $(i,3)$-insensitive set under any subspace coordinate map
is $(i,3)$-insensitive.  Finally, if $D$ is $(i,3)$-insensitive and
$V^{(2)}\subseteq D$, then $V\subseteq D$.
\end{lemma}
\begin{proof}
The first assertions follow because membership is constant on each
equivalence class.  Fixing coordinates restricts that equivalence
relation to the unfixed coordinates.  For a subspace coordinate map,
changing a parameter between $i$ and $3$ changes all coordinates of
its variable set between those same two letters, so equivalent
parameter words have equivalent images.
For the final assertion, take a point of $V$ and replace every
parameter equal to $3$ by $i$.  The resulting point lies in
$V^{(2)}$ and is $(i,3)$-equivalent to the original point.  Its
membership in $D$ therefore implies membership of the original point.
\end{proof}

\begin{lemma}[Tiling one insensitive set]\label{cj:tile-one}
\lean{Auto.C_cj_tile_one,Auto.cj_tile_one}
\leanok
For every $d\ge1$ and $0<\varepsilon\le1$ there is an integer
$T_1(d,\varepsilon)$ with the following property.  If
$n\ge T_1(d,\varepsilon)$, $i\in\{1,2\}$, and
$D\subseteq[3]^n$ is $(i,3)$-insensitive, there is a finite family
$\mathcal V$ of pairwise disjoint $d$-dimensional subspaces, all
contained in $D$, such that
\[
 \mu_{3,n}\left(D\setminus\bigcup_{V\in\mathcal V}V\right)
 \le\varepsilon.
\]
A valid choice is obtained from
\[
 \beta=\varepsilon/2,\quad M=F_2(d,\beta),\quad Q=(3+d)^M,
 \quad \vartheta=\frac{\beta}{Q}\,3^{d-M}
\]
by setting
\[
 L=\lceil\vartheta^{-1}\rceil+1,
 \qquad T_1(d,\varepsilon)=LM.
\]
\end{lemma}
\begin{proof}
We remove subspaces greedily, using successive blocks of $M$
coordinates from right to left.  The invariant is the following:
after $t$ blocks have been used, the unremoved set $R_t\subseteq D$
has the property that, for every suffix $y\in[3]^{tM}$, its prefix
section
\[
 (R_t)^y=\{x\in[3]^{n-tM}:x\frown y\in R_t\}
\]
is $(i,3)$-insensitive on its $n-tM$ coordinates.  This holds for
$t=0$, with $R_0=D$.

If $\mu_{3,n}(R_t)\le\varepsilon$, stop.  Otherwise assume $t<L$
and write a word as $x\frown z\frown y$, with respective lengths
$n-(t+1)M$, $M$, and $tM$.  For a pair $(x,y)$, put
\[
 E_{x,y}=\{z\in[3]^M:x\frown z\frown y\in R_t\}.
\]
The invariant and Lemma~\ref{cj:insensitive-facts} imply that
$E_{x,y}$ is $(i,3)$-insensitive.  Its average density, as $(x,y)$
ranges uniformly over the indicated product, is
$\mu_{3,n}(R_t)>2\beta$.  Hence the pairs with
$\mu_{3,M}(E_{x,y})\ge\beta$ have product density at least $\beta$:
if that density is $p$, the average is at most $p+\beta$.

For each such pair, Lemma~\ref{cj:binary-face} supplies a
$d$-dimensional subspace whose binary face lies in $E_{x,y}$.
Lemma~\ref{cj:insensitive-facts} puts the entire subspace in
$E_{x,y}$.  There are at most $Q$ possible subspaces by
Lemma~\ref{cj:bookkeeping}.  Consequently some fixed
$d$-dimensional $W\subseteq[3]^M$ has
\[
 S=\{(x,y):x\frown W\frown y\subseteq R_t\}
\]
of product density at least $\beta/Q$.  Remove every set
$x\frown W\frown y$ for $(x,y)\in S$.  These sets are pairwise
disjoint $d$-dimensional subspaces: different pairs have different
fixed prefix or suffix, and $W$ varies only in the middle block.
They are all in $R_t$, so they are disjoint from all subspaces removed
at earlier stages.  The density removed at this stage is
\[
 \frac{|S|3^d}{3^n}
 =\frac{|S|}{3^{n-M}}\,3^{d-M}
 \ge\vartheta.
\]

It remains to check the invariant; the residual need not be
insensitive on \emph{all} its coordinates.  For a fixed old suffix
$y$, define
\[
 S_y=\{x:(x,y)\in S\}
 =\bigcap_{z\in W}\{x:x\frown z\frown y\in R_t\}.
\]
Each set in this intersection is an insensitive prefix section by
the induction hypothesis, so $S_y$ is insensitive.  For any fixed
new suffix $z\frown y$, the section of $R_{t+1}$ is
\[
 \{x:x\frown z\frown y\in R_{t+1}\}
 =\begin{cases}
   \{x:x\frown z\frown y\in R_t\}\setminus S_y,&z\in W,\\
   \{x:x\frown z\frown y\in R_t\},&z\notin W.
  \end{cases}
\]
Both alternatives are $(i,3)$-insensitive.  This proves the invariant
for the next stage, including when the remaining prefix has length
zero.

All $L$ stages would be possible dimensionally, since $n\ge LM$.
If the algorithm performed $L$ removal stages without stopping, it
would remove total density at least $L\vartheta>1$, which is
impossible.  It therefore stops with residual density at most
$\varepsilon$.  The collection of all removed subspaces is the
required finite family.  If the initial density is already at most
$\varepsilon$, the empty family suffices.
\end{proof}

\begin{lemma}[Tiling an intersection of two insensitive sets]\label{cj:tile-two}
\lean{Auto.C_cj_tile_two,Auto.cj_tile_two}
\leanok
For $d\ge1$ and $0<\varepsilon\le1$, put
\[
 M=T_1(d,\varepsilon/2),\qquad
 T_2(d,\varepsilon)=T_1(M,\varepsilon/2).
\]
If $n\ge T_2(d,\varepsilon)$, $D_1\subseteq[3]^n$ is
$(1,3)$-insensitive, and $D_2\subseteq[3]^n$ is
$(2,3)$-insensitive, then $D=D_1\cap D_2$ contains a finite family
$\mathcal V$ of pairwise disjoint $d$-dimensional subspaces such that
\[
 \mu_{3,n}\left(D\setminus\bigcup_{V\in\mathcal V}V\right)
 \le\varepsilon.
\]
\end{lemma}
\begin{proof}
Apply Lemma~\ref{cj:tile-one} to $D_1$, with target dimension $M$
and error $\varepsilon/2$.  This gives pairwise disjoint
$M$-dimensional subspaces $W\subseteq D_1$ whose union misses at
most $\varepsilon/2$ of the ambient density of $D_1$.
Inside each $W$, the pullback of $D_2\cap W$ is
$(2,3)$-insensitive by Lemma~\ref{cj:insensitive-facts}.
Its ambient dimension is $M=T_1(d,\varepsilon/2)$, so a second
application of Lemma~\ref{cj:tile-one} tiles it by $d$-dimensional
subspaces, leaving relative density at most $\varepsilon/2$ in $W$.

All resulting subspaces are contained in $D_1\cap D_2$ and are
pairwise disjoint, since the outer subspaces $W$ are disjoint.
The part of $D$ outside their union consists of a subset of the
uncovered portion of $D_1$, together with the residuals inside the
outer subspaces.  Its ambient density is therefore at most
\[
 \frac{\varepsilon}{2}
 +\sum_W\mu_{3,n}(W)\,\frac{\varepsilon}{2}
 \le\varepsilon.
\]
This proves the assertion.
\end{proof}



\subsection{Equal-slices averaging and the three-letter theorem}

The remaining argument specializes the density-increment method of
\cite{Polymath} to three letters.  Its two changes of probability measure
are proved below by direct counting.  Thus no Ramsey theorem for
parameter sets, and no density Hales--Jewett theorem, is used as an input.

\begin{definition}[Equal-slices probability and total variation]\label{cj:equal-slices}
\lean{Auto.equalSlices,Auto.dTV}
Retain the uniform probability $\mu_{k,n}$ from
Definition~\ref{cj:def-words}.  For $k\in\{2,3\}$ and $n\ge0$,
the \emph{equal-slices probability} $\nu_{k,n}$ is obtained
by first choosing uniformly a vector $(a_1,\ldots,a_k)$ of nonnegative
integers with $a_1+\cdots+a_k=n$, and then choosing uniformly a word
having exactly $a_j$ occurrences of the letter $j$ for every $j$.
For $n=0$ both measures give mass one to the empty word.

For probabilities $P,Q$ on a finite set $X$, put
\[
 d_{\mathrm{TV}}(P,Q)=\frac12\sum_{x\in X}|P(\{x\})-Q(\{x\})|.
\]
For integers $a\ge0$ and $r\ge0$, the falling factorial is
$(a)_r=a(a-1)\cdots(a-r+1)$, with $(a)_0=1$ and $(a)_r=0$ if
$r>a$.
\end{definition}

\begin{lemma}[Elementary counting and total variation]\label{cj:probability-facts}
\lean{Auto.cj_probability_facts_atom,Auto.cj_probability_facts_binary_face_mass,Auto.cj_probability_facts_tv_function,Auto.cj_probability_facts_mixture}
\leanok
If $x\in[k]^n$ has $a_j$ occurrences of $j$, then
\begin{equation}\label{eq:cj-equal-atom}
 \nu_{k,n}(\{x\})
 =\binom{n+k-1}{k-1}^{-1}\binom{n}{a_1,\ldots,a_k}^{-1}
 =\frac{(k-1)!\,a_1!\cdots a_k!}{(n+k-1)!}.
\end{equation}
In particular,
\begin{equation}\label{eq:cj-binary-face-mass}
 \nu_{3,n}([2]^n)=\frac{2}{n+2}.
\end{equation}
For any $0\le h\le1$ on a finite set,
\begin{equation}\label{eq:cj-tv-function}
 |\mathbb E_P h-\mathbb E_Qh|\le d_{\mathrm{TV}}(P,Q).
\end{equation}
Moreover, if $P=\sum_j a_jP_j$, where $a_j\ge0$ and $\sum_j a_j=1$,
then
\[
 d_{\mathrm{TV}}(P,Q)\le\sum_j a_jd_{\mathrm{TV}}(P_j,Q).
\]
\end{lemma}
\begin{proof}
A vector of $k$ nonnegative integers summing to $n$ is specified by
placing $k-1$ separators among a row of $n$ indistinguishable marks.
There are $\binom{n+k-1}{k-1}$ possibilities.  There are
$n!/(a_1!\cdots a_k!)$ words with a prescribed count vector, proving
\eqref{eq:cj-equal-atom}.  Exactly $n+1$ of the
$\binom{n+2}{2}$ count vectors for three letters have $a_3=0$;
this gives \eqref{eq:cj-binary-face-mass}.

The signed masses $P(\{x\})-Q(\{x\})$ sum to zero.  Their positive
sum and the absolute value of their negative sum are therefore both
$d_{\mathrm{TV}}(P,Q)$.  Multiplication by numbers in $[0,1]$ makes
the resulting sum lie between the negative and positive sums.  This
proves \eqref{eq:cj-tv-function}.  The last assertion follows by
applying the triangle inequality inside the defining finite sum.
\end{proof}

\begin{definition}[Random coordinate restriction]\label{cj:restriction}
\lean{Auto.restrictMap,Auto.restrictPre}
Let $n\ge r$.  Choose an injection $\iota:[r]\longrightarrow[n]$
uniformly among the $(n)_r$ injections.  For a word $y$ on the
complement of its image, define
\[
 \phi_{\iota,y}:[k]^r\longrightarrow[k]^n
\]
by putting the $j$th input letter in coordinate $\iota(j)$ and putting
the prescribed letters $y$ in all remaining coordinates.  Its image
is a coordinate subspace, hence a combinatorial subspace in the sense
of Definition~\ref{cj:def-words}.  For $A\subseteq[k]^n$, write
$A_{\iota,y}=\phi_{\iota,y}^{-1}(A)$.
\end{definition}

\begin{lemma}[Inserting a fixed number of letters into uniform background]
\label{cj:transfer-uniform}
\lean{Auto.cj_transfer_uniform}
\leanok
Fix $k\in\{2,3\}$, $r\ge1$, and $\zeta>0$.  There is
$N=N(k,r,\zeta)\ge r$ such that the following holds for every $n\ge N$.
Choose $\iota$ as in Definition~\ref{cj:restriction}, give the
remaining coordinates independent uniform letters, and independently
choose $w\in[k]^r$ with any probability distribution $\lambda$.
The law of $\phi_{\iota,y}(w)$ is within $\zeta$ in total variation
of $\mu_{k,n}$.  The same $N$ works for every $\lambda$.
\end{lemma}
\begin{proof}
First take $\lambda$ to be concentrated on a single word $w$, and
let $c_j$ be the number of its letters equal to $j$.  For an output
word $z$, let $a_j$ be the number of its letters equal to $j$.
There are $\prod_j(a_j)_{c_j}$ injections whose specified letters
agree with $z$.  Once the injection has been chosen, the probability
that the remaining uniform letters agree with $z$ is $k^{-(n-r)}$.
The density of the output law with respect to $\mu_{k,n}$ is consequently
\begin{equation}\label{eq:cj-uniform-likelihood}
 L_{n,w}(z)=k^r\frac{\prod_{j=1}^k(a_j)_{c_j}}{(n)_r}.
\end{equation}
It lies in $[0,k^r]$, because the injections counted in the numerator
are some of the $(n)_r$ possible injections.

For a uniform random $z$, the random variable $a_j$ is the sum of
$n$ independent indicators, each having mean $1/k$.  Expanding the
square and using independence gives
\[
 \mathbb E a_j=\frac nk,
 \qquad
 \mathbb E\left(a_j-\frac nk\right)^2
       =n\frac1k\left(1-\frac1k\right).
\]
For every $\tau>0$, the elementary inequality
$\one_{\{|X|>\tau\}}\le |X|^2/\tau^2$ therefore yields
\[
 \mathbb P\left(\left|\frac{a_j}{n}-\frac1k\right|>\tau\right)
 \le\frac{k-1}{k^2n\tau^2}.
\]
Thus all $a_j/n$ approach $1/k$ in probability.  Since $r$ and the
$c_j$ are fixed, divide the numerator and denominator of
\eqref{eq:cj-uniform-likelihood} by $n^r$ to see that $L_{n,w}$
approaches $1$ in probability.  This also follows directly by writing
each factor as $a_j/n-s/n$ and each denominator factor as $1-s/n$.
For any $u>0$,
\[
 \mathbb E_{\mu_{k,n}}|L_{n,w}-1|
 \le u+(k^r+1)\,
       \mathbb P_{\mu_{k,n}}(|L_{n,w}-1|>u).
\]
First let $n$ increase and then let $u$ decrease to zero.  It follows
that the total variation distance, which is one half of this
expectation, tends to zero.

There are only $k^r$ choices of $w$.  Choose a single $N$ that works
for every one of them.  Any distribution $\lambda$ gives a mixture
of these output laws, so the mixture inequality of
Lemma~\ref{cj:probability-facts} proves the assertion.
\end{proof}

\begin{lemma}[Inserting letters into equal-slices background]
\label{cj:transfer-equal}
\lean{Auto.cj_transfer_equal,Auto.cj_transfer_equal_average}
\leanok
Let $m\ge r\ge1$.  Choose $\iota:[r]\to[m]$ uniformly, choose
$y$ on the remaining $m-r$ coordinates with probability
$\nu_{3,m-r}$, and independently choose $w\in[3]^r$ with any
distribution.  If $P$ is the law of $\phi_{\iota,y}(w)$, then
\begin{equation}\label{eq:cj-equal-transfer}
 d_{\mathrm{TV}}(P,\nu_{3,m})
 \le 1-\frac{(m-r+1)(m-r+2)}{(m+1)(m+2)}
 \le\frac{2r}{m+1}\le\frac{3r}{m}.
\end{equation}
In particular, if $w$ is uniform, then for every $H\subseteq[3]^m$,
\begin{equation}\label{eq:cj-average-uniform-slice}
 \left|\mathbb E_{\iota,y}
       \mu_{3,r}(H_{\iota,y})-\nu_{3,m}(H)\right|
 \le\frac{3r}{m}.
\end{equation}
\end{lemma}
\begin{proof}
Fix $w$, with letter counts $c_1,c_2,c_3$, and let $z$ have
letter counts $a_1,a_2,a_3$.  If any $a_j<c_j$, the output cannot
be $z$.  Otherwise its probability is
\[
 \frac{\prod_{j=1}^3(a_j)_{c_j}}{(m)_r}
 \frac{2!\,\prod_{j=1}^3(a_j-c_j)!}{(m-r+2)!}.
\]
Divide this by the equal-slices atom from
\eqref{eq:cj-equal-atom}.  All the factorials depending on the
$a_j$ cancel.  The resulting density is
\[
 L_w(z)=C_{m,r}\,\one_{\{a_j\ge c_j\text{ for all }j\}},
 \qquad
 C_{m,r}=\frac{(m+2)_r}{(m)_r}
 =\frac{(m+1)(m+2)}{(m-r+1)(m-r+2)}.
\]
Since this density has integral one, the indicated event has
$\nu_{3,m}$-probability $C_{m,r}^{-1}$.  Its total variation distance
from $\nu_{3,m}$ is exactly $1-C_{m,r}^{-1}$: the positive part of
$L_w-1$ and the negative part have that same integral.
The bound for arbitrary $w$-distributions follows by taking mixtures.
Finally,
\[
 1-C_{m,r}^{-1}
   =\frac{r(2m-r+3)}{(m+1)(m+2)}
   \le\frac{2r}{m+1}.
\]
Equation~\eqref{eq:cj-average-uniform-slice} is
\eqref{eq:cj-tv-function} applied to the indicator of $H$.
\end{proof}

\begin{definition}[Collapsing the third letter]\label{cj:collapse}
\lean{Auto.collapse,Auto.collapseSet,Auto.collapseBoth,Auto.binaryTrace}
For $i\in\{1,2\}$ and $x\in[3]^m$, let $x^{3\to i}\in[2]^m$
be obtained by replacing every $3$ by $i$.  For $B\subseteq[2]^m$,
put
\[
 C_i(B)=\{x\in[3]^m:x^{3\to i}\in B\},
 \qquad C(B)=C_1(B)\cap C_2(B).
\]
\end{definition}

\begin{lemma}[A quantitative supply of binary pairs]\label{cj:binary-pairs}
\lean{Auto.cj_binary_pairs_insensitive,Auto.cj_binary_pairs_bound,Auto.cj_binary_pairs_forbidden}
\leanok
Let $B\subseteq[2]^m$ and $b=\nu_{2,m}(B)$.  Each $C_i(B)$ is
$(i,3)$-insensitive, and
\begin{equation}\label{eq:cj-pair-bound}
 \nu_{3,m}(C(B))\ge\frac{b^2}{2}.
\end{equation}
If $A\subseteq[3]^m$ contains no combinatorial line and
$B=A\cap[2]^m$, then
\begin{equation}\label{eq:cj-forbidden-third}
 A\cap C(B)\subseteq[2]^m.
\end{equation}
\end{lemma}
\begin{proof}
Changing any occurrences of $i$ and $3$ leaves $x^{3\to i}$
unchanged.  This proves the asserted insensitivity.

Choose a uniform permutation $\pi$ of $[m]$.  For $0\le s\le m$,
let $u_s(\pi)$ have letter $1$ at
$\pi(1),\ldots,\pi(s)$ and letter $2$ elsewhere.  If $s$ is chosen
uniformly from $\{0,\ldots,m\}$, then $u_s(\pi)$ has law
$\nu_{2,m}$: its number of ones is uniform, and conditional on that
number their positions are uniform.  Set
\[
 h(\pi)=\frac1{m+1}\sum_{s=0}^m\one_B(u_s(\pi)).
\]
Then $\mathbb E_\pi h=b$.  For independent uniform ranks $s,t$,
\[
 \mathbb P(u_s(\pi),u_t(\pi)\in B)
 =\mathbb E_\pi h(\pi)^2\ge b^2,
\]
because $\mathbb E(h-b)^2\ge0$.

Put $a=\min(s,t)$ and $c=\max(s,t)$.  Make a ternary word $z$
by giving the first $a$ positions in the order $\pi$ letter $1$,
the next $c-a$ positions letter $3$, and the remaining positions
letter $2$.  Its two collapsed words are precisely $u_a(\pi)$ and
$u_c(\pi)$.  Hence the probability that $z\in C(B)$ is at least
$b^2$.

For $a<c$, the count vector $(a,m-c,c-a)$ has probability
$2/(m+1)^2$ under this construction; for $a=c$ it has probability
$1/(m+1)^2$.  Under $\nu_{3,m}$ every count vector has probability
$2/((m+1)(m+2))$.  Conditional on a count vector, both constructions
are uniform on its words.  Therefore the probability of every word
under the constructed law is at most
$(m+2)/(m+1)$ times its equal-slices probability.  Consequently
\[
 \nu_{3,m}(C(B))\ge\frac{m+1}{m+2}b^2\ge\frac{b^2}{2}.
\]
Finally, if $z\in A\cap C(B)$ has at least one $3$, the words
$z^{3\to1},z^{3\to2},z$ form a combinatorial line, with variable
set equal to the positions of the $3$'s.  All three words belong
to $A$, a contradiction.  This proves
\eqref{eq:cj-forbidden-third}.
\end{proof}

\begin{proposition}[A three-letter density increment]
\label{cj:density-increment}
\lean{Auto.C_cj_density_increment,Auto.cj_density_increment}
\leanok
For every $0<\delta\le1$ there is $h(\delta)>0$ with the following
property.  For every integer $d\ge1$ there is an integer
$N(d,\delta)\ge d$ such that, whenever $n\ge N(d,\delta)$ and
$A\subseteq[3]^n$ has uniform density $\alpha\ge\delta$ and
contains no combinatorial line, some $d$-dimensional combinatorial
subspace $V$ satisfies
\[
 \frac{|A\cap V|}{|V|}\ge\alpha+h(\delta).
\]
One may take $h(\delta)=\delta^4/65536$.
\end{proposition}
\begin{proof}
Fix $\delta$ and $d$.  Set
\begin{equation}\label{eq:cj-increment-constants}
 \theta=\frac{\delta^2}{64},\qquad
 t=\frac{\delta\theta}{8},\qquad
 \kappa=\frac t8,\qquad
 \varepsilon=\frac{\delta\kappa}{8},\qquad
 h=\frac\varepsilon2=\frac{\delta^4}{65536}.
\end{equation}
In particular, $\kappa\le\delta/4$ and
$\varepsilon\le\delta/4$.

Use Lemma~\ref{cj:tile-two} to choose $r\ge d$ so that every
intersection of a $(1,3)$-insensitive set and a $(2,3)$-insensitive
set in $[3]^r$ can be covered by disjoint $d$-dimensional
combinatorial subspaces contained in the intersection, with uncovered
uniform measure at most $t/8$.  Choose $m\ge3r$ so large that
\begin{equation}\label{eq:cj-dimension-choices}
 \frac{3r}{m}\le\min\left(\frac\varepsilon2,\frac t8\right),
 \qquad
 \frac{2}{m+2}\le\frac{\delta\theta}{2}.
\end{equation}
Finally choose $N(d,\delta)\ge m$ so that
Lemma~\ref{cj:transfer-uniform}, with $k=3$, $r=m$ and error
$\varepsilon$, holds for all $n\ge N(d,\delta)$.
All these choices depend on $d$ and $\delta$, and not on $A$ or
its actual density $\alpha$.

Let $n\ge N(d,\delta)$ and let $A$ satisfy the hypotheses.
Choose a random coordinate restriction $\phi:[3]^m\to[3]^n$
with uniform letters in the fixed coordinates, and put
\[
 a(\phi)=\nu_{3,m}(\phi^{-1}(A)),\qquad
 b(\phi)=\nu_{2,m}(\phi^{-1}(A)\cap[2]^m).
\]
Apply Lemma~\ref{cj:transfer-uniform} first with inserted-word
law $\nu_{3,m}$ and then with inserted-word law $\nu_{2,m}$,
viewed as a probability on $[3]^m$ supported on $[2]^m$.
It gives
\begin{equation}\label{eq:cj-two-averages}
 \mathbb E a\ge\alpha-\varepsilon,
 \qquad
 \mathbb E b\ge\alpha-\varepsilon.
\end{equation}

Suppose first that some restriction $\phi$ has
$a(\phi)\ge\alpha+\varepsilon$.  Inside its copy of $[3]^m$,
choose a random $r$-coordinate restriction with equal-slices
background, as in Lemma~\ref{cj:transfer-equal}.  By
\eqref{eq:cj-average-uniform-slice} and
\eqref{eq:cj-dimension-choices}, the average uniform density of
$\phi^{-1}(A)$ on these $r$-dimensional subspaces is at least
$\alpha+\varepsilon/2$.  One of them has at least this density.
Partition that subspace by fixing $r-d$ of its free coordinates.
The resulting $d$-dimensional subspaces have equal size, so one
has density at least $\alpha+\varepsilon/2=\alpha+h$.
Composing its coordinate maps with $\phi$ gives the required
subspace of the original cube.

It remains to handle the case in which every restriction satisfies
$a<\alpha+\varepsilon$.  Let
$p=\mathbb P(a<\alpha-\kappa)$.  From
\eqref{eq:cj-two-averages} and the upper bound on $a$,
\[
 \alpha-\varepsilon
 \le\mathbb E a
 \le p(\alpha-\kappa)+(1-p)(\alpha+\varepsilon).
\]
It follows that
\[
 p\le\frac{2\varepsilon}{\kappa+\varepsilon}
      \le\frac{2\varepsilon}{\kappa}=\frac\delta4.
\]
On the other hand, if $q=\mathbb P(b\ge\delta/4)$, then
$0\le b\le1$ and \eqref{eq:cj-two-averages} give
\[
 \delta-\varepsilon
 \le\mathbb E b\le q+(1-q)\frac\delta4,
 \qquad
 q\ge\frac{3\delta/4-\varepsilon}{1-\delta/4}
       \ge\frac\delta2.
\]
Thus there is a restriction for which
\begin{equation}\label{eq:cj-good-face}
 \nu_{3,m}(A')\ge\alpha-\kappa,
 \qquad
 \nu_{2,m}(A'\cap[2]^m)\ge\frac\delta4,
 \qquad A'=\phi^{-1}(A).
\end{equation}
A combinatorial line in $A'$ would map to one in $A$, so $A'$ is
also line-free.

Set $B=A'\cap[2]^m$ and
$C_i=C_i(B)$, $C=C_1\cap C_2$ as in
Definition~\ref{cj:collapse}.  Lemma~\ref{cj:binary-pairs} and
\eqref{eq:cj-good-face} imply
\[
 \nu_{3,m}(C)\ge\frac{(\delta/4)^2}{2}
 =\frac{\delta^2}{32}\ge\theta.
\]
Since $A'$ is line-free,
\eqref{eq:cj-forbidden-third} gives
\begin{equation}\label{eq:cj-small-intersection}
 \nu_{3,m}(A'\cap C)
 \le\frac{2}{m+2}
 \le\frac{\delta\theta}{2}
 \le\frac\alpha2\nu_{3,m}(C).
\end{equation}
For the next calculation write $\nu=\nu_{3,m}$ and
$\beta=\alpha-\kappa$.  We have $0<\beta\le1$.  The total
$\beta$-excess of $A'$ is nonnegative by
\eqref{eq:cj-good-face}.  Removing $C$ and using
\eqref{eq:cj-small-intersection}, we obtain
\begin{align*}
 \nu(A'\cap C^c)-\beta\nu(C^c)
 &=[\nu(A')-\beta]-[\nu(A'\cap C)-\beta\nu(C)]\\
 &\ge(\alpha/2-\kappa)\nu(C)
 \ge\frac{\delta\theta}{4}=2t.
\end{align*}
The complement of $C$ is the disjoint union
\[
 D^{(1)}=C_1^c,
 \qquad
 D^{(2)}=C_1\cap C_2^c.
\]
Each of these is an intersection of a $(1,3)$-insensitive set and
a $(2,3)$-insensitive set; for $D^{(1)}$ the second set is the whole
cube.  Because the excesses add over this partition, one of the two
sets, denoted by $D$, satisfies
\begin{equation}\label{eq:cj-structured-excess}
 \nu(A'\cap D)-\beta\nu(D)\ge t.
\end{equation}

Choose a random $r$-coordinate restriction $\psi:[3]^r\to[3]^m$
with equal-slices background.  Write $A''=\psi^{-1}(A')$ and
$D'=\psi^{-1}(D)$ for its two traces.  Applying
Lemma~\ref{cj:transfer-equal} separately to $A'\cap D$ and $D$,
and using $\beta\le1$, shows that the expected value of
\[
 \mu_{3,r}(A''\cap D')-\beta\mu_{3,r}(D')
\]
is at least $t-2(3r/m)\ge3t/4$.  Fix a restriction for which
this bound holds.  The two insensitivities defining $D$ are
preserved by coordinate restriction, so $D'$ has the required
form for the tiling lemma.  Since $\kappa=t/8$,
\begin{equation}\label{eq:cj-uniform-excess}
 \mu_{3,r}(A''\cap D')-\alpha\mu_{3,r}(D')
 \ge\frac{3t}{4}-\kappa\mu_{3,r}(D')
 \ge\frac{5t}{8}\ge\frac t2.
\end{equation}

By the choice of $r$, there are disjoint $d$-dimensional subspaces
$V_1,\ldots,V_s$ contained in $D'$ whose union $U$ leaves
$\mu_{3,r}(D'\setminus U)\le t/8$.  Subtracting the contribution
of this remainder from \eqref{eq:cj-uniform-excess} gives
\[
 \mu_{3,r}(A''\cap U)-\alpha\mu_{3,r}(U)\ge\frac{3t}{8}.
\]
In particular $U$ is nonempty.  The left side is
\[
 \sum_{j=1}^s\mu_{3,r}(V_j)
       \left(\frac{|A''\cap V_j|}{|V_j|}-\alpha\right).
\]
As $\sum_j\mu_{3,r}(V_j)=\mu_{3,r}(U)\le1$, at least one tile
has relative density at least $\alpha+3t/8$.  This is at least
$\alpha+h$, because $h=\delta t/128\le3t/8$.
The image of that tile under $\phi\circ\psi$ is a
$d$-dimensional combinatorial subspace of $[3]^n$ with the same
relative density.  This proves the proposition.
\end{proof}

\begin{theorem}[Density Hales--Jewett for three letters]\label{cj:dhj3}
\lean{Auto.cj_dhj3}
\leanok
For every $0<\delta\le1$ there is an integer $H_3(\delta)$ such
that, for every $n\ge H_3(\delta)$, each set $A\subseteq[3]^n$
with $|A|\ge\delta3^n$ contains a combinatorial line.
\end{theorem}
\begin{proof}
Fix $\delta$ and let $h=h(\delta)>0$ be supplied by
Proposition~\ref{cj:density-increment}.  Choose an integer $K$
with $Kh>1$, put $n_0=1$, and recursively choose
\[
 n_{j+1}=N(n_j,\delta)\qquad(0\le j<K).
\]
Set $H_3(\delta)=n_K$.  First suppose $n=n_K$ and that $A$ is
line-free with density at least $\delta$.  Apply the proposition
with target dimension $n_{K-1}$.  Identifying the resulting
subspace with $[3]^{n_{K-1}}$, the trace of $A$ is still line-free
and its density has increased by at least $h$.  Apply the
proposition again with target dimension $n_{K-2}$.  Continue in
this way.  Composing the coordinate descriptions of combinatorial
subspaces preserves their dimensions and sends their combinatorial
lines to combinatorial lines of the original cube.  Therefore all
$K$ applications are legitimate, and the final trace has density
at least $\delta+Kh>1$, which is impossible.

For $n>n_K$, partition $[3]^n$ into coordinate subspaces of
dimension $n_K$ by fixing the other coordinates.  Their relative
$A$-densities average to $\mu_{3,n}(A)$, so some trace has density
at least $\delta$.  The case just proved gives a line in that
trace, hence in $A$.
\end{proof}


\section{Geometry, forms, and elementary estimates}\label{sec:prelim}

\begin{definition}[Triangle coordinates]\label{def:geometry}
\lean{Auto.J,Auto.M,Auto.Lv,Auto.Ang,Auto.u,Auto.v,Auto.σ,Auto.μinc,Auto.proj,Auto.FCMeasure.sigma}
Fix real numbers $\alpha,\beta$ with $\beta\ne0$ and set
\[
 J=\begin{pmatrix}0&-1\\1&0\end{pmatrix},\qquad
 M=\alpha I+\beta J,\qquad
 L_0=0,\quad L_1=I,\quad L_2=M.
\]
Let $\mathbb T=\R/(2\pi\Z)$,
$u_\theta=(\cos\theta,\sin\theta)$, and $v_\theta=Ju_\theta$.
We write $\dd\sigma(\theta)=\dd\theta/(2\pi)$ and also use
$\sigma$ for its pushforward to the unit circle.
The configuration space is $\R^2\times\mathbb T$, with
$p=(x,\theta)$, $\dd\mu(p)=\dd x\dd\sigma(\theta)$, and
\[
 \pi_jp=x+L_ju_\theta\qquad(j=0,1,2).
\]
All measurable functions are represented by Borel versions when
pointwise compositions are needed; equalities of functions in
$L^p$ are understood almost everywhere.
\end{definition}

\begin{lemma}[Similarities and normalization]\label{lem:geometry}
\lean{Auto.lem_geometry_differences,Auto.lem_geometry_commute,Auto.lem_geometry_normalization,Auto.lem_geometry_measurePreserving}
\leanok
Each $L_i-L_j$, $i\ne j$, is an invertible similarity.  The matrix
$M$ commutes with every rotation.  Every ordered noncollinear
triangle is a translate of a rotation of a positive multiple of
$\{0,e_1,Me_1\}$ for a matrix $M$ of the form in
Definition~\ref{def:geometry}.  For each $j$ the map
$(x,\theta)\mapsto(\pi_j(x,\theta),\theta)$ preserves $\mu$.
\end{lemma}
\begin{proof}
For real $a,b$,
$(aI+bJ)^T(aI+bJ)=(a^2+b^2)I$.  Apply this to $M$ and $M-I$,
and also to $I$ and the negatives of these matrices.  The condition
$\beta\ne0$ ensures invertibility of both $M$ and $M-I$.
Every planar rotation is $\cos t\,I+\sin t\,J$, which commutes
with $M$.  Given vertices $q_0,q_1,q_2$, put
$\ell=|q_1-q_0|>0$ and choose a rotation $U$ with
$Ue_1=(q_1-q_0)/\ell$.  Write
$\ell^{-1}U^{-1}(q_2-q_0)=(\alpha,\beta)$.  Noncollinearity is
exactly $\beta\ne0$, and
$\{q_0,q_1,q_2\}=q_0+\ell U\{0,e_1,Me_1\}$.
The last assertion follows from translation invariance of
Lebesgue measure in $x$, separately for each $\theta$.
\end{proof}

\begin{definition}[Counting form and its endpoint norms]\label{def:form}
\lean{Auto.FCMeasure,Auto.form,Auto.endNorm,Auto.formNorm,Auto.FCMeasure.tv,Auto.FCMeasure.dil}
For a finite complex Borel measure $\nu$ on $\R^2$ set
\[
 \Lambda_\nu(f_0,f_1,f_2)
  =\int_{\R^2}\int_{\R^2}
    \prod_{j=0}^2 f_j(x+L_ju)\dd x\dd\nu(u)
\]
whenever the integral is absolutely convergent against $|\nu|$.
For $i\in\{0,1,2\}$ let
\[
 \|\Lambda_\nu\|_i
   =\sup_{\|f_i\|_\infty\le1,\ \|f_j\|_2\le1\ (j\ne i)}
          |\Lambda_\nu(f_0,f_1,f_2)|,
 \qquad \|\Lambda_\nu\|=\|\Lambda_\nu\|_2.
\]
The total variation norm of $\nu$ is $\|\nu\|_{\TV}=|\nu|(\R^2)$.
For $s>0$, $\nu_s$ is the pushforward under $u\mapsto su$.
An $L^1$ function is identified with the measure having that density.
\end{definition}

\begin{lemma}[Basic endpoint bounds and scaling]\label{lem:form-basic}
\lean{Auto.lem_form_basic_integrable,Auto.lem_form_basic_bound,Auto.lem_form_basic_scaling,Auto.lem_form_basic_ae,Auto.lem_form_basic_translate}
\leanok
For each endpoint permutation the defining integral is absolutely
convergent and
\[
 |\Lambda_\nu(f_0,f_1,f_2)|
 \le\|\nu\|_{\TV}\|f_i\|_\infty\prod_{j\ne i}\|f_j\|_2,
 \qquad
 \|\Lambda_{\nu_s}\|_i=\|\Lambda_\nu\|_i.
\]
The form does not depend on the choice of almost-everywhere versions
of its inputs.  Simultaneously translating all inputs leaves the
form unchanged.
\end{lemma}
\begin{proof}
For fixed $u$, bound the $i$th factor by its essential supremum
and apply Cauchy--Schwarz to the other two translated factors in
$x$.  Their $L^2$ norms are unchanged by translation.  Integrate
this bound against $|\nu|$.  The same argument, or Fubini applied
to a null set in any one translated variable, proves independence
of versions.

For the scaling assertion put $p_i=\infty$ and $p_j=2$ for
$j\ne i$, with $1/\infty=0$, and define
$U_s^{(p)}f(x)=s^{-2/p}f(x/s)$.  These are isometries on the
indicated spaces.  The substitutions $x=sy$ and $u=sv$ give
\[
 \Lambda_{\nu_s}(U_s^{(p_0)}f_0,U_s^{(p_1)}f_1,
                   U_s^{(p_2)}f_2)=\Lambda_\nu(f_0,f_1,f_2),
\]
because $\sum_j1/p_j=1$.  Take the supremum and then apply the
same argument with $s^{-1}$.  A common translation is removed
by a translation of $x$.
\end{proof}

\begin{lemma}[Convolution as translations of the inputs]\label{lem:convolution-translate}
\lean{Auto.lem_convolution_translate,Auto.lem_convolution_translate_norm,Auto.lem_convolution_translate_general}
\leanok
For finite measures $\nu,\kappa$ and endpoint inputs,
\begin{equation}\label{eq:convolution-translate}
 \Lambda_{\nu*\kappa}(f_0,f_1,f_2)
 =\int\Lambda_\nu\big(f_0,f_1(\cdot+w),f_2(\cdot+Mw)\big)
       \dd\kappa(w).
\end{equation}
Consequently
$\|\Lambda_{\nu*\kappa}\|_i
 \le\|\kappa\|_{\TV}\|\Lambda_\nu\|_i$.
More generally, any translation-invariant norm bound for the
inputs of $\Lambda_\nu$ is inherited by $\Lambda_{\nu*\kappa}$,
with the extra factor $\|\kappa\|_{\TV}$.
\end{lemma}
\begin{proof}
In the convolution integral replace the displacement $u$ by
$v+w$.  The three vertices become $x$, $x+v+w$, and
$x+Mv+Mw$.  Fubini is justified by the absolute bound of
Lemma~\ref{lem:form-basic}, first against $|\nu|$ and then
against $|\kappa|$.  This proves the identity.  Translations
preserve each input norm, which proves the stated consequences
whenever the corresponding absolute integral or continuous
extension is defined.
\end{proof}

\begin{definition}[Exponential and frequency-scale notation]\label{def:fourier}
\lean{Auto.e,Auto.scaleR}
Write $\e(t)=\exp(2\pi it)$, keeping the Fourier convention of
Definition~\ref{ext:fourier-convention}.  For a function $h$ and
$R>0$, put $h_R(u)=R^2h(Ru)$.
\end{definition}

\begin{lemma}[Fourier representation]\label{lem:fourier-form}
\lean{Auto.lem_fourier_form,Auto.lem_fourier_form_multiplier}
\leanok
For Schwartz inputs and a finite measure $\nu$,
\begin{align*}
 \Lambda_\nu(f_0,f_1,f_2)
  =\iint &\widehat f_0(-\xi_1-\xi_2)
        \widehat f_1(\xi_1)\widehat f_2(\xi_2)\\
   &\hspace{12mm}\cdot\widehat\nu(-\xi_1-M^T\xi_2)
             \dd\xi_1\dd\xi_2.
\end{align*}
In particular, if $\psi$ is smooth, compactly supported, and even,
the form with kernel $\sigma*(\check\psi)_R$ has multiplier
$\widehat\sigma(b)\psi(b/R)$, where $b=\xi_1+M^T\xi_2$.
\end{lemma}
\begin{proof}
Apply Fourier inversion from Lemma~\ref{ext:fourier} to the last two inputs and integrate the
first input in $x$.  The resulting phase in $u$ is $\e(u\cdot b)$,
so its integral is $\widehat\nu(-b)$.  Absolute integrability
follows, for example, from
$\|f_0\|_1
  \|\widehat f_1\|_1\|\widehat f_2\|_1\|\nu\|_{\TV}<\infty$.
For the last assertion use symmetry of $\sigma$, evenness of
$\psi$, and the convolution identity obtained directly by Fubini.
The substitution $v=Ru$ in the Fourier integral gives
$\widehat{h_R}(\xi)=\widehat h(\xi/R)$, supplying the stated scale.
\end{proof}

\begin{lemma}[Density of the circle difference measure]\label{lem:circle-difference}
\lean{Auto.lem_circle_difference}
\leanok
If $N$ is an invertible planar similarity, the pushforward of
$\sigma\otimes\sigma$ under
$(\theta,\varphi)\mapsto N(u_\theta-u_\varphi)$ has a
nonnegative density $w_N\in L^1(\R^2)$ with integral one.
\end{lemma}
\begin{proof}
The absolute Jacobian of this map is
$|\det N|\,|\det(v_\theta,v_\varphi)|
 =|\det N|\,|\sin(\varphi-\theta)|$.
It is nonzero except on the two one-dimensional sets
$\varphi=\theta$ and $\varphi=\theta+\pi$ modulo $2\pi$,
which have zero product measure.  The complement has a countable
cover by open sets on which the map is a diffeomorphism onto its
image, by the inverse function theorem and second countability.
Partition this complement into measurable pieces subordinate to
that cover.  On each piece the change-of-variables theorem shows
that the pushforward is absolutely continuous.  Their countable
sum is absolutely continuous as well, and the omitted set has
zero measure.  The Radon--Nikodym density is nonnegative and
has integral equal to the original mass, namely one.
\end{proof}

\begin{lemma}[Weighted circle compactness and small-set control]\label{lem:circle-compact}
\lean{Auto.circleOp,Auto.lem_circle_compact_small,Auto.lem_circle_compact}
\leanok
Let $N$ be an invertible planar similarity and $W$ a measurable
function with $|W|\le1$.  Define
\[
 Tf(z)=\int W(z,\theta)f(z+Nu_\theta)\dd\sigma(\theta).
\]
There is a nondecreasing function $\omega_N:[0,\infty)\to[0,1]$,
with $\omega_N(s)\to0$ as $s\downarrow0$, independent of $W$,
such that
\[
 \|\sqrt e\,T\|_{L^2\to L^2}\le\omega_N(\|e\|_1)
 \quad\text{if }0\le e\le1,\ e\in L^1.
\]
If $W(z,\theta)=0$ for $z$ outside a bounded set, then $T$ is
compact on $L^2(\R^2)$.
\end{lemma}
\begin{proof}
Cauchy--Schwarz in the probability measure $\sigma$ and translation
invariance give $\|T\|\le1$.
Its adjoint is
\[
 T^*g(y)=\int\overline{W(y-Nu_\varphi,\varphi)}
                         g(y-Nu_\varphi)\dd\sigma(\varphi).
\]
Consequently
\begin{align*}
 TT^*g(z)=\iint &W(z,\theta)
  \overline{W(z+Nu_\theta-Nu_\varphi,\varphi)}\\
 &\hspace{7mm}\cdot g(z+Nu_\theta-Nu_\varphi)
                   \dd\sigma(\theta)\dd\sigma(\varphi).
\end{align*}
By Lemma~\ref{lem:circle-difference} this is an integral operator
with a measurable kernel $K(z,y)$ satisfying
$|K(z,y)|\le w_N(y-z)$.  One precise way to obtain the kernel is
to push forward the weighted measure in $(z,\theta,\varphi)$
under $(z,\theta,\varphi)\mapsto(z,z+Nu_\theta-Nu_\varphi)$;
its total variation is bounded by
$w_N(y-z)\dd z\dd y$, so Radon--Nikodym applies on bounded
$z$ sets and then consistently on all of $\R^2$.
The resulting identity first holds in pairings with bounded compactly
supported tests.  The domination by $w_N$ and
Lemma~\ref{ext:young} extend it to $L^2$ by density.

For $H>0$ split the kernel according to $w_N(y-z)\le H$ and
$w_N(y-z)>H$.  In the first part the kernel of
$\sqrt e\,TT^*\sqrt e$ is bounded in modulus by
$H\sqrt{e(z)e(y)}$, a positive rank-one kernel of operator
norm $H\|e\|_1$.  In the second part, Young's inequality and
$\sqrt e\le1$ bound the operator norm by
$\|w_N\one_{\{w_N>H\}}\|_1$.  Thus
\begin{equation}\label{eq:circle-modulus}
 \|\sqrt e\,T\|^2
  \le H\|e\|_1+\|w_N\one_{\{w_N>H\}}\|_1.
\end{equation}
Take the minimum with one and then the infimum over $H>0$ of
the square root of the right side.  The resulting modulus is
nondecreasing.  It tends to zero because the second term can
first be made small by increasing $H$, and the first can then
be made small by decreasing $\|e\|_1$.

If $W$ has bounded $z$ support $X$, both variables of $K$ are
in $X$.  The part with $w_N\le H$ is a bounded kernel on
$X\times X$, hence is compact by Lemma~\ref{ext:compact-facts}.
The remaining part has norm tending to zero as $H\to\infty$.
The same lemma makes their operator-norm limit $TT^*$ compact.
Lemma~\ref{ext:ttstar} now gives compactness of $T$.
\end{proof}

\begin{lemma}[The trilinear $L^2$ estimate]\label{lem:triple-l2}
\lean{Auto.lem_triple_l2,Auto.lem_triple_l2_sigma}
\leanok
There is a constant $C_M$ such that
\begin{equation}\label{eq:triple-l2}
 |\Lambda_\sigma(f_0,f_1,f_2)|
     \le C_M\prod_{j=0}^2\|f_j\|_2.
\end{equation}
The integral is absolutely convergent for these inputs.
For every finite measure $\kappa$, the same assertions hold
for $\sigma*\kappa$ with the additional factor
$\|\kappa\|_{\TV}$.
\end{lemma}
\begin{proof}
Consider $F(\theta,\varphi)=Mu_\varphi-u_\theta$.  At
$(\theta_0,\theta_0)$ its derivative has columns
$-v_{\theta_0}$ and $Mv_{\theta_0}$, whose determinant has
absolute value $|\beta|>0$.  There is a sufficiently small
interval $I$ about each $\theta_0$ on which $F$ is one-to-one
on $I\times I$ and its absolute Jacobian is bounded below.
For completeness, choose $I$ so that $DF$ differs from the
invertible matrix $DF(\theta_0,\theta_0)$ by less than half
its least singular value.  Integrating this difference along
segments in $I\times I$ proves injectivity and the lower
Jacobian bound.  Compactness of the circle supplies a finite
partition into intervals with these properties.

For such an interval let
$B_I(f_1,f_2)(x)=\int_I f_1(x+u_\theta)
 f_2(x+Mu_\theta)\dd\sigma(\theta)$.
Initially take bounded compactly supported inputs.  Expand
$\|B_I\|_2^2$, take absolute values, and apply Cauchy--Schwarz
on $(x,\theta,\varphi)$ to the factors
\[
 |f_1(x+u_\theta)f_2(x+Mu_\varphi)|,
 \qquad |f_1(x+u_\varphi)f_2(x+Mu_\theta)|.
\]
The two resulting square integrals are identical by interchanging
$\theta$ and $\varphi$.  It follows that
\begin{align*}
 \|B_I(f_1,f_2)\|_2^2
 &\le\int_{I^2}\int
      |f_1(x+u_\theta)|^2|f_2(x+Mu_\varphi)|^2
                 \dd x\dd\sigma(\theta)\dd\sigma(\varphi)\\
 &=\int |f_1(y)|^2
       \int_{I^2}|f_2(y+F(\theta,\varphi))|^2
                     \dd\sigma(\theta)\dd\sigma(\varphi)\dd y\\
 &\le C_M\|f_1\|_2^2\|f_2\|_2^2.
\end{align*}
The last inequality is the change of variables through $F$.
Sum the $L^2$ norms over the finite partition and pair with
$f_0$ by Cauchy--Schwarz.  Applying the same bound to absolute
values, followed by monotone approximation, proves both absolute
convergence and the estimate for arbitrary $L^2$ inputs.
Finally use the identity and translation argument in
Lemma~\ref{lem:convolution-translate}, first for absolute
values against $|\kappa|$.
\end{proof}

\begin{lemma}[Small support in an $L^2$ input]\label{lem:small-support}
\lean{Auto.lem_small_support}
\leanok
There is a nondecreasing modulus $\omega_M(s)\to0$ as
$s\downarrow0$ with the following property.  Suppose $i,j,k$
are distinct, $f_i\in L^\infty$, $f_j,f_k\in L^2$, and
$f_j$ is supported on a measurable set of area at most $s$.
Then, for every finite measure $\kappa$,
\[
 |\Lambda_{\sigma*\kappa}(f_0,f_1,f_2)|
 \le\|\kappa\|_{\TV}\omega_M(s)
              \|f_i\|_\infty\|f_j\|_2\|f_k\|_2.
\]
The assertion for $\sigma$ is included by taking
$\kappa=\delta_0$.
\end{lemma}
\begin{proof}
For $\sigma$, bound the $i$th input by its supremum and set
$z=x+L_ju_\theta$.  The remaining integral is bounded by
\[
 \|f_i\|_\infty\int |f_j(z)|
   \left(\int |f_k(z+(L_k-L_j)u_\theta)|
                          \dd\sigma(\theta)\right)\dd z.
\]
Since $f_j$ is supported on a set $A$ of measure at most $s$,
Cauchy--Schwarz and Lemma~\ref{lem:circle-compact}, with
$e=\one_A$, bound this by
$\omega_{L_k-L_j}(s)\|f_i\|_\infty\|f_j\|_2\|f_k\|_2$.
Take the maximum over the six possible edge maps, which are
invertible similarities by Lemma~\ref{lem:geometry}.  For a
convolution by $\kappa$, translate the inputs as in
Lemma~\ref{lem:convolution-translate}; the small support is
merely translated and its area is unchanged.  Integrate the
bound against $|\kappa|$.
\end{proof}

\begin{remark}\label{rem:moduli}
The symbol $\omega_M$ may be increased by a constant depending
only on $M$ whenever finitely many of the edge moduli are
combined.  It will always be nondecreasing and tend to zero
at zero.  All constants denoted by $C_M$ may change from one
inequality to the next.  These conventions do not assert any
quantitative rate of decay.
\end{remark}


\section{Finite-dimensional scalar kernels}

\begin{definition}[Degrees and cores]\label{rig:def-core}
\lean{Auto.degree,Auto.IsCoreData,Auto.IsCore}
Let $X_0,X_1,X_2\subset\R^2$ be bounded measurable sets, and let
$E\subset\R^2\times \mathbb T$ be measurable, with
$\pi_j(p)\in X_j$ for almost every $p\in E$ and every $j$.
The degree of $z\in\R^2$ at vertex $j$ is
\[
 d_{j,E}(z)=\int_{\mathbb T}
 \one_E(z-L_ju_\theta,\theta)\dd\sigma(\theta).
\]
The data $(E;X_0,X_1,X_2)$ form a \emph{core of degree $c$} if
$\mu(E)>0$ and $d_{j,E}\ge c>0$ almost everywhere on $X_j$, for every $j$.
Sets and their defining relations may throughout be modified on null sets.
\end{definition}

\begin{lemma}[Integration along a vertex]\label{rig:degree-integration}
\lean{Auto.rig_degree_integration,Auto.rig_degree_integration_bounds,Auto.rig_degree_integration_total,Auto.rig_degree_integration_null}
\leanok
For every nonnegative measurable function $g$ on $\R^2$,
\[
 \int_E g(\pi_jp)\dd\mu(p)
 =\int_{\R^2}g(z)d_{j,E}(z)\dd z.
\]
In particular, $0\le d_{j,E}\le1$,
$\int d_{j,E}=\mu(E)$, and pullback by $\pi_j$ sends planar null sets
to $\mu$-null sets.
\end{lemma}
\begin{proof}
Apply Tonelli's theorem, and at fixed $\theta$ make the translation
$z=x+L_ju_\theta$. Its Jacobian is one. The bound on the degree follows
because $\sigma$ is a probability measure. Taking $g=1$ gives its integral,
and taking $g$ to be the indicator of a null set gives the last assertion.
\end{proof}

\begin{lemma}[Extraction of a positive-degree core]\label{rig:prune}
\lean{Auto.rig_prune}
\leanok
If $\mu(E)>0$, then for every
\[
 0<c<\frac{\mu(E)}{\sum_{j=0}^2|X_j|}
\]
there are measurable $Y_j\subset X_j$ and $F\subset E$ such that
$(F;Y_0,Y_1,Y_2)$ is a core of degree $c$. Moreover,
\[
 \mu(E\setminus F)\le c\sum_{j=0}^2|X_j|.
\]
\end{lemma}
\begin{proof}
Set $E^{(0)}=E$ and $X_j^{(0)}=X_j$. Having defined these objects at
step $m$, delete
\[
 V_j^{(m)}=\{z\in X_j^{(m)}:d_{j,E^{(m)}}(z)<c\},
 \qquad X_j^{(m+1)}=X_j^{(m)}\setminus V_j^{(m)},
\]
and let
\[
 E^{(m+1)}=E^{(m)}\cap
 \bigcap_{j=0}^2\pi_j^{-1}(X_j^{(m+1)}).
\]
By Lemma~\ref{rig:degree-integration} and the union bound,
\[
 \mu(E^{(m)}\setminus E^{(m+1)})
 \le\sum_j\int_{V_j^{(m)}}d_{j,E^{(m)}}(z)\dd z
 \le c\sum_j|V_j^{(m)}|.
\]
For each $j$, the deleted vertex sets are disjoint as $m$ varies.
Consequently the sets
$Y_j=\bigcap_mX_j^{(m)}$ and $F=\bigcap_mE^{(m)}$ satisfy the stated loss
bound. Its right-hand side is smaller than $\mu(E)$, so $\mu(F)>0$.
For almost every $z\in Y_j$, survival at every step gives
$d_{j,E^{(m)}}(z)\ge c$ for every $m$. Continuity of the integral under
decreasing bounded integrands gives
$d_{j,F}(z)=\lim_m d_{j,E^{(m)}}(z)\ge c$.
\end{proof}

\begin{definition}[The scalar incidence operator]\label{rig:def-kernel}
\lean{Auto.HE,Auto.incidenceOp,Auto.incidenceOp_ae,Auto.kernelSpace}
On a core $(E;X_0,X_1,X_2)$ let
\[
 \mathcal H_E=\bigoplus_{j=0}^2L^2(X_j;\R),\qquad
 D_E:\mathcal H_E\longrightarrow L^2(E,\mu),\qquad
 D_E f=\sum_{j=0}^2 f_j\circ\pi_j,
 \qquad K_E=\ker D_E.
\]
The norm on $\mathcal H_E$ is the square root of the sum of the three
squared $L^2$ norms. Each component is extended by zero when regarded as
a function on all of $\R^2$.
\end{definition}

\begin{lemma}[Finite-dimensionality and a distance estimate]\label{rig:kernel-gap}
\lean{Auto.rig_kernel_gap_finiteDimensional,Auto.rig_kernel_gap}
\leanok
For every core $E$, the space $K_E$ is finite dimensional. There is
$C_E<\infty$ such that
\[
 \operatorname{dist}_{\mathcal H_E}(f,K_E)
 \le C_E\norm{D_Ef}_{L^2(E,\mu)}\qquad(f\in\mathcal H_E).
\]
The same constant is valid for every translate of $E$ in the base variable.
\end{lemma}
\begin{proof}
The inequality $|a_0+a_1+a_2|^2\le3\sum_j|a_j|^2$, followed by
Lemma~\ref{rig:degree-integration}, shows that $\norm{D_E}\le\sqrt3$.
That lemma also identifies the blocks of $D_E^*D_E$. The diagonal block
at vertex $j$ is multiplication by $d_{j,E}$. For $i\ne j$ the block is
\[
 T_{ji}f_i(z)=\int_{\mathbb T}
 \one_E(z-L_ju_\theta,\theta)
 f_i(z+(L_i-L_j)u_\theta)\dd\sigma(\theta),\qquad z\in X_j.
\]
Its weight has absolute value at most one and bounded support in $z$.
Each $L_i-L_j$ is a nonzero similarity, so
Lemma~\ref{lem:circle-compact} makes every off-diagonal block compact.

Write $D_E^*D_E=A+K$, where $A$ is the diagonal multiplication operator
and $K$ is compact and self-adjoint. If $E$ has degree $c$, then
$cI\le A\le I$. Thus
\[
 B=A^{-1/2}D_E^*D_EA^{-1/2}=I+A^{-1/2}KA^{-1/2}
\]
is a nonnegative compact perturbation of the identity.
Lemma~\ref{ext:compact-gap} implies that $\ker B$ is finite
dimensional and that some $\gamma>0$ satisfies
\[
 \langle Bg,g\rangle\ge
 \gamma\operatorname{dist}(g,\ker B)^2.
\]
Here $\ker B=A^{1/2}K_E$. Substituting $g=A^{1/2}f$ and using
$\norm{A^{1/2}(f-k)}\ge\sqrt c\norm{f-k}$ for $k\in K_E$ gives
\[
 \norm{D_Ef}_2^2\ge\gamma c\operatorname{dist}(f,K_E)^2.
\]
It also proves finite-dimensionality of $K_E$.
Base translation induces unitary identifications of both Hilbert spaces
and intertwines the incidence operators, so it preserves this estimate.
\end{proof}

\begin{lemma}[Uniform integrability of kernel elements]\label{rig:kernel-ui}
\lean{Auto.rig_kernel_ui}
\leanok
There is a nondecreasing function $\omega:[0,\infty)\to[0,\infty)$,
depending only on $M$, with $\omega(s)\to0$ as $s\downarrow0$, such that
the following holds. On any core $G$ of degree $t>0$, if $f\in K_G$,
$0\le e\le1$, and $e\in L^1(\R^2)$, then
\[
 \norm{\sqrt e\,f_j}_2
 \le C_Mt^{-1}\omega(\norm e_1)\norm f_{\mathcal H_G}
 \qquad(j=0,1,2).
\]
The modulus is independent of the core and of its vertex sets.
\end{lemma}
\begin{proof}
Take $\omega$ to be the maximum of the finitely many moduli supplied
by Lemma~\ref{lem:circle-compact} for $L_i-L_j$, $i\ne j$.
The $j$th block of $D_G^*D_Gf=0$ reads
\[
 d_{j,G}f_j=-\sum_{i\ne j}T_{ji}f_i.
\]
Multiplication by $\sqrt e$, the lower bound $d_{j,G}\ge t$ on the
support of $f_j$, and that lemma give
\[
 t\norm{\sqrt e\,f_j}_2
 \le\sum_{i\ne j}\norm{\sqrt e\,T_{ji}f_i}_2
 \le\omega(\norm e_1)\sum_{i\ne j}\norm{f_i}_2.
\]
The last sum is bounded by $\sqrt2\norm f_{\mathcal H_G}$.
\end{proof}

\begin{lemma}[Stability under removal of a small set]\label{rig:kernel-stability}
\lean{Auto.rig_kernel_stability}
\leanok
Fix a core $(E;X_0,X_1,X_2)$ of degree $c$ and a number $t>0$.
There is $\eta_0>0$ such that, whenever
$(F;Y_0,Y_1,Y_2)$ is a core of degree $t$, with
$F\subset E$, $Y_j\subset X_j$, and $\mu(E\setminus F)<\eta_0$,
the restriction map
\[
 R:K_E\longrightarrow K_F,\qquad Rf=(f_j|_{Y_j})_{j=0}^2,
\]
is an isomorphism. The threshold may be chosen using only
$c,t,C_E$, and the fixed modulus in Lemma~\ref{rig:kernel-ui}; in
particular, the same threshold works after a base translation.
\end{lemma}
\begin{proof}
Set $\eta=\mu(E\setminus F)$. Restriction takes $K_E$ to $K_F$ because
$F\subset E$. Every incidence of a vertex in $X_j\setminus Y_j$ is
absent from $F$. Therefore
\[
 c|X_j\setminus Y_j|
 \le\int_{X_j\setminus Y_j}d_{j,E}(z)\dd z\le\eta.
\]
Applying Lemma~\ref{rig:kernel-ui} on $E$, with
$e=\one_{X_j\setminus Y_j}$, proves
\[
 \sum_j\norm{f_j\one_{X_j\setminus Y_j}}_2^2
 \le C_Mc^{-2}\omega(\eta/c)^2\norm f_{\mathcal H_E}^2
 \qquad(f\in K_E).
\]
For sufficiently small $\eta$, this makes $R$ injective.

For the reverse dimension inequality, extend $f\in K_F$ by zero to
$\widetilde f\in\mathcal H_E$. Define
$e_j=d_{j,E\setminus F}$, so $0\le e_j\le1$ and $\int e_j=\eta$.
Since $D_E\widetilde f=0$ almost everywhere on $F$,
\begin{align*}
 \norm{D_E\widetilde f}_2^2
 &\le3\sum_j\int_{E\setminus F}|\widetilde f_j(\pi_jp)|^2\dd\mu(p)\\
 &=3\sum_j\norm{\sqrt{e_j}\,\widetilde f_j}_2^2
 \le C_Mt^{-2}\omega(\eta)^2\norm f_{\mathcal H_F}^2.
\end{align*}
Lemma~\ref{rig:kernel-gap} now gives
\[
 \operatorname{dist}(\widetilde f,K_E)
 \le C_MC_Et^{-1}\omega(\eta)\norm f_{\mathcal H_F}.
\]
Choose $\eta_0$ so that the coefficient is less than one. The orthogonal
projection of the zero extension onto $K_E$ is then injective on $K_F$:
if its value were zero, the displayed inequality would force $f=0$.
Both kernels are finite dimensional by Lemma~\ref{rig:kernel-gap}.
The two injections imply equality of their dimensions, and the already
injective restriction is consequently surjective. The stated dependence
of the threshold follows from the estimates just proved.
\end{proof}

\section{Analytic pieces and measurable rigidity}

\begin{lemma}[A common core for nearby translates]\label{rig:common-core}
\lean{Auto.rig_common_core}
\leanok
Let $E$ be a core of degree $c$, and define
\[
 E_h=\{(x,\theta):(x+h,\theta)\in E\},
 \qquad \eta(h)=\mu(E\setminus E_h).
\]
For all sufficiently small $h\in\R^2$, there is a core
$(F_h;Y_{0,h},Y_{1,h},Y_{2,h})$ of degree $c/4$, with
\[
 F_h\subset E\cap E_h,\qquad
 Y_{j,h}\subset X_j\cap(X_j-h),
\]
such that
\[
 \mu(E\setminus F_h)=\mu(E_h\setminus F_h)\le4\eta(h),
 \qquad |X_j\setminus Y_{j,h}|\longrightarrow0\quad(h\to0).
\]
\end{lemma}
\begin{proof}
Translation continuity in $L^1(\R^2\times \mathbb T)$ gives
$\mu(E\triangle E_h)\to0$; in particular $\eta(h)\to0$.
Start with $E\cap E_h$ on the original vertex sets $X_j$, and perform
the deletion process from Lemma~\ref{rig:prune} at threshold $t=c/4$.
Write $V_j$ for the union of the deleted vertices and
$V=\sum_j|V_j|$. If $L$ denotes the measure of further incidences removed
after the initial removal of $E\setminus E_h$, the same accounting as
in that lemma gives $L\le tV$.
On the other hand, the degrees in the original core yield
\[
 cV\le\sum_j\int_{V_j}d_{j,E}(z)\dd z
 \le3\mu(E\setminus F_h)=3(\eta(h)+L).
\]
Every removed configuration is counted at most three times in this sum.
Thus $cV\le3\eta(h)+3cV/4$, so
$V\le12\eta(h)/c$ and
$\mu(E\setminus F_h)\le\eta(h)+(c/4)V\le4\eta(h)$.
For small $h$ this is less than $\mu(E)$, and the surviving set is a
positive-measure core. Surviving vertices have positive degree in
$E\cap E_h$, so their sets may be taken inside
$X_j\cap(X_j-h)$ up to null sets. Since $\mu(E_h)=\mu(E)$, the two
displayed configuration losses are equal. The bound on $V$ proves the
convergence of the vertex losses.
\end{proof}

\begin{lemma}[Translation matrices]\label{rig:translation-matrices}
\lean{Auto.rig_translation_matrices}
\leanok
Let $\phi^1,\ldots,\phi^d$ be an orthonormal basis of $K_E$, and write
\[
 \Phi_j(z)=(\phi_j^1(z),\ldots,\phi_j^d(z)),
\]
with the components extended by zero. For sufficiently small $h$ there
are sets $Y_{j,h}$ as in Lemma~\ref{rig:common-core} and matrices
$C(h)\in\operatorname{GL}_d(\R)$ such that
\[
 \Phi_j(z+h)=\Phi_j(z)C(h)
 \quad\text{for almost every }z\in Y_{j,h}.
\]
On a sufficiently small neighborhood of zero these matrices form a
continuous local homomorphism:
\[
 C(0)=I,\qquad C(h+k)=C(h)C(k)
\]
whenever $h,k,h+k$ lie in that neighborhood.
\end{lemma}
\begin{proof}
The case $d=0$ is vacuous, so assume $d\ge1$.
The translated rows $\Phi_j(\,\cdot+h)$ give an orthonormal basis for
$K_{E_h}$. By Lemmas~\ref{rig:common-core} and
\ref{rig:kernel-stability}, restriction identifies both $K_E$ and
$K_{E_h}$ with $K_{F_h}$ for all sufficiently small $h$. The threshold
for $E_h$ is the same as for $E$. Two bases of the same finite-dimensional
space differ by a unique invertible matrix, which defines $C(h)$ and
gives the asserted row identities.

Orthonormality says
\[
 \sum_j\int_{X_j}\Phi_j(z)^{\mathsf T}\Phi_j(z)\dd z=I.
\]
Consequently
$A_h=\sum_j\int_{Y_{j,h}}\Phi_j^{\mathsf T}\Phi_j\to I$, by absolute
continuity of the integrals and the vertex-loss estimate. Multiplying
the row identity on the left by $\Phi_j(z)^{\mathsf T}$ and integrating
gives
\[
 A_hC(h)=\sum_j\int_{Y_{j,h}}
 \Phi_j(z)^{\mathsf T}\Phi_j(z+h)\dd z.
\]
The right-hand side minus $A_h$ tends to zero by Cauchy--Schwarz and
$L^2$ translation continuity of each row. Hence $C(h)\to I$ as $h\to0$.

To verify composition, set
\[
 Z_j=Y_{j,h}\cap(Y_{j,k}-h)\cap Y_{j,h+k}.
\]
For almost every $z\in Z_j$, the three row identities give
\[
 \Phi_j(z)C(h+k)=\Phi_j(z+h+k)
 =\Phi_j(z+h)C(k)=\Phi_j(z)C(h)C(k).
\]
The complements of $Z_j$ in $X_j$ have measure tending to zero as
$h,k\to0$. Indeed, besides the three vertex losses, it is enough to
allow for $|X_j\triangle(X_j-h)|$, which tends to zero by translation
continuity of the indicator of $X_j$.
Thus the Gram matrix $\sum_j\int_{Z_j}\Phi_j^{\mathsf T}\Phi_j$ is
invertible for all sufficiently small $h,k$. Multiplying and integrating
the last row identity yields $C(h+k)=C(h)C(k)$.
In particular $C(0)=I$ and $C(-h)=C(h)^{-1}$ where defined.
Continuity at any nearby $h_0$ follows from continuity at zero and
$C(h_0+h)=C(h_0)C(h)$.
\end{proof}

\begin{proposition}[Analytic pieces for a scalar kernel]\label{rig:pieces}
\lean{Auto.rig_pieces}
\leanok
Every positive-measure core $E$ has a measurable subset $E'\subset E$
of positive measure such that each component of an orthonormal basis of
$K_E$, evaluated at its vertex, agrees on $E'$ with a real analytic
function defined on all of $\R^2$.
\end{proposition}
\begin{proof}
If $K_E=\{0\}$, the conclusion is vacuous and we take $E'=E$.
Otherwise apply Lemma~\ref{ext:matrix-group} to the local homomorphism
of Lemma~\ref{rig:translation-matrices}.  It gives commuting matrices
$T_1,T_2$ and an everywhere real analytic function
\[
 P(h)=\exp(h_1T_1+h_2T_2)
\]
which equals $C(h)$ near zero.  Use the rows $\Phi_j$ from
Lemma~\ref{rig:translation-matrices}, taking finite Borel
representatives for them, and write $z_j=\pi_jp$.
The translation identity implies
\[
 \Phi_j(z_j)=\Phi_j(z_j-h_j)P(h_j)
\]
whenever $z_j\in Y_{j,h_j}+h_j$, apart from a planar null set.
Lemma~\ref{rig:degree-integration} shows that the exceptional
configurations in $E$ have measure at most
\[
 |X_j\setminus(Y_{j,h_j}+h_j)|
 \le |X_j\triangle(X_j+h_j)|+|X_j\setminus Y_{j,h_j}|,
\]
which tends to zero as $h_j\to0$.
Choose a ball $B$ centered at zero so small that, for every triple
$(h_0,h_1,h_2)\in B^3$, all three row identities hold simultaneously
on a subset of $E$ of measure at least $\mu(E)/2$.

The set of $(p,h_0,h_1,h_2)$ satisfying these identities is jointly
measurable: this uses the Borel representatives and the explicit
continuous function $P$, not any measurable choice of the cores.
Its measure in $E\times B^3$ is at least $\mu(E)|B|^3/2>0$.
For each fixed $p$ make the change of variables
$b_j=z_j-h_j$, $j=0,1,2$, in the six translation variables.
The Jacobian has absolute value one. Fubini's theorem gives a fixed
triple $(b_0,b_1,b_2)$ for which the set of configurations satisfying
\[
 \Phi_j(z_j)=\Phi_j(b_j)P(z_j-b_j),\qquad j=0,1,2,
\]
has positive $\mu$-measure. Take this set to be $E'$.
For each $j$, the row-valued function
$z\mapsto\Phi_j(b_j)P(z-b_j)$ is real analytic on all of $\R^2$.
Its components supply the required analytic representatives.
\end{proof}

\begin{lemma}[Rigidity for analytic fields]\label{rig:analytic-rigidity}
\lean{Auto.rig_analytic_rigidity}
\leanok
Let $a_j:\R^2\to\R^2$ be real analytic and suppose that, for every
$(x,\theta)\in\R^2\times \mathbb T$,
\[
 \sum_{j=0}^2a_j(x+L_ju_\theta)=0,
 \qquad
 v_\theta\cdot
 \bigl(a_1(x+u_\theta)+M^{\mathsf T}a_2(x+Mu_\theta)\bigr)=0.
\]
Then
\[
 a_1(x+u_\theta)+M^{\mathsf T}a_2(x+Mu_\theta)=0
 \qquad\text{for every }(x,\theta).
\]
\end{lemma}
\begin{proof}
Fix the notation
\[
 A=a_0(x),\quad B=a_1(x+u_\theta),\quad
 C=a_2(x+Mu_\theta),\quad u=u_\theta,\quad v=Ju.
\]
Since $u,v$ are an orthonormal basis, the second hypothesis says that
$B+M^{\mathsf T}C=tu$ for a smooth real scalar function $t=t(x,\theta)$.
Combining this with $A+B+C=0$ gives
\begin{equation}\label{rig:algebraic-identities}
 (M^{\mathsf T}-I)C=A+tu,
 \qquad -(M^{\mathsf T}-I)B=M^{\mathsf T}A+tu.
\end{equation}
Put $K=Da_0(x)$, with column-vector convention for the derivative,
and define
\[
 X_1=\partial_\theta-v\cdot\nabla_x,
 \qquad X_2=\partial_\theta-Mv\cdot\nabla_x.
\]
The chain rule gives $X_1B=0$ and $X_2C=0$.
Applying $X_1$ to the second identity in
\eqref{rig:algebraic-identities}, and $X_2$ to the first, yields
\begin{equation}\label{rig:transport-identities}
 (X_1t)u+tv=M^{\mathsf T}Kv,
 \qquad (X_2t)u+tv=KMv.
\end{equation}
Take the inner product of the second identity with $v$. Since
$u\cdot v=0$ and $|v|=1$, it gives
\[
 t(x,\theta)=v^{\mathsf T}KMv.
\]
In particular $t(x,\theta+\pi)=t(x,\theta)$.
In the expanded versions of \eqref{rig:transport-identities}, the
two spatial transport terms
$(v\cdot\nabla_xt)u$ and $(Mv\cdot\nabla_xt)u$ are unchanged by
$\theta\mapsto\theta+\pi$. The terms
$\partial_\theta(tu)$, $M^{\mathsf T}Kv$, and $KMv$ change sign.
Subtracting each equation at $\theta+\pi$ from the equation at $\theta$
therefore gives
\[
 \partial_\theta(tu)=M^{\mathsf T}Kv=KMv.
\]
As $v$ ranges over the unit circle, this proves
$M^{\mathsf T}K=KM$. Integration in $\theta$ gives
$tu=KMu+c(x)$. Both $tu$ and $KMu$ change sign under
$\theta\mapsto\theta+\pi$, so $c(x)=0$.

Thus $KM$ sends every line through the origin into itself. A real
linear map on $\R^2$ with this property is scalar: applying it to
$e_1,e_2,e_1+e_2$ shows that its two diagonal eigenvalues agree.
Hence $KM=\lambda(x)I$. Since $M$ is invertible, $K$ commutes with $M$.
The identity $M^{\mathsf T}K=KM$ consequently becomes
\[
 (M^{\mathsf T}-M)K=0.
\]
But $M^{\mathsf T}-M=-2\beta J$ is invertible because $\beta\ne0$.
It follows that $K=0$, and then $t=0$. Therefore
$B+M^{\mathsf T}C=tu=0$, as required.
\end{proof}

\begin{theorem}[Measurable stationarity forces zero displacement frequency]
\label{rig:exact}
\lean{Auto.rig_exact}
\leanok
Let $(E;X_0,X_1,X_2)$ be a core of positive degree, and let
$a_j\in L^2(X_j;\R^2)$ satisfy, for almost every $p=(x,\theta)\in E$,
\[
 \sum_{j=0}^2a_j(\pi_jp)=0,
 \qquad
 v_\theta\cdot
 \bigl(a_1(\pi_1p)+M^{\mathsf T}a_2(\pi_2p)\bigr)=0.
\]
Then
\[
 a_1(\pi_1p)+M^{\mathsf T}a_2(\pi_2p)=0
 \qquad\text{for almost every }p\in E.
\]
\end{theorem}
\begin{proof}
Suppose instead that the asserted vector is nonzero on a measurable
set $G\subset E$ of positive measure. By Lemma~\ref{rig:prune}, $G$
contains a positive-measure core $F$ on bounded vertex subsets.
Restrict the fields to these subsets. Each of the two scalar-coordinate
triples of the fields belongs to $K_F$ by the first stationary equation.
Express these triples in an orthonormal basis of $K_F$ and apply
Proposition~\ref{rig:pieces}. On a subset $F'\subset F$ of positive
measure, the fields at their vertices agree with real analytic fields
$\widetilde a_j:\R^2\to\R^2$.

The two stationary expressions for these analytic fields are real
analytic on the connected real analytic manifold $\R^2\times \mathbb T$.
They vanish on $F'$, except on a null set.
Lemma~\ref{ext:analytic-zero}, applied to each scalar component on
$\R^2\times\mathbb T$, shows that both stationary equations hold
everywhere for the fields $\widetilde a_j$.
Lemma~\ref{rig:analytic-rigidity} makes their displacement-frequency
vector identically zero. On $F'$ it equals the displacement-frequency
vector of the original fields, which was nonzero almost everywhere
because $F'\subset G$. This contradiction proves the theorem.
\end{proof}


\section{Translation correlations and local cancellation}

\begin{definition}[Admissible sequences and translates]\label{ann:sequence}
\lean{Auto.IsAdmissibleSeq,Auto.annT,Auto.annWeakStar}
Fix a bounded box $Q\subset\R^2$ and a constant $C<\infty$.
An admissible sequence consists of $R_n\to\infty$ and complex-valued
functions $f_{j,n}$, $j\in\{0,1,2\}$, such that
\[
 |f_{j,n}|\le\one_Q,\qquad \operatorname{Lip}(f_{j,n})\le CR_n.
\]
For $s\in\R^2$ write
\[
 T^s_{j,n}(z)=f_{j,n}(z+s/R_n).
\]
Weak-* convergence on a vertex space or on configuration space means
convergence against every $L^1$ test function on that space.
\end{definition}

\begin{lemma}[Simultaneous correlation limits]\label{ann:extract}
\lean{Auto.ann_extract}
\leanok
Every admissible sequence has a subsequence with the following properties.
For all real shifts $s,t,s_0,s_1,s_2\in\R^2$, the limits
\begin{align*}
 T^s_{j,n}\overline{T^t_{j,n}}&\stackrel{*}{\rightharpoonup}
 k_j(z;s,t),\\
 \prod_{j=0}^2T^{s_j}_{j,n}(\pi_jp)
 &\stackrel{*}{\rightharpoonup}h(p;s_0,s_1,s_2)
\end{align*}
exist. The limits have representatives that are continuous in all shift
variables off a fixed null set in the respective underlying space.
They are bounded by one, and their Lipschitz constants in the shifts
are bounded in terms of $C$. The three-point convergence also holds
when the shifts $s_j$ are replaced by bounded measurable functions of $p$.
\end{lemma}
\begin{proof}
First take all shifts in $\mathbb Q^2$. There are countably many
two-point and three-point products. Each is bounded by one in
$L^\infty$. Lemma~\ref{ext:weakstar} gives all these limits on one
subsequence by its countable diagonal construction.
For example,
\[
 |T^s_{j,n}(z)-T^t_{j,n}(z)|\le C|s-t|.
\]
The elementary product estimate therefore gives uniform Lipschitz
bounds for every correlation as a function of its shifts. These
inequalities pass to the weak-* limits. Since there are only
countably many rational shifts, they hold simultaneously outside
one null set. Complete the rational shift variables by continuity
at every point outside that set. The same inequalities, tested
against an arbitrary $L^1$ function, show that the original products
converge weak-* at every real choice of shifts.

The resulting limit functions are jointly measurable: they are
measurable in the underlying variable at rational shifts and are
continuous in the shift variables. For bounded measurable shifts,
round each coordinate to a rational grid of mesh $\varepsilon$.
The rounded shifts take finitely many values. Convergence for them
follows by testing on their measurable level sets. Both the original
products and their limits change by at most $O_C(\varepsilon)$ when
the shifts are rounded. Letting $\varepsilon\downarrow0$ proves the
last assertion.
\end{proof}

\begin{lemma}[Measurable spectral measures]\label{ann:spectral}
\lean{Auto.ann_spectral}
\leanok
For a subsequence as in Lemma~\ref{ann:extract}, there are measurable
kernels of finite positive Borel measures $z\mapsto\nu_{j,z}$ on
$\R^2$, of mass at most one, such that
\[
 k_j(z;s,t)=\int_{\R^2}\e((s-t)\cdot\xi)\,\dd\nu_{j,z}(\xi)
\]
for almost every $z$ and all $s,t\in\R^2$.
\end{lemma}
\begin{proof}
For fixed $s,t$, changing variables by $z\mapsto z+t/R_n$ in a
compactly supported test function shows that
\[
 k_j(z;s,t)=k_j(z;s-t,0)
\]
almost everywhere. Translation continuity of the test in $L^1$
justifies this assertion. Take rational shifts first and then use
Lemma~\ref{ann:extract} to obtain one exceptional null set for all
real shifts. Write $c_{j,z}(s)=k_j(z;s,0)$.

For any finite family of rational shifts $s_a$ and rational complex
coefficients $c_a$, the nonnegative functions
\[
 \left|\sum_a c_aT^{s_a}_{j,n}(z)\right|^2
\]
have weak-* limit
$\sum_{a,b}c_a\overline{c_b}c_{j,z}(s_a-s_b)$. It is nonnegative
almost everywhere. Countability and continuity extend this conclusion
to all shifts and coefficients on one common full-measure set.
Thus $c_{j,z}$ is a Lipschitz positive-definite function there,
by the shift bound in Lemma~\ref{ann:extract}.
Theorem~\ref{ext:bochner} supplies a finite positive measure $\nu_{j,z}$
with the asserted Fourier representation. Its mass is
$c_{j,z}(0)$, which lies in $[0,1]$. Set $\nu_{j,z}=0$ on the
exceptional set.

For each shift $s$, $z\mapsto c_{j,z}(s)$ is measurable.
Set it equal to zero on the same exceptional set.
Lemma~\ref{ext:measurable-spectral} now proves that
$z\mapsto\nu_{j,z}(B)$ is measurable for every Borel set $B$,
which is the kernel assertion.
\end{proof}

\begin{definition}[Polynomial translates]\label{ann:polynomials}
\lean{Auto.charE,Auto.trigPoly,Auto.tensorPoly,Auto.polyTranslate,Auto.tensorTranslate}
Write $e_s(\xi)=\e(s\cdot\xi)$. A trigonometric polynomial on
$\R^2$ is a finite complex linear combination of the functions $e_s$.
For $a=\sum_s c_se_s$ put
\[
 F_{j,n}[a](z)=\sum_s c_s T^s_{j,n}(z).
\]
For a finite sum of elementary tensors
$G=\sum_{d=1}^D b_d\otimes c_d$, with $b_d,c_d$ trigonometric
polynomials, put
\[
 G_n(p)=\sum_{d=1}^D
 F_{1,n}[b_d](\pi_1p)F_{2,n}[c_d](\pi_2p).
\]
When $p$ is fixed, abbreviate $\nu_j=\nu_{j,\pi_jp}$.
\end{definition}

\begin{lemma}[Factorization of two-vertex quadratic correlations]
\label{ann:factor}
\lean{Auto.ann_factor_one,Auto.ann_factor_two}
\leanok
For the notation of Definition~\ref{ann:polynomials},
\begin{align*}
 |F_{0,n}[a](\pi_0p)|^2
 &\stackrel{*}{\rightharpoonup}\int|a|^2\,\dd\nu_0,\\
 |G_n(p)|^2
 &\stackrel{*}{\rightharpoonup}
 \int|G|^2\,\dd(\nu_1\otimes\nu_2).
\end{align*}
\end{lemma}
\begin{proof}
The first assertion follows by expanding the square and using
Lemma~\ref{ann:spectral}. Pullback by $\pi_0(x,\theta)=x$ preserves
the required convergence, by integration of the test in $\theta$.

For the second assertion, first consider a term in the expansion
of $|G_n|^2$. It has the form
\[
 A_n(\pi_1p)B_n(\pi_2p),
\]
where $A_n$ and $B_n$ are products of two polynomial translates.
Their one-vertex weak-* limits $A$ and $B$ are the corresponding
integrals against $\nu_{1,z}$ and $\nu_{2,z}$.
Take a bounded compactly supported measurable test $\varphi(p)$.
In coordinates $z=x+u_\theta$, its pairing with the displayed
product is
\[
 \int A_n(z)\,(T_\varphi B_n)(z)\,\dd z,
 \qquad
 T_\varphi b(z)=\int
 \varphi(z-u_\theta,\theta)b(z+(M-I)u_\theta)\,\dd\sigma(\theta).
\]
The weight of $T_\varphi$ is bounded and has bounded support in $z$.
Since $M-I$ is a nonzero similarity, Lemma~\ref{lem:circle-compact}
implies that $T_\varphi$ is compact on $L^2$.
Only fixed bounded vertex boxes enter this integral. On these
boxes the uniformly bounded sequences $A_n,B_n$ converge weakly
in $L^2$ to $A,B$. Lemma~\ref{ext:compact-facts} gives strong convergence of
$T_\varphi B_n$ to $T_\varphi B$, so the pairing converges to
$\int A\,T_\varphi B$. This is exactly the pairing of
$A(\pi_1p)B(\pi_2p)$ with $\varphi$.

Summing the finitely many expanded terms gives
\[
 \sum_{d,e}
 \left(\int b_d\overline{b_e}\,\dd\nu_1\right)
 \left(\int c_d\overline{c_e}\,\dd\nu_2\right)
 =\int|G|^2\,\dd(\nu_1\otimes\nu_2).
\]
Uniform boundedness of the products extends convergence from
bounded compactly supported tests to all $L^1$ tests.
\end{proof}

\begin{lemma}[The limiting bilinear correlation form]\label{ann:form}
\lean{Auto.ann_form}
\leanok
For almost every configuration $p$, there is a complex bilinear form
\[
 \ell_p:L^2(\nu_0)\times L^2(\nu_1\otimes\nu_2)\longrightarrow\C
\]
such that
\begin{equation}\label{ann:cs}
 |\ell_p(a,G)|\le
 \|a\|_{L^2(\nu_0)}\|G\|_{L^2(\nu_1\otimes\nu_2)}.
\end{equation}
For trigonometric polynomials and their finite tensor sums,
$\ell_p(a,G)$ is the weak-* limit of
$F_{0,n}[a](\pi_0p)G_n(p)$. In particular,
\[
 \ell_p(e_{s_0},e_{s_1}\otimes e_{s_2})
 =h(p;s_0,s_1,s_2).
\]
\end{lemma}
\begin{proof}
Expand $F_{0,n}[a]G_n$ into three-point products and apply
Lemma~\ref{ann:extract} to define its limit on the polynomial
spaces. For every measurable configuration set $B$ of finite measure,
Cauchy--Schwarz gives
\[
 \left|\int_B F_{0,n}[a](\pi_0p)G_n(p)\,\dd\mu(p)\right|
 \le
 \left(\int_B |F_{0,n}[a](\pi_0p)|^2\,\dd\mu\right)^{1/2}
 \left(\int_B |G_n|^2\,\dd\mu\right)^{1/2}.
\]
Pass to the limit using Lemma~\ref{ann:factor}.  Write
$l(p)=\ell_p(a,G)$, $a_0(p)=\|a\|_{L^2(\nu_0)}^2$, and
$b_0(p)=\|G\|_{L^2(\nu_1\otimes\nu_2)}^2$ for the resulting
three measurable functions.  For every complex number $q$,
Cauchy--Schwarz followed by $2uv\le u^2+v^2$ gives
\[
 \int_B\bigl(|q|^2a_0+b_0-2\operatorname{Re}(ql)\bigr)\dd\mu
 \ge0.
\]
For each rational complex $q$ the integrand is therefore nonnegative
almost everywhere: otherwise a bounded subset of its negative set
would contradict this inequality.  Remove the union of the resulting
countably many null sets and then use continuity in $q$.
The inequality holds for all $q$ at each remaining point.
If $a_0>0$, minimize the quadratic expression in $q$ to obtain
$|l|^2\le a_0b_0$; if $a_0=0$, varying $q$ forces $l=0$.
This proves \eqref{ann:cs} without a differentiation theorem.
Initially use rational shifts and rational complex coefficients,
so that one null set suffices.
The continuity of the limits in the shifts, and the Fourier
representation in Lemma~\ref{ann:spectral}, extend the inequality
to all polynomials and all finite tensor sums.

Characters with rational shifts span a dense subspace of
$L^2(\nu_j)$ by Lemma~\ref{ext:fourier-uniqueness}.
Their tensor products are dense in $L^2(\nu_1\otimes\nu_2)$ as well.
Indeed, if $H$ is orthogonal to every $e_s\otimes e_t$ with
$s,t\in\Q^2$, then for each $s$ the function
\[
 \xi_2\longmapsto\int H(\xi_1,\xi_2)
                       \overline{e_s(\xi_1)}\dd\nu_1(\xi_1)
\]
belongs to $L^2(\nu_2)$ by Cauchy--Schwarz and is orthogonal to
every rational character.  It is zero almost everywhere.  Countability
of $s$ and Fubini show that almost every section $H(\cdot,\xi_2)$
is orthogonal to every rational character in $L^2(\nu_1)$; hence
$H=0$ almost everywhere.  Inequality
\eqref{ann:cs} therefore gives the unique bounded bilinear
extension. It also shows that its definition is independent of
all polynomial or tensor representations.
\end{proof}

\begin{lemma}[Infinitesimal translation and rotation invariance]
\label{ann:invariance}
\lean{Auto.ann_invariance}
\leanok
Outside a fixed null set of configurations, for every
$s\in\R^2$, $h\in\R$, $a\in L^2(\nu_0)$, and
$G\in L^2(\nu_1\otimes\nu_2)$,
\begin{equation}\label{ann:inv}
 \ell_p(a,G)=
 \ell_p\bigl(e_sa,
 (e_{s+h v_\theta}\otimes e_{s+hMv_\theta})G\bigr).
\end{equation}
\end{lemma}
\begin{proof}
Begin with polynomial arguments and fixed parameters $s,h$.
Translate the configuration by
\[
 (x,\theta)\longmapsto(x+s/R_n,\theta+h/R_n).
\]
This transformation preserves $\dd x\,\dd\sigma(\theta)$.
For a compactly supported continuous test, the resulting change
in the test tends to zero in $L^1$. Taylor expansion gives
\[
 L_ju_{\theta+h/R_n}-L_ju_\theta
 =hL_jv_\theta/R_n+O_M(h^2/R_n^2).
\]
Multiplication by the Lipschitz bound $CR_n$ shows that the
resulting error in every input tends to zero. Thus the transformed
product has the same weak-* limit as the original product, and
its limiting shifts at the three vertices are respectively
$s$, $s+hv_\theta$, and $s+hMv_\theta$ in addition to the original
polynomial shifts. The variable-shift assertion of
Lemma~\ref{ann:extract} identifies its limit and proves
\eqref{ann:inv} for the polynomial arguments.

Take rational polynomial data and rational $s,h$ first. Their
countability gives one exceptional null set. At each remaining
$p$, characters vary continuously in $L^2$ of each finite spectral
measure, by dominated convergence. Inequality~\eqref{ann:cs}
therefore extends the identity to real $s,h$, to all polynomial
data, and finally by density to both Hilbert spaces. Character
multipliers have norm one throughout this extension.
\end{proof}

\begin{lemma}[Atoms of one-dimensional projections]\label{ann:fibers}
\lean{Auto.annH,Auto.ann_fibers,Auto.ann_fibers_hyperplane}
\leanok
Let $z\mapsto\nu_z$ be a measurable kernel of finite positive
measures on $\R^2$, of mass at most one, and let $N$ be invertible.
For almost every $(z,\theta)$, every positive-mass fiber of
\[
 \xi\longmapsto (Nv_\theta)\cdot\xi
\]
is concentrated at a single point atom of $\nu_z$.
Consequently, for almost every configuration $p$, the restriction of
$\nu_1\otimes\nu_2$ to
\begin{equation}\label{ann:hyperplane}
 H_p=\{(\xi_1,\xi_2):
 v_\theta\cdot\xi_1+Mv_\theta\cdot\xi_2=0\}
\end{equation}
is concentrated on countably many pairs of point atoms.
\end{lemma}
\begin{proof}
For distinct $\xi,\xi'$, the equation
$(Nv_\theta)\cdot(\xi-\xi')=0$ holds at only two angles modulo
$2\pi$. Lemma~\ref{ext:kernel-integration} makes
$\dd z\dd\nu_z(\xi)\dd\nu_z(\xi')$ a sigma-finite measure.
Tonelli with the independent angular measure therefore gives, for almost every
$(z,\theta)$,
\[
 \iint\one_{\{\xi\ne\xi',\,
 (Nv_\theta)\cdot(\xi-\xi')=0\}}
 \,\dd\nu_z(\xi)\,\dd\nu_z(\xi')=0.
\]
For completeness, integrate first over bounded sets of $z$ and
then exhaust $\R^2$. On a positive-mass fiber the product measure
is thus supported on the diagonal. A finite positive measure
with this property is a point mass: otherwise some Borel subset
has measure strictly between zero and the total mass, and its
product with its complement gives positive off-diagonal mass.

Apply this conclusion to $N=I$ at vertex $1$ and to $N=M$ at
vertex $2$. The map $(x,\theta)\mapsto(\pi_jp,\theta)$ preserves
measure, so the conclusions hold for almost every $p$ simultaneously.
Let $\eta_1$ and $\eta_2$ be the pushforwards of $\nu_1$ and
$\nu_2$ by $v_\theta\cdot\xi$ and $Mv_\theta\cdot\xi$.
The product mass of $H_p$ is
\[
 \sum_{t:\,\eta_1(\{t\})>0}
 \eta_1(\{t\})\eta_2(\{-t\}).
\]
Only countably many terms can be positive: a finite measure has
finitely many atoms of mass at least $1/n$ for each $n\ge1$.
By the first part,
each positive term is carried by one pair of point atoms.
This proves the assertion about the restricted measure, not just
its total mass.
\end{proof}

\begin{lemma}[Countably many measurable selectors cover all atoms]
\label{ann:selectors}
\lean{Auto.ann_selectors}
\leanok
For a measurable kernel $z\mapsto\nu_z$ of finite positive Borel
measures on $\R^2$, there is a countable family of measurable maps
$a_m:\R^2\to\R^2$ such that every point atom of every nonzero
$\nu_z$ occurs among the values $a_m(z)$.
\end{lemma}
\begin{proof}
Where the mass is nonzero, normalize $\nu_z$ to a probability
measure. For each of its two coordinate marginals and every
$r\in\mathbb Q\cap(0,1)$, form the quantile
\[
 q_{k,r}(z)=\inf\{t\in\R:
 \nu_z(\{\xi:\xi_k\le t\})/\nu_z(\R^2)\ge r\},
 \qquad k=1,2.
\]
These values are finite, since the normalized distribution functions
tend to zero and one at the two ends of the real line.
They are measurable: on the set $m(z)=\nu_z(\R^2)>0$,
\[
 \{z:q_{k,r}(z)<a\}
 =\bigcup_{\substack{t\in\Q\\t<a}}
   \{z:\nu_z(\{\xi:\xi_k\le t\})\ge r\,m(z)\}.
\]
Indeed, the threshold set defining the infimum is an upper interval,
so it meets $(-\infty,a)$ precisely when it contains a rational point
less than $a$.  Every set in the displayed union is measurable.
Take every pair $(q_{1,r},q_{2,s})$ with rational $r,s\in(0,1)$,
and set these maps to zero where $\nu_z=0$.

If $(a,b)$ is a point atom, then $a$ and $b$ are atoms of the
respective marginals. At either coordinate the distribution
function has a positive jump. Every quantile level strictly
inside that jump returns the coordinate in question. Choose a
rational level inside each jump. The resulting pair is $(a,b)$.
It is harmless that other pairs of quantiles need not be atoms.
\end{proof}

\begin{lemma}[Rigidity for all contributing atomic triples]
\label{ann:atomic-rigidity}
\lean{Auto.ann_atomic_rigidity}
\leanok
For almost every configuration $p$, every triple of point atoms
$\xi_j$ of $\nu_{j,\pi_jp}$ satisfying
\begin{equation}\label{ann:stationarity}
 \xi_0+\xi_1+\xi_2=0,\qquad
 v_\theta\cdot(\xi_1+M^T\xi_2)=0
\end{equation}
also satisfies $\xi_1+M^T\xi_2=0$.
\end{lemma}
\begin{proof}
Apply Lemma~\ref{ann:selectors} to the three spectral kernels.
There are countably many triples of the resulting measurable
vector fields. Fix one such triple $a_0,a_1,a_2$ and let $E$ be
the measurable set where its values at the vertices satisfy
\eqref{ann:stationarity} but
$a_1(\pi_1p)+M^Ta_2(\pi_2p)\ne0$.
If $\mu(E)>0$, restrict to a positive-measure subset on which
$x$ and all three field values are bounded. This is possible by
a countable exhaustion. If the respective bounds are $L$ and $K$,
take $X_j=\{z:|z|\le L+\|L_j\|,\ |a_j(z)|\le K\}$.
All vertices of this subset then lie in the bounded measurable sets
$X_j$, and the fields are bounded on them. Lemma~\ref{rig:prune}
provides a positive-measure core in this subset. The restrictions
of the fields to its vertex sets are square integrable, and they
satisfy the stationary equations there. Theorem~\ref{rig:exact}
forces $a_1(\pi_1p)+M^Ta_2(\pi_2p)=0$ almost everywhere on that
core, a contradiction. Hence $\mu(E)=0$.

Remove the union of these exceptional sets over the countable
family of triples. At every remaining configuration, each triple
of point atoms is represented by a triple of the selectors, so
the asserted implication holds for all atomic triples at once.
\end{proof}

\begin{lemma}[Invariance under displacement smoothing]
\label{ann:displacement}
\lean{Auto.ann_displacement}
\leanok
For almost every configuration $p$ and every $y\in\R^2$,
\[
 \ell_p(1,e_y\otimes e_{My})=\ell_p(1,1).
\]
\end{lemma}
\begin{proof}
Fix $p$ outside all the exceptional sets in the preceding lemmas.
In \eqref{ann:inv}, put $s=0$ and average $h$ over $[0,H]$.
The resulting multiplier on the second Hilbert space is
\[
 \frac1H\int_0^H
 \e\bigl(h(v_\theta\cdot\xi_1+Mv_\theta\cdot\xi_2)\bigr)\,\dd h.
\]
It is bounded by one and converges pointwise to $\one_{H_p}$.
Dominated convergence gives strong convergence after multiplication
by any $G\in L^2(\nu_1\otimes\nu_2)$.
Boundedness of $\ell_p$ therefore gives
\begin{equation}\label{ann:restrict-H}
 \ell_p(a,G)=\ell_p(a,\one_{H_p}G).
\end{equation}

By Lemma~\ref{ann:fibers}, the measure on $H_p$ is carried by
countably many pairs of point atoms. For one such pair
$(\xi_1,\xi_2)$, let $E_{\xi_1,\xi_2}$ be the indicator of that
point in the product frequency space. Put $h=0$ in
\eqref{ann:inv} and average $s$ over $[0,H]^2$.
The scalar character contributed by this second argument can be
moved into the first argument by bilinearity. Strong dominated
convergence in $L^2(\nu_0)$ then gives
\[
 \ell_p(a,E_{\xi_1,\xi_2})
 =\ell_p\bigl(a\one_{\{-\xi_1-\xi_2\}},E_{\xi_1,\xi_2}\bigr).
\]
This is zero unless $\xi_0=-\xi_1-\xi_2$ is a point atom of
$\nu_0$. In the remaining case the three atoms satisfy
\eqref{ann:stationarity}, so Lemma~\ref{ann:atomic-rigidity}
gives $\xi_1+M^T\xi_2=0$. Consequently
\[
 e_y(\xi_1)e_{My}(\xi_2)
 =\e\bigl(y\cdot(\xi_1+M^T\xi_2)\bigr)=1.
\]
Approximate $\one_{H_p}G$ by its restrictions to finitely many
of the atomic pairs. The approximation converges in the product
$L^2$ space, and \eqref{ann:cs} passes the preceding identity to
the limit. Thus, in fact,
\[
 \ell_p(a,(e_y\otimes e_{My})G)=\ell_p(a,G)
\]
for every $a,G,y$. Taking $a=1$ and $G=1$ proves the assertion.
\end{proof}

\begin{proposition}[Cancellation for bounded Lipschitz inputs]
\label{ann:local}
\lean{Auto.ann_local}
\leanok
Let $\Psi\in L^1(\R^2)$ satisfy $\int\Psi=0$, and put
$\Psi_R(u)=R^2\Psi(Ru)$. For every bounded box $Q$ and every
$C<\infty$,
\[
 \sup_{\substack{|f_j|\le\one_Q\\
                  \operatorname{Lip}(f_j)\le CR}}
 |\Lambda_{\sigma*\Psi_R}(f_0,f_1,f_2)|\longrightarrow0
 \qquad (R\to\infty).
\]
\end{proposition}
\begin{proof}
It suffices to consider an arbitrary sequence of admissible
inputs with $R_n\to\infty$ and find a subsequence on which the
forms tend to zero. Otherwise a sequence on which their absolute
values are bounded below would contradict that property.
Take the subsequence and correlation forms constructed above.
Choose $w\in C_c(\R^2)$ with $w=1$ on $Q$. The convolution
identity in Lemma~\ref{lem:convolution-translate} gives
\begin{align*}
 &\Lambda_{\sigma*\Psi_{R_n}}(f_{0,n},f_{1,n},f_{2,n})\\
 &\quad=\int\Psi(y)\int
 w(x)f_{0,n}(x)
 f_{1,n}(x+u_\theta+y/R_n)
 f_{2,n}(x+Mu_\theta+My/R_n)\,\dd\mu(p)\,\dd y.
\end{align*}
For each fixed $y$ the inner integral tends, by
Lemma~\ref{ann:form}, to
\[
 \int w(x)\ell_p(1,e_y\otimes e_{My})\,\dd\mu(p).
\]
All these integrals are bounded by $\|w\|_1$, independently of
$n,y$. Dominated convergence with $\Psi\in L^1$ is therefore
applicable. Lemma~\ref{ann:displacement} makes the limit independent
of $y$, and $\int\Psi=0$ makes the resulting expression zero.
\end{proof}

\section{Localization and a global annular estimate}

\begin{lemma}[From bounded local inputs to endpoint operator norms]
\label{ann:localize}
\lean{Auto.ann_localize}
\leanok
Let $\mu_n=\sigma*\kappa_n$, where
$\sup_n\|\kappa_n\|_{\mathrm{TV}}<\infty$, and suppose that for
some fixed $U<\infty$,
\[
 |\mu_n|(\{u:|u|>U\})\longrightarrow0.
\]
Let $P_n$ be convolution by measures supported in the closed unit
ball, with uniformly bounded total variation. Suppose that for
every fixed triple $\chi_0,\chi_1,\chi_2\in C_c^\infty(\R^2)$,
\begin{equation}\label{ann:bounded-local-hyp}
 \sup_{\|f_j\|_\infty\le1}
 |\Lambda_{\mu_n}(\chi_0P_nf_0,\chi_1P_nf_1,\chi_2P_nf_2)|
 \longrightarrow0.
\end{equation}
Then the form
$\Lambda_{\mu_n}(P_nf_0,P_nf_1,P_nf_2)$ has operator norm
tending to zero on every permutation of
$L^2\times L^2\times L^\infty$.
\end{lemma}
\begin{proof}
Choose a nonnegative smooth partition of unity
$\sum_{q\in\mathbb Z^2}\chi_{0,q}(x)=1$, where the functions
are translates of one compactly supported function. For $j=1,2$
choose fixed compactly supported smooth cutoffs, between zero
and one, whose translates $\chi_{j,q}$ equal one at $x+L_ju$
whenever $x\in\operatorname{supp}\chi_{0,q}$ and $|u|\le U$.
Thus
\[
 0\le\sum_q\prod_{j=0}^2\chi_{j,q}(x+L_ju)\le1,
\]
and this sum is one when $|u|\le U$.
The total-variation endpoint estimate of
Lemma~\ref{lem:form-basic} bounds the error in inserting this
sum by
\[
 o_n(1)\|f_i\|_\infty\prod_{j\ne i}\|f_j\|_2
\]
for any fixed placement $i$ of the bounded input. The uniform
convolution bounds for $P_n$ have been absorbed into the error.

For each localized term, restrict the two $L^2$ inputs to boxes
$B_j+q$ containing the closed unit neighborhoods of the supports
of $\chi_{j,q}$. This restriction is exact because the convolution
measure defining $P_n$ is supported in the unit ball. Denote
their local $L^2$ norms by $r_{j,q}$. If either is zero, the term
is zero. Otherwise divide the two restricted inputs by these
norms, and divide the bounded input by its global supremum norm.
It is enough to bound the resulting localized expression when
the two local $L^2$ norms and the global $L^\infty$ norm are one.

Fix $T\ge1$. For either normalized $L^2$ input $g$, write
\[
 g=g\one_{\{|g|\le T\}}+g\one_{\{|g|>T\}}.
\]
The second term has $L^2$ norm at most one and is supported on
a set of area at most $T^{-2}$. In a form containing this term,
expand its convolution under $P_n$ as an integral of translates.
Each translated term, after multiplication by the cutoff, still
has support of area at most $T^{-2}$ and $L^2$ norm at most one.
Lemma~\ref{lem:small-support}, including its assertion for
$\sigma*\kappa_n$, gives an error
\[
 C\omega_M(T^{-2})
\]
uniformly in $n$ and $q$. The other projected $L^2$ input and the
projected bounded input have uniformly bounded norms. The same
argument removes the large part of the other $L^2$ input, with
the same type of bound.

The two remaining inputs have absolute value at most $T$.
Hypothesis~\eqref{ann:bounded-local-hyp} now bounds the normalized
localized form by
\[
 C\omega_M(T^{-2})+T^2\varepsilon_n,
 \qquad \varepsilon_n\longrightarrow0.
\]
The number $\varepsilon_n$ is independent of $q$: simultaneous
translation of the inputs and cutoffs reduces each $q$ to the
same fixed cutoff triple, since $P_n$ and $\Lambda_{\mu_n}$
commute with simultaneous translations.

Restore the local norms. If $j,k$ are the two $L^2$ indices,
Cauchy--Schwarz and bounded overlap of the boxes give
\[
 \sum_q r_{j,q}r_{k,q}
 \le\left(\sum_q r_{j,q}^2\right)^{1/2}
      \left(\sum_q r_{k,q}^2\right)^{1/2}
 \le C\|f_j\|_2\|f_k\|_2.
\]
This also justifies the sum of localized forms. First let
$n\to\infty$ with $T$ fixed, and then let $T\to\infty$.
The small-support modulus tends to zero, which proves the
claimed endpoint norm convergence for this placement. The
argument applies to each of the three choices of $i$.
\end{proof}

\begin{corollary}[Global cancellation on one annulus]\label{ann:global}
\lean{Auto.ann_global}
\leanok
Let $\psi\in C_c^\infty(\R^2)$ be supported in a compact
annulus disjoint from the origin. Put
\[
 K_R=\sigma*\bigl(R^2\check\psi(R\,\cdot)\bigr).
\]
For every fixed $0<C_0<\infty$ there is a function $a(R)\to0$,
with $0\le a(R)\le\|\check\psi\|_1$, such that whenever
$f_i\in L^\infty$ and $f_j\in L^2$ for $j\ne i$, and all three
inputs have Fourier supports in $\overline{B(0,C_0R)}$,
\[
 |\Lambda_{K_R}(f_0,f_1,f_2)|
 \le a(R)\|f_i\|_\infty\prod_{j\ne i}\|f_j\|_2,
 \qquad i=0,1,2.
\]
For the bounded input, Fourier support is understood through the
test-function condition in Definition~\ref{ext:bounded-support}.
\end{corollary}
\begin{proof}
Choose a Schwartz function $p$ whose Fourier transform is one
on a neighborhood of $\overline{B(0,C_0)}$. Then
$p_R=R^2p(R\,\cdot)$ reproduces each of the specified inputs
by Lemma~\ref{ext:reproduction}, applied at scale $R$.

Choose $\chi\in C_c^\infty(B(0,1))$ equal to one near zero and
put $q_R=\chi p_R$. Rapid decay gives
\[
 \|q_R-p_R\|_1=o(1),\qquad
 \sup_R\|q_R\|_1<\infty,\qquad
 \|\nabla q_R\|_1=O(R).
\]
Apply Lemma~\ref{ann:localize} along an arbitrary sequence
$R_n\to\infty$ with $P_nf=q_{R_n}*f$ and
$\mu_n=K_{R_n}$. Its convolution factor
$\kappa_n=R_n^2\check\psi(R_n\,\cdot)$ has bounded $L^1$ norm.
Moreover,
\[
 |K_{R_n}|(\{|u|>2\})
 \le\int_{|y|>1}R_n^2|\check\psi(R_ny)|\,\dd y
 \longrightarrow0.
\]
For fixed smooth compactly supported cutoffs and bounded inputs,
the functions $\chi_jP_nf_j$ have common bounded support,
uniformly bounded magnitude, and Lipschitz constants $O(R_n)$.
The integral of $\check\psi$ is $\psi(0)=0$, so
Proposition~\ref{ann:local} proves the hypothesis
\eqref{ann:bounded-local-hyp}. The localized lemma consequently
gives endpoint norm convergence with $q_R$ in all input slots.

Finally replace $q_R$ by $p_R$ in the three slots. Young's
inequality gives an operator difference at most
$\|q_R-p_R\|_1=o(1)$ on both $L^2$ and $L^\infty$.
The uniform total-variation bound on $K_R$ and
Lemma~\ref{lem:form-basic} show that the replacement costs
$o(1)$ in each endpoint form norm. Since $p_R$ reproduces the
inputs, this proves the stated estimate. Take $a(R)$ to be the
maximum of the actual restricted endpoint norms over the three
placements. The total-variation bound gives
$a(R)\le\|\check\psi\|_1$, and the argument proves $a(R)\to0$.
\end{proof}


\section{Summation over frequency annuli}\label{tail:section}

Throughout this section, the annular estimate of Corollary~\ref{ann:global}
is used for each of the three choices of the bounded input. All constants below
may depend on $M$ and on the smooth functions fixed in this section.

\begin{definition}[Annular kernels]\label{tail:annular-definition}
\lean{Auto.tailPsi,Auto.tailQ,Auto.tailK,Auto.tailForm,Auto.aux_tail_annular_definition}
Fix a real radial function $\chi\in C_c^\infty(\R^2)$ which equals one
on $B(0,1)$ and vanishes outside $B(0,2)$. Define
\[
 \psi(b)=\chi(b/2)-\chi(b),\qquad
 q_R(u)=R^2\check\chi(Ru),\qquad
 k_r(u)=r^2\check\psi(ru).
\]
Thus $\psi$ is supported in $\{1\le |b|\le4\}$ and
$\int q_R=\chi(0)=1$. Write
\[
 \Lambda_r=\Lambda_{\sigma*k_r}.
\]
The identities $k_r=q_{2r}-q_r$ and
$\|k_r\|_1=\|\check\psi\|_1$ hold for every $r>0$.
\end{definition}

\begin{lemma}[Frequency representation with compact spectral supports]
\label{tail:fourier}
\lean{Auto.tail_fourier}
\leanok
Suppose that $f_0,f_1,f_2\in L^2(\R^2)$ have compact Fourier supports.
Then
\begin{align*}
 \Lambda_r(f_0,f_1,f_2)
  =\iint &\widehat f_0(-\xi_1-\xi_2)
          \widehat f_1(\xi_1)\widehat f_2(\xi_2)\\
     &\hspace{10mm}\cdot\widehat\sigma(\xi_1+M^T\xi_2)
          \psi\big((\xi_1+M^T\xi_2)/r\big)
                \dd\xi_1\dd\xi_2.
\end{align*}
The double integral is absolutely convergent.
\end{lemma}
\begin{proof}
For Schwartz functions, this is Lemma~\ref{lem:fourier-form}, applied
with the kernel $\sigma*k_r$ of Definition~\ref{tail:annular-definition}.

For fixed $\xi_2$, Cauchy--Schwarz gives
\[
 \int |\widehat f_0(-\xi_1-\xi_2)\widehat f_1(\xi_1)|\dd\xi_1
 \le \|f_0\|_2\|f_1\|_2.
\]
The compact support and the $L^2$ integrability of $\widehat f_2$ imply
$\widehat f_2\in L^1$. Consequently the absolute value of the double
integral is bounded by
\[
 \|\widehat\sigma\psi(\cdot/r)\|_\infty
       \|f_0\|_2\|f_1\|_2\|\widehat f_2\|_1.
\]
Approximate the three Fourier transforms by smooth compactly supported
functions, with all approximants supported in fixed compact enlargements
of their original supports. They converge in $L^2$ and, for the third
input, also in $L^1$. The displayed bound passes the Fourier integral to
the limit. Lemma~\ref{lem:triple-l2}, applied after convolution with
$k_r$, passes the physical-space form to the same limit.
\end{proof}

\begin{definition}[Common frequency coordinates]\label{tail:coordinates}
\lean{Auto.tailA,Auto.tailT}
Put $A=M^T$ and define invertible linear maps
\[
 T_0=(A-I)^{-1},\qquad T_1=-A^{-1},\qquad T_2=I.
\]
For a frequency triple with $\xi_0+\xi_1+\xi_2=0$, write
$z_j=T_j\xi_j$ and $b=\xi_1+A\xi_2$.
\end{definition}

\begin{lemma}[Frequency separation and modulation invariance]
\label{tail:separation}
\lean{Auto.tailV,Auto.tail_separation_identities,Auto.tail_separation_comparable,Auto.tail_separation_modulation,Auto.tail_separation_translate}
\leanok
For the coordinates of Definition~\ref{tail:coordinates},
\[
 z_0-z_2=-(A-I)^{-1}b,\quad
 z_1-z_2=-A^{-1}b,\quad
 z_1-z_0=A^{-1}(A-I)^{-1}b.
\]
In particular, whenever $\psi(b/r)\ne0$, every pairwise distance
$|z_i-z_j|$ is bounded above and below by positive multiples of $r$.
For $w\in\R^2$ set
\[
 v_0=(A-I)w,\qquad v_1=-Aw,\qquad v_2=w.
\]
Replacing each $f_j(x)$ by $\e(-v_j\cdot x)f_j(x)$ leaves
$\Lambda_\nu(f_0,f_1,f_2)$ unchanged for every finite measure $\nu$.
This replacement translates all three $z$-frequency supports by $-w$.
\end{lemma}
\begin{proof}
Substitute $\xi_0=-\xi_1-\xi_2$ in the definition of $z_0$ and
subtract. The matrices $A$ and $A-I$, and hence their inverses,
commute. This gives the three displayed identities. Each matrix on the
right is invertible, so each norm is comparable to $|b|$. The support
condition in Definition~\ref{tail:annular-definition} gives
$r\le|b|\le4r$.

The modulation factors in the integrand of the form multiply to
\[
 \e\big(-x\cdot(v_0+v_1+v_2)-u\cdot(v_1+M^Tv_2)\big)=1.
\]
Indeed both vectors in parentheses vanish. Finally
$T_jv_j=w$ for every $j$, so the asserted translation of supports
follows from the Fourier transform of a modulation.
\end{proof}

\begin{definition}[Smooth frequency partitions and spectral mass]
\label{tail:partition-definition}
\lean{Auto.latticePt,Auto.tailP,Auto.tailSquare,Auto.tailMu,Auto.tailContributing}
Choose a nonnegative $\varphi\in C_c^\infty(\R^2)$ with
$\sum_{q\in\mathbb Z^2}\varphi(z-q)=1$. Fix a closed square
$Q_*$ containing its support. For a small fixed $h>0$, define
\[
 \widehat{P_{j,q,r}f}(\xi)
    =\varphi\big(T_j\xi/(hr)-q\big)\widehat f(\xi).
\]
Let $Q_{q,r}=hr(q+Q_*)$. For a fixed triple of $L^2$
inputs define the finite positive measure
\[
 \mu=\sum_{j=0}^2(T_j)_*
          \big(|\widehat f_j(\xi)|^2\dd\xi\big).
\]
A triple $(q_0,q_1,q_2)$ is called contributing at scale $r$ if there
are frequencies $\xi_j$ with $\sum_j\xi_j=0$, with
$T_j\xi_j\in Q_{q_j,r}$ for all $j$, and with
$\psi((\xi_1+A\xi_2)/r)\ne0$.
\end{definition}

\begin{lemma}[Properties of the smooth partition]\label{tail:partition}
\lean{Auto.tail_partition}
\leanok
The parameter $h$ can be chosen depending only on $M$ and $\varphi$ so
that the following statements hold uniformly in $r>0$.
\begin{enumerate}
\item The square family $\{Q_{q,r}\}$ has bounded overlap. The three
squares in every contributing triple have mutual distance at least
$d r$ and their union has diameter at most $D r$, for fixed $d,D>0$.
Every square in a specified slot occurs in at most $K$ contributing
triples.
\item There is a constant $C$ such that
\[
 \|P_{j,q,r}f_j\|_2^2\le C\mu(Q_{q,r}),\qquad
 \|P_{j,q,r}f_j\|_\infty\le C\|f_j\|_\infty,
 \qquad
 \sum_q\|P_{j,q,r}f_j\|_2^2\le C\|f_j\|_2^2.
\]
\item For $L^2$ inputs the following expansion is absolutely convergent:
\[
 \Lambda_r(f_0,f_1,f_2)
  =\sum_{(q_0,q_1,q_2)\,\mathrm{contributing}}
     \Lambda_r(P_{0,q_0,r}f_0,P_{1,q_1,r}f_1,P_{2,q_2,r}f_2).
\]
\end{enumerate}
\end{lemma}
\begin{proof}
The square side length is a fixed multiple of $hr$, and its centers
form a lattice of mesh $hr$. This proves bounded overlap. By
Lemma~\ref{tail:separation}, the witnessing points of a contributing
triple have pairwise distances between $c_Mr$ and $C_Mr$. Choose $h$
small enough that the diameters of the support squares are less than
$c_Mr/4$. Every pair of entire squares is then separated by a positive
multiple of $r$, and the union lies in a ball of radius a fixed multiple
of $r$. Once one square is fixed, the centers of either of its partners
lie within $O_M(r)$. There are only $O_{M,h}(1)$ lattice centers in this
region, which bounds the number of partner pairs.

Plancherel gives the first $L^2$ bound, since the multiplier vanishes
outside $T_j^{-1}Q_{q,r}$ and has uniformly bounded magnitude. The same
argument followed by the bounded-overlap estimate proves the sum of
squares bound. A translation, dilation, and fixed invertible linear
change of variables in frequency show that the convolution kernel of
$P_{j,q,r}$ has $L^1$ norm independent of $q$ and $r$. Young's inequality
gives the $L^\infty$ estimate.

Put $u_{j,q}=\|P_{j,q,r}f_j\|_2$. By the convolved version of
Lemma~\ref{lem:triple-l2}, each summand has absolute value at most
$C u_{0,q_0}u_{1,q_1}u_{2,q_2}$. The constant is independent of $r$
because $\|k_r\|_1$ is fixed. The bounded partner count and
Cauchy--Schwarz yield
\begin{align*}
 \sum_{\mathrm{contributing}}u_{0,q_0}u_{1,q_1}u_{2,q_2}
 &\le \Big(\sup_q u_{2,q}\Big)
      \Big(\sum_{\mathrm{contributing}}u_{0,q_0}^2\Big)^{1/2}
      \Big(\sum_{\mathrm{contributing}}u_{1,q_1}^2\Big)^{1/2}\\
 &\le C\prod_{j=0}^2\|f_j\|_2.
\end{align*}
This proves absolute convergence.

Finite partial partition sums converge to each input in $L^2$, by
Plancherel and dominated convergence. For finite sums, expand the form
by trilinearity. Lemma~\ref{tail:fourier} applies to every summand,
whose Fourier support is compact, and shows that every noncontributing
term vanishes. Pass to the limit using Lemma~\ref{lem:triple-l2} on
the left and the just-proved absolute convergence on the right.
\end{proof}

\begin{lemma}[One frequency triple]\label{tail:one-triple}
\lean{Auto.tail_one_triple}
\leanok
Suppose that $\|f_j\|_2,\|f_j\|_\infty\le B$. For a contributing
triple let $m\le n\le p$ be the three numbers
$\mu(Q_{q_j,r})$, in increasing order. For a sufficiently large fixed
frequency-ball constant in Corollary~\ref{ann:global}, its modulus
$a(r)$ satisfies
\[
 |\Lambda_r(P_{0,q_0,r}f_0,P_{1,q_1,r}f_1,P_{2,q_2,r}f_2)|
       \le C_B a(r)^{1/2}n p^{1/4}.
\]
\end{lemma}
\begin{proof}
Lemma~\ref{lem:triple-l2} and Lemma~\ref{tail:partition} first give
$C(mnp)^{1/2}$. Translate all three $z$-frequency squares by the negative
of the center of one square. Lemma~\ref{tail:separation} realizes this
as a form-preserving modulation. Lemma~\ref{tail:partition} shows that
the translated supports lie in a fixed ball of radius $C_Mr$ in each
original frequency plane. Apply Corollary~\ref{ann:global}, putting
the input corresponding to mass $p$ in the bounded slot. Its supremum
norm is at most $CB$, and the other two $L^2$ norms are at most
$C\sqrt m$ and $C\sqrt n$. This gives the second bound
$C_Ba(r)(mn)^{1/2}$.

The minimum of two nonnegative bounds is at most their geometric mean.
The resulting estimate is
$C_Ba(r)^{1/2}(mn)^{1/2}p^{1/4}$, which is at most the claimed bound
because $m\le n$. If a mass is zero, the relevant projected input is
zero in $L^2$ and the conclusion holds directly.
\end{proof}

\begin{definition}[An increasing mass functional]\label{tail:mass-definition}
\lean{Auto.tailH}
For a finite positive Borel measure $\mu$ on $\R^2$ and $t>0$, put
\[
 H(t)=\int_{\R^2}\mu(B(z,t))^{1/4}\dd\mu(z).
\]
This is nondecreasing and satisfies
$0\le H(t)\le\mu(\R^2)^{5/4}$.
\end{definition}

\begin{lemma}[Geometric mass increment]\label{tail:increment}
\lean{Auto.tail_increment}
\leanok
For the square families and contributing triples of
Lemma~\ref{tail:partition}, there are fixed constants $0<c<C$ such that
\[
 \sum_{\mathrm{contributing}}n p^{1/4}
       \le C_M\big(H(Cr)-H(cr)\big),
\]
where $m\le n\le p$ are the masses of the three squares in each triple.
\end{lemma}
\begin{proof}
Choose $c<d/4$, with $d$ from Lemma~\ref{tail:partition}. For a fixed
triple, the neighborhoods $N_j=Q_{q_j,r}+B(0,cr)$ are pairwise disjoint.
Choose a ball $B_*$ containing their union. Since the diameter of that
union is $O_M(r)$, the constant $C$ can be fixed so that
$B_*\subset B(z,Cr)$ for every $z$ in any of the three squares.
At least two of the disjoint neighborhoods have mass at most
$\mu(B_*)/2$. Select their indices $i,j$.

For $z\in Q_{q_i,r}\cup Q_{q_j,r}$, the smaller ball $B(z,cr)$ is
contained in the corresponding selected neighborhood. Therefore
\begin{align*}
 \mu(B(z,Cr))^{1/4}-\mu(B(z,cr))^{1/4}
 &\ge (1-2^{-1/4})\mu(B_*)^{1/4}\\
 &\ge (1-2^{-1/4})p^{1/4}.
\end{align*}
The selected squares are disjoint, and their combined mass is at least
$n$, irrespective of which two indices were selected. Integrating the
last inequality over their union bounds $n p^{1/4}$ by a fixed multiple
of that integral.

Sum over contributing triples. Each point belongs to only boundedly
many support squares, and each square has only boundedly many partner
pairs. Hence the multiplicity of the selected regions is uniformly
bounded. The integrand is nonnegative, so the sum is at most a constant
times its integral over the whole plane, which is $H(Cr)-H(cr)$.
\end{proof}

\begin{lemma}[Dyadic telescoping of the mass increments]
\label{tail:telescoping}
\lean{Auto.tail_telescoping}
\leanok
For $0<c<C$, every nondecreasing bounded function $H$ satisfies
\[
 \sum_{k=0}^{\infty}
       \big(H(C2^kR)-H(c2^kR)\big)
       \le D\sup_{t>0}H(t),
\]
provided $H\ge0$ and $D$ is any positive integer with $2^Dc\ge C$.
\end{lemma}
\begin{proof}
By monotonicity, the summand is at most
$H(c2^{k+D}R)-H(c2^kR)$. In a finite sum over $0\le k\le N$,
all but the first and last $D$ terms cancel. The terms with a minus
sign are nonpositive and may be discarded, leaving at most
$D\sup H$. The original summands are nonnegative, so monotone
convergence gives the infinite-sum bound.
\end{proof}

\begin{proposition}[Uniform annular tails]\label{tail:bound}
\lean{Auto.tail_bound,Auto.tail_bound_tendsto}
\leanok
Let $R\ge1$ be dyadic. If $\|f_j\|_2,\|f_j\|_\infty\le B$, then
\[
 \sum_{r=2^k\ge R}|\Lambda_r(f_0,f_1,f_2)|
       \le C_{M,B}\left(\sup_{r\ge R}a(r)\right)^{1/2}.
\]
In particular this tail tends to zero as $R\to\infty$, uniformly over
these inputs.
\end{proposition}
\begin{proof}
Combine the absolutely convergent expansion in
Lemma~\ref{tail:partition}, the estimate in
Lemma~\ref{tail:one-triple}, and Lemma~\ref{tail:increment}. At each
scale the result is
\[
 |\Lambda_r(f_0,f_1,f_2)|
   \le C_{M,B}a(r)^{1/2}\big(H(Cr)-H(cr)\big).
\]
The same spectral measure $\mu$ is used at every scale, and Plancherel
shows that $\mu(\R^2)=\sum_j\|f_j\|_2^2\le3B^2$. Pull out the
supremum of $a(r)^{1/2}$ and apply Lemma~\ref{tail:telescoping} and
Definition~\ref{tail:mass-definition}. Finally,
$a(r)\to0$ implies $\sup_{r\ge R}a(r)\to0$.
\end{proof}

\section{Approximation of the counting form}\label{tail:approximation-section}

\begin{lemma}[Exhaustion by smooth low frequencies]\label{tail:exhaustion}
\lean{Auto.tail_exhaustion}
\leanok
For every finite measure $\kappa$ and every triple of $L^2$ inputs,
\[
 \Lambda_{\sigma*\kappa*q_S}(f_0,f_1,f_2)
       \longrightarrow\Lambda_{\sigma*\kappa}(f_0,f_1,f_2)
       \qquad(S\to\infty).
\]
\end{lemma}
\begin{proof}
Write $q_S(u)=S^2\check\chi(Su)$ and use
Lemma~\ref{lem:convolution-translate}. Convolution by $q_S$ averages
simultaneous translations of the inputs at vertices $1$ and $2$ by
$y/S$ and $My/S$, with weight $\check\chi(y)\dd y$. For each fixed
$y$, these inputs converge to the originals in $L^2$. The convolved
version of Lemma~\ref{lem:triple-l2} therefore gives convergence of
the corresponding form. The absolute value is bounded uniformly in
$y$ and $S$ by
\[
 C_M\|\kappa\|_{\mathrm{TV}}\prod_j\|f_j\|_2.
\]
Dominated convergence against $|\check\chi(y)|\dd y$ applies, and
$\int\check\chi=1$ gives the stated limit. No positivity of
$\check\chi$ is required.
\end{proof}

\begin{lemma}[Tails after convolution]\label{tail:convolved-tail}
\lean{Auto.tail_convolved_tail,Auto.tail_convolved_tail_sum}
\leanok
For every finite measure $\kappa$, dyadic $R\ge1$, and inputs with
$\|f_j\|_2,\|f_j\|_\infty\le B$,
\[
 \sum_{r=2^k\ge R}|\Lambda_{\sigma*\kappa*k_r}(f_0,f_1,f_2)|
  \le C_{M,B}\|\kappa\|_{\mathrm{TV}}
       \left(\sup_{r\ge R}a(r)\right)^{1/2}.
\]
Moreover the sum without absolute values equals
$\Lambda_{\sigma*\kappa}(f_0,f_1,f_2)
-\Lambda_{\sigma*\kappa*q_R}(f_0,f_1,f_2)$.
\end{lemma}
\begin{proof}
Lemma~\ref{lem:convolution-translate} expresses each summand as an
integral against $\kappa$ of $\Lambda_r$ applied to $f_0$ and
translates of $f_1,f_2$. These translations preserve the two norms.
Apply Proposition~\ref{tail:bound}, first to finite sums, integrate
against $|\kappa|$, and then use monotone convergence.

For a finite consecutive dyadic sum the kernel identity in
Definition~\ref{tail:annular-definition} gives
\[
 \sum_{k=0}^{N}\sigma*\kappa*k_{2^kR}
   =\sigma*\kappa*(q_{2^{N+1}R}-q_R).
\]
Apply Lemma~\ref{tail:exhaustion} and the absolute convergence just
proved to pass to the limit.
\end{proof}

\begin{proposition}[Smoothing for inputs bounded in both norms]
\label{tail:bounded-smoothing}
\lean{Auto.tail_bounded_smoothing}
\leanok
Let $\rho\in L^1(\R^2)$ be nonnegative with $\int\rho=1$, and set
$\rho_\varepsilon(u)=\varepsilon^{-2}\rho(u/\varepsilon)$. For each
fixed $B<\infty$,
\[
 \sup_{\|f_j\|_2,\|f_j\|_\infty\le B}
 |\Lambda_{\sigma*\rho_\varepsilon-\sigma}(f_0,f_1,f_2)|
      \longrightarrow0\qquad(\varepsilon\downarrow0).
\]
\end{proposition}
\begin{proof}
Fix a dyadic $R\ge1$. Apply Lemma~\ref{tail:convolved-tail} to
$\kappa=\rho_\varepsilon(u)\dd u$ and to $\kappa=\delta_0$.
Each has total variation one. Thus the two high-frequency tails
contribute together at most
\[
 2C_{M,B}\left(\sup_{r\ge R}a(r)\right)^{1/2}.
\]
Their remaining difference has kernel
$\sigma*(\rho_\varepsilon*q_R-q_R)$. For fixed $R$, translation
continuity in $L^1$ and dominated convergence give
\begin{align*}
 \|\rho_\varepsilon*q_R-q_R\|_1
 &\le \int\rho(y)\|q_R(\,\cdot-\varepsilon y)-q_R\|_1\dd y
       \longrightarrow0.
\end{align*}
The integrand is bounded by $2\|q_R\|_1\rho(y)$, so this argument
requires only the stated $L^1$ assumptions on $\rho$. The total
variation estimate in Lemma~\ref{lem:form-basic} bounds the low-frequency
form by $B^3\|\rho_\varepsilon*q_R-q_R\|_1$.

Given a desired error, first choose $R$ sufficiently large to make the
tail bound small, and then choose $\varepsilon$ sufficiently small to
make this last term small. The estimates are uniform in the inputs.
\end{proof}

\begin{theorem}[Uniform smoothing in the endpoint form norm]
\label{tail:smoothing}
\lean{Auto.tail_smoothing}
\leanok
For every probability density $\rho\in L^1(\R^2)$,
\[
 \lim_{\varepsilon\downarrow0}\sup_{s>0}
      \|\Lambda_{(\sigma*\rho_\varepsilon-\sigma)_s}\|=0.
\]
The same conclusion holds for every permutation of the
$L^2\times L^2\times L^\infty$ norm.
\end{theorem}
\begin{proof}
Take an arbitrary sequence $\varepsilon_n\downarrow0$ and put
\[
 \kappa_n=\rho_{\varepsilon_n}(u)\dd u-\delta_0,
 \qquad \mu_n=\sigma*\kappa_n.
\]
Then $\|\kappa_n\|_{\mathrm{TV}}\le2$, and
\[
 |\mu_n|(\{|u|>2\})
   \le\int_{|w|>1}\rho_{\varepsilon_n}(w)\dd w
   =\int_{|y|>1/\varepsilon_n}\rho(y)\dd y\longrightarrow0.
\]
Use Lemma~\ref{ann:localize} with $P_n$ equal to the identity.
For every fixed triple of smooth compactly supported cutoffs $\chi_j$
and every $\|f_j\|_\infty\le1$, the functions $\chi_jf_j$ have
$L^2$ norms at most $\|\chi_j\|_2$ and supremum norms at most
$\|\chi_j\|_\infty$. Proposition~\ref{tail:bounded-smoothing}
therefore verifies its localized hypothesis, with a bound $B$ depending
only on these fixed cutoffs. The lemma gives convergence of the
endpoint form norm to zero, for each of the three bounded slots.
Since the sequence $\varepsilon_n$ was arbitrary, convergence holds
as $\varepsilon\downarrow0$.

Finally, dilation invariance from Lemma~\ref{lem:form-basic} gives
\[
 \|\Lambda_{(\sigma*\rho_\varepsilon-\sigma)_s}\|
       =\|\Lambda_{\sigma*\rho_\varepsilon-\sigma}\|
       \qquad(s>0).
\]
This proves the asserted uniformity in $s$.
\end{proof}



\section{From smoothing to triangles}\label{sec:density}

\begin{lemma}[Continuity under small changes of scale]\label{den:continuity}
\lean{Auto.den_continuity}
\leanok
For every endpoint placement $i$,
\[
 \lim_{r\to1}\|\Lambda_{\sigma_r-\sigma}\|_i=0.
\]
\end{lemma}
\begin{proof}
Fix a nonnegative smooth compactly supported probability density
$\rho$.  For $t>0$ let $k_t=\sigma*\rho_t$, where
$\rho_t(u)=t^{-2}\rho(u/t)$.  This is an $L^1$ density.
The identity
\[
 \sigma_r-\sigma=(\sigma-k_t)_r+((k_t)_r-k_t)+(k_t-\sigma)
\]
and Lemma~\ref{lem:form-basic} give
\[
 \|\Lambda_{\sigma_r-\sigma}\|_i
 \le2\|\Lambda_{\sigma-k_t}\|_i+\|(k_t)_r-k_t\|_1.
\]
Here $(k_t)_r(u)=r^{-2}k_t(u/r)$ as a density.  With $t$ fixed,
the last norm tends to zero by continuity of dilations in $L^1$.
Theorem~\ref{tail:smoothing} makes the first term arbitrarily
small as $t\downarrow0$.  Taking these limits in that order
proves the assertion.
\end{proof}

\begin{definition}[Density and triangle counts]\label{den:definitions}
\lean{Auto.squareAt,Auto.upperBanachDensity,Auto.triangleCount,Auto.IsMissed}
For a Lebesgue measurable set $E\subset\R^2$, its upper Banach
density is
\[
 d^*(E)=\limsup_{L\to\infty}\sup_{z\in\R^2}
           \frac{|E\cap(z+[0,L]^2)|}{L^2}.
\]
For a bounded, compactly supported, nonnegative function $f$ and
$s>0$, write
\[
 D_s(f)=\Lambda_{\sigma_s}(f,f,f)
       =\int_{\mathbb T}\int_{\R^2}
          f(x)f(x+su_\theta)f(x+sMu_\theta)
                                \dd x\dd\sigma(\theta).
\]
A scale $s$ is \emph{missed by $E$} if there are no
$x\in\R^2$, $\theta\in\mathbb T$ for which all three displayed
vertices belong to $E$.
\end{definition}

\begin{lemma}[Large dense squares and exact copies]\label{den:dense-squares}
\lean{Auto.den_dense_squares_exists,Auto.den_dense_squares_borel,Auto.den_dense_squares_missed,Auto.den_dense_squares_copy}
\leanok
If $0<\delta<d^*(E)$, there are squares $Q$ of arbitrarily
large side length for which $|E\cap Q|\ge\delta|Q|$.
For each such square there is a Borel set $A\subset E\cap Q$
of the same measure.  If $s$ is missed by $E$, then
$D_s(\one_A)=0$.  If $D_s(\one_A)>0$, then $A$, and hence
$E$, contains a translate of a rotation of
$s\{0,e_1,Me_1\}$.
\end{lemma}
\begin{proof}
Choose $\delta'$ strictly between $\delta$ and $d^*(E)$.
The definition of limsup supplies arbitrarily large $L$ for
which the supremum over $z$ of the density in $z+[0,L]^2$
exceeds $\delta'$.  Approximate this supremum closely enough
to obtain density at least $\delta$.  Existence of the Borel
subset follows from completeness and regularity of Lebesgue
measure.

The integrand defining $D_s(\one_A)$ is a nonnegative Borel
function.  It is identically zero when the scale is missed.
If its integral is positive, it equals one at some
configuration.  The vertices there have the form
$x+s\{0,u_\theta,Mu_\theta\}$, which is the required rotated
triangle because $M$ commutes with rotations.
\end{proof}

\begin{definition}[Word notation for the geometric application]\label{den:lines}
\lean{Auto.lineVarSet,Auto.subsum}
For $n\ge1$, use $\mathcal W_n=\{0,1,2\}^n$, obtained from
$[3]^n$ by the relabeling $j\mapsto j-1$.  Words, combinatorial
lines, their variable sets, and their fixed letters are as in
Definition~\ref{cj:def-words}, with this relabeling.
For positive numbers $\lambda_1,\ldots,\lambda_n$, set
$s_J=\sum_{j\in J}\lambda_j$.
\end{definition}

\begin{lemma}[Number of lines]\label{den:line-number}
\lean{Auto.den_line_number_fixed,Auto.den_line_number}
\leanok
For a fixed nonempty $J\subset[n]$ there are $3^{n-|J|}$
combinatorial lines with variable set $J$.  The total number
of lines in $\mathcal W_n$ is $4^n-3^n$.
\end{lemma}
\begin{proof}
There are three independent choices for each fixed coordinate.
Summing over the nonempty variable sets gives
\[
 \sum_{J\ne\varnothing}3^{n-|J|}
   =\sum_{j=1}^n\binom nj3^{n-j}=(3+1)^n-3^n.
\]
The variable set and fixed letters are uniquely determined by
the line, so this sum counts each line once.
\end{proof}

\begin{lemma}[A weighted lower bound from density Hales--Jewett]
\label{den:weighted-lower}
\lean{Auto.den_weighted_lower}
\leanok
Let $0<\delta\le1$ and choose
$n\ge H_3(\delta/4)$ from Theorem~\ref{cj:dhj3}.
Let $\lambda_1,\ldots,\lambda_n>0$ and put
\[
 S=\max(1,\|M\|)\sum_{j=1}^n\lambda_j.
\]
Suppose $A$ is a Borel subset of a square $Q$ with
$|A|\ge\delta|Q|$, and that
$Q'=Q+[-S,S]^2$ satisfies $|Q'|\le2|Q|$.
Then, for $f=\one_A$,
\begin{equation}\label{den:lower-bound}
 \sum_{\varnothing\ne J\subset[n]}3^{n-|J|}D_{s_J}(f)
       \ge\frac{\delta}{2}|Q|.
\end{equation}
\end{lemma}
\begin{proof}
Fix $\theta$.  Put $d_0=0$, $d_1=u_\theta$, and
$d_2=Mu_\theta$.  For a word $w\in\mathcal W_n$ let
$a_w(\theta)=\sum_j\lambda_jd_{w_j}$, and define the marked
fraction
\[
 m_\theta(x)=3^{-n}\sum_{w\in\mathcal W_n}
                     f(x+a_w(\theta)).
\]
Every coordinate of $a_w(\theta)$ has absolute value at most
$S$.  Hence $f(x+a_w(\theta))$ is supported in $Q'$ as a
function of $x$.  Translation invariance gives
\[
 0\le m_\theta\le1,
 \qquad \int_{Q'}m_\theta(x)\dd x=|A|\ge\delta|Q|.
\]
Let $G_\theta=\{x\in Q':m_\theta(x)\ge\delta/4\}$.
Splitting the last integral over $G_\theta$ and its complement
yields
\[
 \delta|Q|
 \le |G_\theta|+\frac{\delta}{4}|Q'|
 \le |G_\theta|+\frac{\delta}{2}|Q|,
\]
so $|G_\theta|\ge\delta|Q|/2$.

At each $x\in G_\theta$, at least $(\delta/4)3^n$ words are
marked by the condition $f(x+a_w(\theta))=1$.
Relabeling $\{1,2,3\}$ as $\{0,1,2\}$, Theorem~\ref{cj:dhj3}
therefore supplies a whole marked line.
If $\mathcal L_n$ is the finite family of all
lines, this means
\[
 \sum_{\ell\in\mathcal L_n}\prod_{w\in\ell}
               f(x+a_w(\theta))\ge\one_{G_\theta}(x).
\]
Integrate over all $x$; all summands are nonnegative.
For a line with variable set $J$, let
$b=\sum_{j\notin J}\lambda_jd_{a_j}$.  Its three offsets
are $b$, $b+s_Ju_\theta$, and $b+s_JMu_\theta$.
The translation $y=x+b$ turns its integral into
\[
 \int f(y)f(y+s_Ju_\theta)f(y+s_JMu_\theta)\dd y.
\]
This expression is independent of the fixed letters.  By
Lemma~\ref{den:line-number} it occurs $3^{n-|J|}$ times.
Integrating next in $\theta$ proves
\eqref{den:lower-bound}.  All functions in this argument are
Borel, and all exchanges of finite sums and integrals are
justified by nonnegativity or bounded compact support.
\end{proof}

\begin{lemma}[Selecting scales with close subsums]\label{den:select-scales}
\lean{Auto.den_select_scales}
\leanok
Let $\mathcal S\subset(0,\infty)$ be unbounded.  Given
$n\ge1$ and $\tau>0$, one can choose
$\lambda_1<\cdots<\lambda_n$ in $\mathcal S$ such that,
for every nonempty $J\subset[n]$ and $k=\max J$,
\[
 \|\Lambda_{\sigma_{s_J}-\sigma_{\lambda_k}}\|<\tau.
\]
\end{lemma}
\begin{proof}
By Lemma~\ref{den:continuity}, choose $\eta>0$ so that
$1\le r<1+\eta$ implies
$\|\Lambda_{\sigma_r-\sigma}\|<\tau$.
Choose any $\lambda_1\in\mathcal S$.  Having chosen the
previous scales, use unboundedness to choose
\[
 \lambda_k>\max\left(\lambda_{k-1},
                  \eta^{-1}\sum_{j<k}\lambda_j\right)
      \quad\text{in }\mathcal S.
\]
For $k=\max J$ this implies
\[
 1\le\frac{s_J}{\lambda_k}
    \le1+\frac{\sum_{j<k}\lambda_j}{\lambda_k}<1+\eta.
\]
Finally
$\sigma_{s_J}-\sigma_{\lambda_k}
  =(\sigma_{s_J/\lambda_k}-\sigma)_{\lambda_k}$,
so the norm equality in Lemma~\ref{lem:form-basic} gives
the desired inequality.  When $J=\{1\}$ it is immediate.
\end{proof}

\begin{theorem}[All sufficiently large planar triangles]\label{den:main}
\lean{Auto.den_main,Auto.den_main_general}
\leanok
If $E\subset\R^2$ is Lebesgue measurable with $d^*(E)>0$,
then every sufficiently large scale of
$\{0,e_1,Me_1\}$ has a translate of a rotation contained
in $E$.  Consequently the same conclusion holds for any
prescribed noncollinear planar triangle.
\end{theorem}
\begin{proof}
Fix $0<\delta<d^*(E)$.  Let $n$ be the integer in
Lemma~\ref{den:weighted-lower}.  Suppose for contradiction
that the set of missed scales is unbounded.  Apply
Lemma~\ref{den:select-scales} to that set with
$\tau=\delta/(4\cdot4^n)$ to obtain missed scales
$\lambda_1<\cdots<\lambda_n$.

These scales determine the finite number
$S=\max(1,\|M\|)\sum_j\lambda_j$.  By
Lemma~\ref{den:dense-squares}, choose a square $Q$ of
sufficiently large side length that
$|E\cap Q|\ge\delta|Q|$ and $|Q+[-S,S]^2|\le2|Q|$.
For example, the latter condition holds when its side
length is at least $2S/(\sqrt2-1)$.  Take a Borel set
$A\subset E\cap Q$ with $|A|=|E\cap Q|$ and set
$f=\one_A$.

For every chosen scale, $D_{\lambda_k}(f)=0$, by
Lemma~\ref{den:dense-squares}.  If $J\ne\varnothing$ and
$k=\max J$, the endpoint norm bound and the scale selection
give
\[
 0\le D_{s_J}(f)
   =|D_{s_J}(f)-D_{\lambda_k}(f)|
   \le \tau\|f\|_2^2\|f\|_\infty
   \le\frac{\delta|Q|}{4\cdot4^n}.
\]
On summing and using Lemma~\ref{den:line-number},
\[
 \sum_{J\ne\varnothing}3^{n-|J|}D_{s_J}(f)
    \le\frac{(4^n-3^n)\delta|Q|}{4\cdot4^n}
    <\frac{\delta}{4}|Q|.
\]
This contradicts the lower bound $\delta|Q|/2$ in
Lemma~\ref{den:weighted-lower}.  The missed scales are
therefore bounded, proving the first assertion.

For a general triangle, use Lemma~\ref{lem:geometry} to write
$\Delta=q_0+\ell U\{0,e_1,Me_1\}$ with $\ell>0$ and
$U$ a rotation.  A copy at scale $\lambda\ell$ of the
normalized triangle is a translated rotated copy of
$\lambda\Delta$, since products and inverses of planar
rotations are rotations.  Divide the normalized threshold
by $\ell$.  This proves the second assertion and the
introductory geometric conclusion.
\end{proof}

\section{The Fourier inverse consequence}\label{sec:inverse}

\begin{proposition}[A smooth-kernel Fourier estimate]\label{inv:smooth}
Suppose $k\in C_c^\infty(\R^2)$ and
$g\in L^2(\R^2)\cap L^\infty(\R^2)$ with
$\widehat g\in L^\infty(\R^2)$.  Then
\[
 |\Lambda_k(f_0,f_1,g)|
 \le C_M\|\widehat k\|_1\|\widehat g\|_\infty
                    \|f_0\|_2\|f_1\|_2
 \qquad(f_0,f_1\in L^2).
\]
\end{proposition}
\begin{proof}
First take $f_0,f_1$ Schwartz and let $A=M^T$.
The function paired with $g$ in the original form is
\[
 h(z)=\int k(u)f_0(z-Mu)f_1(z+(I-M)u)\dd u.
\]
Because $u$ is restricted to a compact set and the two
inputs are Schwartz, $h$ is Schwartz.  Fourier pairing with
the $L^2$ function $g$ is therefore justified by Plancherel.
The calculation of Lemma~\ref{lem:fourier-form} gives the
same Fourier expression with $\widehat g(\xi_2)$; it is
absolutely integrable by the following bound.

Make the change of variables $b=\xi_1+A\xi_2$, retaining
$\xi_2$.  Its Jacobian is one.  At fixed $b$,
Cauchy--Schwarz gives
\begin{align*}
 &\int\big|\widehat f_0(-b+(A-I)\xi_2)
              \widehat f_1(b-A\xi_2)\big|\dd\xi_2\\
 &\qquad\le
   |\det(A-I)|^{-1/2}|\det A|^{-1/2}
       \|\widehat f_0\|_2\|\widehat f_1\|_2.
\end{align*}
Both matrices are invertible by Lemma~\ref{lem:geometry}.
Bound $\widehat g$ by its essential supremum and integrate
against $|\widehat k(-b)|$.  Plancherel proves the estimate
for Schwartz tests.  Approximate arbitrary $L^2$ tests in
$L^2$ by Schwartz functions.  The original form converges
by Lemma~\ref{lem:form-basic}, since $g\in L^\infty$ and
$k\in L^1$, so the same bound holds in the limit.
\end{proof}

\begin{corollary}[Fourier inverse estimate]\label{inv:consequence}
Let $\rho$ be a smooth compactly supported probability
density and put $k_t=\sigma*\rho_t$.  For
$g\in L^2\cap L^\infty$ with bounded Fourier transform,
\begin{align*}
 \sup_{\|f_0\|_2,\|f_1\|_2\le1}
       |\Lambda_\sigma(f_0,f_1,g)|
 &\le \|\Lambda_{\sigma-k_t}\|\,\|g\|_\infty\\
 &\quad+C_M\|\widehat{k_t}\|_1\|\widehat g\|_\infty.
\end{align*}
In particular, if $g_n\in L^2\cap L^\infty$,
$\sup_n\|g_n\|_\infty<\infty$, and
$\|\widehat g_n\|_\infty\to0$, then
\[
 \sup_{\|f_0\|_2,\|f_1\|_2\le1}
       |\Lambda_\sigma(f_0,f_1,g_n)|\longrightarrow0.
\]
No uniform $L^2$ bound on $g_n$ is needed for this conclusion.
\end{corollary}
\begin{proof}
The convolution $k_t$ is smooth and compactly supported,
so Proposition~\ref{inv:smooth} applies and
$\widehat{k_t}\in L^1$.  Split
$\Lambda_\sigma=\Lambda_{\sigma-k_t}+\Lambda_{k_t}$.
The first term has its defining endpoint norm, and the
second has the bound in the proposition.  This proves the
first assertion.  For the sequence, fix $t$ and take the
upper limit in $n$.  The second term tends to zero, while
the first is bounded by
$\sup_n\|g_n\|_\infty\|\Lambda_{\sigma-k_t}\|$.
Let $t\downarrow0$ and use Theorem~\ref{tail:smoothing}.
\end{proof}

\begin{thebibliography}{9}
\bibitem{FK}
H. Furstenberg and Y. Katznelson,
\emph{A density version of the Hales--Jewett theorem},
J. Anal. Math. \textbf{57} (1991), 64--119.
\href{https://doi.org/10.1007/BF03041066}{doi:10.1007/BF03041066}.

\bibitem{Polymath}
D. H. J. Polymath,
\emph{A new proof of the density Hales--Jewett theorem},
Ann. of Math. (2) \textbf{175} (2012), 1283--1327.
\href{https://doi.org/10.4007/annals.2012.175.3.6}{doi:10.4007/annals.2012.175.3.6}.

\bibitem{DKT}
P. Dodos, V. Kanellopoulos, and K. Tyros,
\emph{A simple proof of the density Hales--Jewett theorem},
Int. Math. Res. Not. IMRN \textbf{2014} (2014), no. 12, 3340--3352.
\href{https://doi.org/10.1093/imrn/rnt041}{doi:10.1093/imrn/rnt041}.
\end{thebibliography}
\end{document}
